% =============================================================================
% RAPPORT V1 COMPLET AVEC LES REMARQUES DE VADIM (collé le 03/09/2026)
% -----------------------------------------------------------------------------
% Usage : Vadim colle ici l'intégralité du rapport LaTeX v1 tel qu'exporté
% d'Overleaf, remarques incluses (\textcolor{red}{...} et remarques en clair
% dans le texte). Claude s'y réfère pendant la passe v2, section par section,
% dans l'ordre : §4 -> §3 -> §2 -> intro/conclusion -> annexes.
%
% Ce fichier N'EST PAS compilé. Il sert de source de vérité du texte v1 et des
% remarques. Les réécritures v2 sont livrées dans les réponses (et, si utile,
% dans docs/rapport/v2/<section>.tex).
%
% Voir aussi : REMARQUES_V2_2026-09-03.md (remarques déjà extraites de la
% partie reçue le 03/09 : abstract, intro, §2.1-2.4 partiel) et
% SKILL_REDACTION.md (règles de style v2).
% =============================================================================

% >>> COLLER LE RAPPORT COMPLET CI-DESSOUS <<<

\documentclass[11pt,a4paper]{article}
\usepackage[utf8]{inputenc}
\usepackage[T1]{fontenc}
\usepackage[french]{babel}
\usepackage{amsmath,amssymb}
\usepackage{geometry}
\usepackage{hyperref}
\usepackage{xcolor}
\usepackage{booktabs} 
\usepackage[most]{tcolorbox}


\geometry{margin=2.5cm}

\title{\Large \textbf{Rapport d'Avancement de Stage de Recherche}\\ \huge \textbf{Auto-Amélioration des Agents LLM en Environnement Multi-tour}\\ \large Criteo AI Lab}
\author{\textbf{Vadim Lagresle} \\ Encadrants : Alberto Lumbreras, Patrick Gallinari, Alain Rakotomamonjy, Sylvain Lamprier}
\date{\today}

\begin{document}

\maketitle

\vspace{2cm}
\begin{abstract}
\noindent Ce rapport présente une étude expérimentale \textcolor{red}{/partiellemnt exploratoire} du post-entraînement d'agents LLM via apprentissage par renforcement en environnement interactif multi-tour. Le passage à l'échelle des environnements d'entraînement, par leur diversité, leur difficulté et l'utilisation de vérificateurs, est devenu le principal moteur des capacités agentiques, reformulant ainsi le goulot d'étranglement de la donnée annotée. Ne pouvant prétendre à cette échelle \textcolor{red}{Non pas besoin d'être négatif ici juste dire que ce qui nous intéresse c'est la stabilisation de l'apprentissage dans ces env, et l'optim de cet apprentissage}, nous étudions ce qui conditionne la stabilité et l'efficience de cet apprentissage dans un environnement contrôlé à récompense vérifiable (TextCraft, écosystème AgentGym-RL) \textcolor{red}{dire qqp que c'est à terme pour adapter ça à des env de shop}. Pour optimiser cet apprentissage sous contrainte matérielle, nous mobilisons l'apprentissage par curriculum, une méthode consistant à structurer et faire évoluer progressivement la difficulté des tâches auxquelles l'agent est confronté \textcolor{red}{et lora pour l'efficience, lora à placer en premier}. Nos contributions sont au nombre de trois. La première constitue un retour d'expérience \textcolor{red}{détaillé} technique détaillant les obstacles rencontrés et les solutions de stabilisation de l'algorithme GRPO en environnement textuel multi-tour \textcolor{red}{qui pourra nous l'espérons permettre à criteo de gagner du temps pour le développement de ses agents shoppers}. Ce bilan documente de manière honnête la gestion des observations, les risques de détournement induits par des récompenses mal définies (reward hacking), et le traitement des instabilités inhérentes à l'adaptation LoRA, offrant ainsi des clés utiles pour la reproductibilité de ce type d'entraînement \textcolor{red}{à mon avis cette phrase est trop longue et complexe pour un abstract}. La deuxième contribution compare différentes approches de curriculum, en opposant des stratégies conçues par l'humain à des régimes autotéliques (pilotés de manière adaptative). Nos résultats démontrent l'efficacité de ces méthodes : le meilleur curriculum atteint un taux de réussite de 82\,\%, dépassant les 75\,\% de l'état de l'art établis sur une infrastructure multi-GPU industrielle, tout en agissant comme un véritable régularisateur d'exploration vis-à-vis d'un contrôle non guidé \textcolor{red}{la phrase est looooooongue avec trop de virgules il faut vraiment pas laisser ce genre de structure dans le rapport, et encore moins dans l'abstract. Qui plus est du coup on ne va pas parler d'agents autoteliques ici, ce sera notre ouverture}. \textcolor{red}{j'aurai peut etre plutot mis en avant les deux principaux travaux qui sont stabilisation LoRA et Comparaison de curriculums dans la stabilisation, et noter par ailleurs qqp que Agentgym qui a précédé agentgym rl ne faisait pas de rl car instable à l'époque !}
Enfin, la troisième contribution \textcolor{red}{je dirai pas contribution c'est trop prétentieux. Nos deux principaux travaux se concentre sur l'optimisation via LoRA et la stabilisation via Curriculum Learning. Enfin, nous proposons une analyse de nos résultats en essayant de quantifier les caoacités de transfer learning et de l'utilité des curriculums au delà de la simple stabilsation mais vers des perspectives plus meta learning et continual learning} propose une analyse critique et une interprétation de nos résultats. Nous y discutons des dynamiques d'apprentissage, notamment le transfert asymétrique de compétences entre planification et contrôle interactif, le rôle de l'amorçage initial, et remettons en question l'hypothèse du plafond pass@$k$ en démontrant que les performances d'un modèle après apprentissage par renforcement (pass@1) peuvent surpasser les capacités exploratoires (pass@20) du modèle de base.
\end{abstract}


\textcolor{red}{ajouter un paragrpahe remerciements : patrick alberto alain pour la supervision depuis le début, sylvain pour les agents autoteliques, flavian et imad pour la découverte du deepshoper, Otmane pour les discussions autour des snis}.






\text




\vspace{0.5cm}

\newpage
\setcounter{tocdepth}{2} % Limite le sommaire aux sections et sous-sections (exclut les subsubsections)
\tableofcontents

\newpage
\section*{Notations et Abréviations}
\addcontentsline{toc}{section}{Notations et Abréviations}

\subsection*{Notations}
\begin{tabular}{ll}
\toprule
Symbole & Signification \\
\midrule
$\pi_\theta$ & politique entraînée, paramétrée par les poids $\theta$ \\
$\pi_{\text{ref}}$ & politique de référence de la pénalité KL (ancre) \\
$\tau$ & trajectoire (épisode complet) \\
$s_t,\ a_t,\ o_t$ & état, action et observation au pas $t$ \\
$T$ & horizon de l'épisode (nombre maximal de tours) \\
$R(\tau)$, $R_i$ & récompense terminale d'une trajectoire ; celle du rollout $i$ \\
$\hat{A}$, $\hat{A}_i$ & estimateur d'avantage ; avantage relatif du rollout $i$ dans son groupe \\
$V_\phi$ & fonction de valeur estimée par le réseau critique (PPO) \\
$r_t(\theta)$ & ratio d'importance $\pi_\theta(a_t \mid s_t)/\pi_{\theta_{\text{old}}}(a_t \mid s_t)$ \\
$G$ & taille de groupe GRPO (rollouts par tâche dans un même groupe) \\
$N$ & nombre de tirages indépendants (best-of-$N$, budget de pass@$k$) \\
$\beta$ & coefficient de la pénalité $D_{\text{KL}}(\pi_\theta \parallel \pi_{\text{ref}})$ \\
$\epsilon$ & rayon de \emph{clipping} de l'objectif PPO/GRPO \\
$r,\ \alpha$ & rang et facteur d'échelle d'un adaptateur LoRA \\
$d$ & profondeur (\emph{depth}) d'une tâche TextCraft, $d \in \{1,\dots,4\}$ \\
\bottomrule
\end{tabular}

\medskip
\noindent\textbf{pass@$k$.} Fraction des tâches de l'ensemble d'évaluation résolues par
au moins un tirage parmi $k$ tirages indépendants de la politique. Le pass@$1$
mesure la fiabilité d'un tirage unique ; le pass@$k$ pour $k$ grand approche le
support de la politique. Toute valeur de pass@$k$ est rapportée avec son budget
$k$ et sa température d'échantillonnage.

\subsection*{Abréviations}
\begin{tabular}{ll}
\toprule
Abréviation & Signification \\
\midrule
LLM & grand modèle de langage (\emph{Large Language Model}) \\
SFT & \emph{Supervised Fine-Tuning} \\
RL & apprentissage par renforcement \\
RLHF & RL à partir de préférences humaines \\
RLVR & RL à récompense vérifiable \\
NLL & log-vraisemblance négative \\
PPO & \emph{Proximal Policy Optimization} \\
GRPO & \emph{Group Relative Policy Optimization} \\
DPO & \emph{Direct Preference Optimization} \\
KL & divergence de Kullback-Leibler \\
(PO)MDP & processus de décision markovien (partiellement observable) \\
LoRA & \emph{Low-Rank Adaptation} \\
CoT & \emph{Chain-of-Thought} \\
PRM & \emph{Process Reward Model} \\
MCP & \emph{Model Context Protocol} \\
SNIS & échantillonnage préférentiel auto-normalisé \\
\bottomrule
\end{tabular}



\newpage
\section{Introduction}

\subsection{Évolution du Signal d'Apprentissage}

Les agents auto-améliorants (\emph{self-improving agents}) comptent aujourd'hui parmi les sujets les plus actifs du domaine de l'apprentissage machine. Cela est notamment visible à travers les cours universitaires dédiés comme le très récent CS329A dispensé à Stanford, le nombre de publications dans les grandes conférences d'apprentissage automatique, et les investissements massifs des constructeurs dans ce domaine. Pour comprendre pourquoi ce sujet attire autant, il peut être utile de retracer l'évolution du signal d'apprentissage des LLMs au cours de la dernière décennie. Au delà de l'opportunisme que présentent les applications agentiques, cette évolution suit les précédentes mutations du domaine, comme une conséquences du changement de la nature de la donnée et du signal qui en est extrait.

La rupture initiale introduite par l'architecture Transformer \cite{vaswani2017attention} a permis de passer à l'échelle sur la tâche d'autocomplétion autorégressive, en prédisant le token suivant sur d'immenses corpus textuels issus du web. Ce pré-entraînement massif confère au modèle une vaste représentation linguistique du monde. Toutefois, ce signal d'apprentissage reste implicite et passif : le modèle apprend à imiter la distribution du texte web sans garantie d'alignement avec les exigences de précision, de logique ou de sécurité attendues par un utilisateur humain.

Pour surmonter ces contraintes, la phase de post-training s'est imposée comme l'étape centrale d'alignement comportemental. Dans un premier temps, le Supervised Fine-Tuning a structuré la forme des réponses en instructions (il manque une référence ici). Par la suite, l'apprentissage par renforcement à partir de préférences humaines (RLHF) \cite{ouyang2022training} a introduit un signal d'évaluation global, permettant de pénaliser les réponses toxiques ou hors sujet. Plus récemment, les efforts se sont concentrés sur l'émergence de capacités de raisonnement complexe, où le modèle est encouragé à produire une trace de réflexion (\emph{Chain-of-Thought}) avant sa réponse \cite{wei2022cot}. Dans ce cadre, la nature du signal a de nouveau évolué vers des récompenses vérifiables de manière déterministe ($0$ ou $1$), notamment en mathématiques ou en programmation, qui permettent d'évaluer la justesse d'une séquence de pensée sans dépendre d'annotations humaines subjectives. Cet apprentissage par renforcement à récompense vérifiable (RLVR) a permis de faire passer à l'échelle les LLMs sur des tâches de raisonnement complexes, et il est au cœur des performances des modèles actuels \cite{openai_o1_2024, deepseekr1_2025}. Son intérêt pratique est double : il s'affranchit de l'entraînement et de l'inférence d'un modèle de récompense, et il renoue avec les questions classiques de l'apprentissage par renforcement, formulées en objectifs, en tâches et en environnements, que la vague des LLMs avait temporairement éclipsées. De nombreuses applications restent hors de portée du RLVR, comme l'évaluation de résumés ou la recommandation ; nous ne les traitons pas ici, et visons à maîtriser les techniques scalables avant d'aborder leurs applications spécifiques.

Dans le cadre du RLVR, les trajectoires d'entraînement se heurtent aujourd'hui à une limite matérielle et économique critique : la rareté et le coût élevé des données humaines de haute qualité dans les domaines exigeants. La donnée annotée constitue désormais le goulot d'étranglement majeur de la discipline, ce qui impose un changement de paradigme vers la génération de données synthétiques et l'apprentissage autonome. Des approches d'auto-amélioration telles que STaR \cite{zelikman2022star} ou ReST \cite{singh2024beyond} ont démontré qu'un modèle pouvait générer ses propres traces de raisonnement, filtrer les trajectoires menant à la bonne réponse, puis se réentraîner sur ses propres succès. Ce processus itératif accroît l'autonomie du modèle tout en réduisant considérablement la dépendance envers la supervision humaine directe.

Une thèse théorique influente interprète cette dynamique avec scepticisme : l'auto-amélioration sur données synthétiques ne créerait pas d'information fondamentalement nouvelle, et s'analyserait comme un mécanisme de \emph{sharpening} \cite{huang2024self}, où le réentraînement concentre la masse de probabilité de la politique sur les trajectoires à haute récompense, amortissant le coût du calcul à l'inférence sans étendre les capacités. Ce rapport traite cette thèse comme une hypothèse de travail à mettre à l'épreuve, non comme un acquis : une part importante de notre étude consiste précisément à mesurer si une structuration adéquate de l'apprentissage permet de dépasser les frontières que cette lecture prédit. L'émergence de l'IA agentique s'inscrit dans la continuité directe de cette évolution : il ne s'agit plus seulement de générer une trace de raisonnement en un seul jet, mais de déployer ces capacités au sein d'environnements numériques interactifs, où l'agent apprend de manière autonome à partir des retours de ses actions.

Le cadre agentique reformule enfin le goulot d'étranglement de la donnée sous une forme nouvelle : ce qui doit désormais passer à l'échelle n'est plus un corpus statique de textes ou d'annotations, mais l'environnement d'interaction lui-même, c'est-à-dire la définition conjointe d'un état initial, d'une surface d'outils, d'une tâche et d'un vérificateur. Les rapports techniques récents des constructeurs de modèles convergent sur ce point. L'équipe Qwen identifie explicitement le passage à l'échelle des environnements (\emph{environment scaling}) comme le verrou central du RL agentique \cite{qwen3coder2025}, et attribue l'essentiel des gains de ses dernières générations à l'extension massive de la difficulté et de la diversité de ses tâches d'entraînement \cite{qwen35_2026}. Kimi K2 fait reposer ses capacités agentiques sur un pipeline de synthèse à grande échelle d'outils, de tâches et de trajectoires vérifiées, couplé à un entraînement RL conjoint sur environnements réels et simulés \cite{kimik2_2025}. Symétriquement, les benchmarks de référence du domaine ($\tau^2$-bench \cite{tau2bench2025}, Toolathlon \cite{toolathlon2025}, MCP-Atlas \cite{mcpatlas2026}) sont moins des jeux de données que des travaux de définition du problème : y construire une tâche, c'est spécifier un monde, ses règles et son critère de succès. La définition de l'environnement devient ainsi un objet scientifique à part entière, au même titre que la curation des corpus l'a été pour le pré-entraînement. Ce constat structure l'ensemble de ce rapport.

\subsection{Rappel et similitudes des Algorithmes d'Entraînement}

Nous partons du principe que les techniques des transformers et du NLP sont connues du lecteur, et renvoyons aux cours de Stanford sur le \emph{language modeling} pour une présentation complète (\href{https://cs336.stanford.edu/}{cs336.stanford.edu}) et ajouter aussi un cours sur les llms de stanford. De même, cette sous-section ne présente que les éléments d'apprentissage par renforcement nécessaires à la lecture du rapport, en particulier GRPO, notre algorithme de travail. L'annexe A développe le formalisme complet, du cadre classique du RL (POMDP) aux méthodes acteur-critique ; pour un traitement approfondi du Deep RL, nous renvoyons aux cours de Stanford (\href{https://cs224r.stanford.edu/}{CS224R}) et de Berkeley (\href{https://rail.eecs.berkeley.edu/deeprlcourse/}{CS285}).

L'évolution de la nature des données et du signal d'apprentissage s'est accompagnée d'une diversité des algorithmes d'optimisation utilisés à chaque phase. Le pré-entraînement repose sur l'architecture Transformer \cite{vaswani2017attention}, qui a supplanté les architectures récurrentes et séquentielles. Ces dernières souffraient d'un goulot d'étranglement computationnel empêchant la parallélisation et d'un phénomène de disparition du gradient sur les longs contextes. Le mécanisme d'auto-attention permet de calculer en parallèle les dépendances entre tous les tokens d'une séquence, indépendamment de leur distance relative. Cette capacité de parallélisation massive sur GPU a rendu possible le traitement de corpus web bruts à très grande échelle.

Dans la même logique, la phase de Supervised Fine-Tuning conserve la même architecture autorégressive et le même principe de prédiction de token que le pré-entraînement, mais s'en distingue par trois ajustements techniques, au-delà du volume restreint et de la haute qualité des données. Tout d'abord, l'injection de tokens de contrôle structurants (\emph{chat templates}) transforme le texte continu en une machine à états alternant les rôles d'utilisateur et d'assistant, ce qui conditionne la génération vers un comportement d'agent conversationnel. Ensuite, le régime d'optimisation est adapté : un taux d'apprentissage nettement plus faible sur un nombre réduit d'époques concentre la masse de probabilité sur les structures de réponse attendues, sans détruire les représentations acquises. Enfin, alors que le pré-entraînement calcule le gradient sur l'intégralité du texte, le SFT applique un masquage de la perte (\emph{loss masking}) sur les tokens de l'instruction $x = (x_1, \dots, x_N)$, afin de ne rétropropager le signal que sur les tokens de la réponse cible $y = (y_1, \dots, y_M)$. Le modèle minimise la log-vraisemblance négative moyenne calculée exclusivement sur les $M$ tokens de la réponse :
\begin{equation}
\mathcal{L}_{\text{SFT}}(\theta) = -\frac{1}{M} \sum_{m=1}^{M} \log \pi_\theta(y_m \mid x, y_{<m})
\end{equation}

Contrairement au SFT, dont la mécanique de mise à jour des poids reste proche de celle du pré-entraînement, l'alignement par renforcement introduit un signal et un cadre d'optimisation fondamentalement différents. Il convient d'abord d'établir la différence fondamentale entre les deux. En fine-tuning, le modèle apprend à reproduire, token par token, une séquence cible : le signal d'apprentissage est le token suivant sachant la séquence qui précède. En RL, le signal est la récompense (\emph{reward}), un scalaire quantifiant la qualité globale d'une réponse produite par l'agent. L'objectif devient de favoriser les trajectoires menant aux récompenses les plus élevées : dans le cas d'une récompense \emph{sparse} binaire, typique des benchmarks mathématiques par exemple, on renforce les trajectoires ayant atteint la bonne réponse et l'on dévalorise les autres. On parle d'apprentissage \emph{par renforcement}, car ce processus renforce la politique du modèle autour d'une certaine manière d'enchaîner les tokens. A reformuler : L'algorithme renforce la distribution de sortie du modèle vers une distribution cible qu'il contient déjà grâce au pre training.  

Pour exploiter un signal de récompense scalaire $R(\tau)$ défini sur une trajectoire entière $\tau$, il faut exprimer un gradient, mais la récompense est rarement dérivable contrairement aux fonctions de pertes utilisées en fine tuning supervisé. Les méthodes de gradient de politique s'appuient sur l'astuce de REINFORCE (dérivée du logarithme), qui réécrit le gradient de l'espérance de la récompense sous une forme calculable par échantillonnage Monte-Carlo :
\begin{equation}
\nabla_\theta \mathbb{E}_{\tau \sim \pi_\theta}[R(\tau)] = \mathbb{E}_{\tau \sim \pi_\theta} \left[ \sum_{t=1}^T \nabla_\theta \log \pi_\theta(a_t \mid s_t) R(\tau) \right]
\end{equation}
Sans cette réécriture, le gradient ne pourrait pas être estimé par échantillonnage.

\textbf{Remarque.} Le SFT et le RL ne sont pas fondamentalement éloignés dans leur mécanique de mise à jour des poids, bien que les théories sous-jacentes diffèrent. Sous une récompense \emph{sparse} et une stratégie strictement \emph{on-policy}, c'est-à-dire lorsque l'entraînement s'effectue sur des générations produites par la politique courante (annexe A), l'algorithme de gradient de politique se réduit exactement à un SFT appliqué aux seules trajectoires ayant atteint la bonne réponse. Cette équivalence illustre à la fois la richesse et l'instabilité du RL : il identifie sans annotation humaine les trajectoires à privilégier, mais le seul signal disponible est la réussite finale, alors que des stratégies erronées peuvent occasionnellement mener à la bonne réponse. Le faible coût d'annotation se paie par une difficulté accrue de stabilisation, ces algorithmes pouvant s'effondrer (\emph{collapse}) en quelques itérations ; les sections 2.3 et 4.2 y reviennent en détail.

Utilisée seule, la récompense $R(\tau)$ constitue un signal brut : elle ne dit pas si une trajectoire est bonne en soi, seulement quel gain elle a rapporté, sans point de comparaison. Le cadre acteur-critique introduit un réseau critique $V_\phi(s)$, entraîné à estimer la fonction de valeur, c'est-à-dire l'espérance du retour obtenu depuis l'état $s_t$ en suivant la politique courante. Cette quantité ne dépend que de l'état, et non de l'action jouée : elle fournit une référence permettant de recentrer le signal sous la forme d'un avantage, dont l'estimateur le plus simple s'écrit $\hat{A}_t = R_t - V_\phi(s_t)$ (l'annexe A détaille les estimateurs employés en pratique). Substituer $\hat{A}_t$ à $R(\tau)$ revient à ne renforcer une action que dans la mesure où elle a fait mieux que la moyenne attendue, ce qui réduit considérablement la variance. À cette structure s'ajoute, pour les LLMs, une pénalité de divergence de Kullback-Leibler ancrant la politique entraînée $\pi_\theta$ à une politique de référence figée $\pi_{\text{ref}}$, généralement le modèle issu du SFT, afin d'éviter la dérive catastrophique et l'oubli des capacités de base :
\begin{equation}
\mathcal{L}_{\text{RL}}(\theta) = -\mathbb{E}_{(s_t, a_t)} \left[ \log \pi_\theta(a_t \mid s_t) \hat{A}_t \right] + \beta \, D_{\text{KL}}(\pi_\theta \parallel \pi_{\text{ref}})
\end{equation}

C'est dans ce cadre acteur-critique que s'inscrit \emph{Proximal Policy Optimization} (PPO) \cite{schulman2017proximal}, conçu pour le contrôle et la robotique bien avant les LLMs modernes, puis repris à partir de 2022 comme l'algorithme central du RLHF \cite{ouyang2022training}. Réutiliser un même lot de trajectoires pour plusieurs pas de gradient améliore considérablement l'efficacité en données, mais impose de corriger l'écart entre la politique courante et celle qui a généré les trajectoires. Un ratio d'importance $r_t(\theta) = \pi_\theta(a_t \mid s_t) / \pi_{\theta_{\text{old}}}(a_t \mid s_t)$ mesure cet écart, et un mécanisme de \emph{clipping} le borne, retirant à la politique tout intérêt à s'éloigner fortement de $\pi_{\theta_{\text{old}}}$ :
\begin{equation}
\mathcal{L}_{\text{PPO}}(\theta) = -\mathbb{E}_{(s_t, a_t)} \left[ \min\left( r_t(\theta)\, \hat{A}_t, \ \text{clip}\big(r_t(\theta),\, 1-\epsilon,\, 1+\epsilon\big)\, \hat{A}_t \right) \right]
\end{equation}

Notre algorithme de travail, \emph{Group Relative Policy Optimization} (GRPO) \cite{shao2024deepseekmath}, en est l'évolution la plus utilisée en RLVR. Il supprime le réseau critique, très coûteux en mémoire GPU, en échantillonnant un groupe de $G$ sorties $\{o_1, \dots, o_G\}$ pour une même requête, et en remplaçant l'estimateur d'avantage par un avantage relatif normalisé au sein du groupe, l'objectif clippé de PPO restant conservé :
\begin{equation}
\hat{A}_i = \frac{R_i - \text{mean}(\mathbf{R})}{\text{std}(\mathbf{R}) + \delta}
\end{equation}
où $R_i$ est la récompense de la sortie $i$ et $\mathbf{R}$ le vecteur des récompenses du groupe. Cette construction a une conséquence directe qui traverse tout le rapport : si toutes les sorties d'un groupe échouent, ou si toutes réussissent, l'écart-type est nul, l'avantage aussi, et la requête ne produit aucun gradient. La quantité de signal d'apprentissage dépend donc de l'adéquation entre la difficulté des tâches et le niveau courant de la politique.

C'est dans ce cadre que s'inscrit le présent travail. Parmi les leviers d'amélioration des agents, ingénierie du harnais, scaling du calcul à l'inférence et optimisation paramétrique, nous plaçons l'amélioration intrinsèque de la politique au cœur de la démarche, guidés par l'ambition de maîtriser l'attribution de crédit en environnement multi-tour, où l'espace des trajectoires croît exponentiellement avec l'horizon $T$. Ne pouvant prétendre au passage à l'échelle des environnements accessible aux constructeurs, nous concentrons l'étude sur un environnement contrôlé à récompense vérifiable, TextCraft, issu de la plateforme AgentGym-RL \cite{xi_agentgym-rl_2025}, qui permet d'isoler les variables décisionnelles.


Nos contributions sont au nombre de trois. La première est la réplication du système de référence sur un GPU unique et un guide de bonnes pratiques (je trouve guide de bonne pratique un peu prétentieux, il faudrauit plutot dire un rapport d'exploration des algorithmes organisé par constats de stabilisation intéressants, mais appeler ça un guide serait faux et prétentieux). de stabilisation du GRPO multi-tour, tiré de l'analyse de nos recettes, runs et hyperparamètres, jusqu'au correctif que nous proposons, une ancre KL mobile. La deuxième est l'étude comparée de quatre régimes de curriculum (mais est ce qu'on a défini ce qu'est un curriculum avant ? Ce serait cool d'en parler quelque part avant, d'y dédier un petit paragraphe avec quelques articles) à recette commune, des paliers conçus par l'humain sur l'horizon, la profondeur des tâches et le budget de génération, jusqu'à un autocurriculum piloté par le progrès d'apprentissage prédit. La troisième est l'analyse transverse de ces régimes : transfert entre difficultés et entre planification mono-tour et contrôle interactif, coût à compute contrôlé, et confrontation à l'hypothèse du plafond pass@$k$. Des ablations (ablations pareil c'est trop prétentieux par rapport à ce qu'on a fait) complémentaires, prompting \emph{few-shot} et générations successives de modèles de base, isolent la contribution de chaque levier ; une piste exploratoire non aboutie, la recomposition de trajectoires par échantillonnage préférentiel auto-normalisé (SNIS), est renvoyée en annexe.
De manière générale c'est trop d'énumération on se perd, il faut réécrire.

La suite de ce rapport s'organise comme suit. La section 2 dresse l'état de l'art selon six entrées : la définition des systèmes agentiques et de leurs harnais ; la construction des environnements et benchmarks, à la fois terrain d'évaluation et matière première (plutot dire comme outil de définition des objectifs de l'ia agentique) de l'entraînement ; la stabilité et la mesure du RL en environnement interactif ; les stratégies de structuration de l'apprentissage, du curriculum aux agents autotéliques ; le transfert entre planification mono-tour et contrôle multi-tour ; et les autres pistes considérées. La section 3 présente le cadre technique, la plateforme AgentGym-RL et l'environnement TextCraft, et le positionnement de notre contribution. La section 4 expose le protocole, les expériences et leurs résultats, organisés en stabilisation, curriculums et transfert, comparaison à compute contrôlé, et ablations d'inférence. La section 5 délimite le périmètre et ouvre les perspectives, avant la conclusion. Trois annexes complètent le document : les fondements du Deep RL (annexe A), la dérivation de l'estimateur SNIS et le bilan de cette piste (annexe B), et une analyse critique du parcours de réplication (annexe C).








%-----------------------------------------------------------------------------

\newpage
\section{État de l'art}

\subsection{Définition des Systèmes Agentiques dans la Continuité des LLMs}

L'émergence de l'IA agentique ne se réduit pas à l'ajout d'interfaces d'appel d'outils aux modèles existants. Elle découle en grande partie du goulot d'étranglement des données décrit en introduction : les corpus statiques s'épuisent et les tâches d'évaluation mono-tour ne suffisent plus à discriminer les modèles. Le déploiement des modèles au sein d'environnements numériques interactifs constitue alors la suite logique des travaux sur l'alignement et l'auto-amélioration. Les rapports techniques des constructeurs documentent ce basculement : l'équipe Qwen attribue les gains agentiques de ses dernières générations au passage à l'échelle de la génération d'environnements et de l'apprentissage par renforcement sur ces environnements \cite{qwen3coder2025, qwen35_2026}, et Kimi K2 comme Magistral décrivent des pipelines analogues de synthèse de tâches et d'entraînement RL \cite{kimik2_2025, magistral2025}. La section 2.2 détaille ce constat.

Nous appelons agent un modèle de langage placé dans une boucle d'interaction avec un environnement. À chaque pas, il reçoit une observation, sélectionne une action au sein d'un espace d'outils exposé par cet environnement, puis observe le résultat de cette action avant de décider du pas suivant, jusqu'à terminaison de l'épisode. Le cadre \emph{ReAct} \cite{yao2023react} en constitue l'instanciation canonique : il alterne explicitement pensées (\texttt{Thought}), actions (\texttt{Action}) et observations (\texttt{Observation}) au sein d'une même séquence autorégressive. Ce cadre distingue formellement l'agentique du raisonnement mono-tour, qui évalue le modèle sur un problème statique à information complète, résolu en un seul rollout autorégressif. L'agentique relève au contraire du processus de décision markovien partiellement observable, le cadre usuel de l'apprentissage par renforcement : le modèle n'accède jamais à l'état complet du système et ne dispose que de la trace des observations accumulées dans son contexte. La difficulté réside alors dans la gestion d'un horizon décisionnel long, où chaque action modifie l'état du système et conditionne les observations ultérieures. Qu'il s'agisse de jeux textuels et d'environnements incarnés \cite{shridhar2020alfworld}, d'interfaces web \cite{zhou2023webarena} ou de bases de code réelles \cite{jimenez2023swebench}, ces environnements se caractérisent par leur nature interactive, évolutive et médiée par des outils.

Ce passage au multi-tour ne modifie pas seulement la nature de la tâche, il modifie la difficulté même de l'optimisation. La trajectoire $\tau$ n'est plus intégralement générée par la politique : elle contient des observations produites par l'environnement, sur lesquelles aucun gradient ne doit être rétropropagé. Il faut donc un masquage analogue à celui du SFT, appliqué cette fois à l'échelle de la trajectoire. Surtout, l'horizon $T$ croît de plusieurs ordres de grandeur tandis que la récompense demeure généralement terminale. L'attribution de crédit, c'est-à-dire l'identification des actions responsables du succès ou de l'échec d'un épisode, devient considérablement plus ardue que dans le cas d'une réponse unique. Les recettes de stabilisation validées en RL mono-tour ne se transfèrent donc pas telles quelles au régime interactif. Elles fournissent un point de départ raisonnable, mais doivent y être revalidées ; la section 2.3 précise ces instabilités et la section 4.2 en documente plusieurs exemples.

L'interaction apporte en retour un avantage propre, et la littérature offre des réponses complémentaires à la tension décrite ci-dessus. D'une part, la boucle de rétroaction de l'environnement, par ses messages d'erreur et ses observations d'état, compense partiellement les déficits de planification \emph{a priori} des LLMs \cite{yao2023react, shinn2023reflexion, kambhampati2024llmmodulo} ; la section 2.5 revient sur cette asymétrie entre boucle ouverte et boucle fermée. D'autre part, une partie de la littérature densifie le signal en estimant une récompense à l'échelle de l'étape (\emph{step-level credit}), au moyen de modèles de récompense de processus (\emph{Process Reward Models}, PRM) \cite{lightman2023verify}, au prix d'un coût d'annotation ou d'estimation supplémentaire. C'est précisément cette tension entre horizon, sparsité du signal et attribution de crédit qui motive l'étude des algorithmes de la section 1.2 dans le contexte agentique.

Au sein de cet écosystème, un modèle de langage n'agit jamais de manière isolée : il reste indissociable de son harnais, la superstructure logicielle qui l'entoure, de la gestion du contexte aux parseurs d'actions, en passant par les protocoles de communication comme MCP \cite{anthropic2024mcp} et les mécanismes de vérification. Il convient de distinguer les workflows rigides, où les chemins d'exécution sont codés en dur, des agents autonomes où le modèle dirige lui-même sa propre boucle de décision \cite{anthropic2024agents}. Les notes méthodologiques des benchmarks récents montrent que les performances mesurées dépendent autant de cet environnement d'exécution et du budget de tokens alloué que des poids du modèle \cite{openai2024swebenchverified}. Cette interdépendance conduit les laboratoires à livrer des paires modèle-harnais développées conjointement, à l'image du co-design entre matériel et logiciel, le modèle étant explicitement post-entraîné sur ses propres outils et prompts système \cite{anthropic2024agents, kimik2_2025}. Le harnais est enfin devenu un objet d'optimisation à part entière : les systèmes industriels s'appuient sur l'orchestration dynamique de sous-agents pour paralléliser la décomposition de tâches complexes \cite{kimik2_2025}, et la recherche explore la découverte automatique d'architectures agentiques, où le code du harnais, gestion de la mémoire, filtrage du contexte ou génération de sous-agents, est lui-même appris ou évolué \cite{hu_adas_2025}.

\subsection{Construction des Environnements et Benchmarks Agentiques}

Si la section précédente a défini l'agent comme une politique en boucle d'interaction, la définition du problème qu'il doit résoudre est un enjeu au moins aussi déterminant. C'est elle qui conditionne à la fois ce que l'on mesure et ce que le modèle peut apprendre. Un environnement agentique se spécifie formellement comme un quadruplet : un état initial (le contenu du monde au début de l'épisode), un ensemble d'outils (les actions exposées à l'agent, aujourd'hui de plus en plus médiées par le protocole MCP), une tâche exprimée en langage naturel, et un vérificateur exécutable qui décide du succès. Chacune de ces composantes est un choix de conception lourd de conséquences. Un vérificateur fondé sur l'exécution, par comparaison d'états finaux ou scripts de test, mesure autre chose qu'une comparaison textuelle à une référence. Un état initial pauvre rend la tâche triviale quelle que soit la sophistication des outils.

Les benchmarks académiques illustrent chacun un axe de difficulté distinct, et leur lecture conjointe dessine une décomposition opérationnelle de la compétence agentique. $\tau$-bench et son extension $\tau^2$-bench \cite{taubench2024, tau2bench2025} introduisent la dimension conversationnelle et normative : l'agent doit résoudre la requête d'un utilisateur simulé, lui-même doté d'outils dans le régime de contrôle dual, tout en respectant une politique de domaine (compagnie aérienne, télécoms). Les performances s'effondrent dès que l'utilisateur doit être activement guidé, avec un pass@1 de l'ordre de 34\,\% en télécoms pour des modèles frontière. Toolathlon \cite{toolathlon2025} cible la diversité et le réalisme des états initiaux : 32 applications réelles et 604 outils MCP, 108 tâches d'environ 20 tours nécessitant de coordonner plusieurs applications sur des états de départ authentiques, tels que des cours Canvas peuplés de dizaines d'étudiants ou des feuilles de calcul financières réelles. Le meilleur modèle évalué n'y dépasse pas 38,6\,\% de réussite. MCP-Atlas \cite{mcpatlas2026} et MCP-Universe \cite{mcpuniverse2025} isolent quant à eux la découverte d'outils : les énoncés ne mentionnent ni serveur ni outil, ce qui contraint l'agent à identifier les bons instruments parmi des distracteurs sémantiquement plausibles, puis à composer des flux multi-serveurs. Ces scores, systématiquement inférieurs à 55\,\% pour les modèles frontière, indiquent que le facteur limitant n'est plus seulement la capacité brute de raisonnement. C'est aussi la conjonction entre horizon long, observabilité partielle, ambiguïté de la surface d'outils et vérification stricte, c'est-à-dire précisément ce que la définition de l'environnement encode.

Du côté de l'entraînement, les rapports techniques des constructeurs révèlent que ces environnements ne sont plus seulement des instruments d'évaluation, mais la matière première du post-training. L'équipe Qwen décrit une trajectoire éloquente : une vingtaine de tâches RL généralistes pour Qwen3 \cite{qwen3_2025}, puis 20\,000 environnements d'exécution parallèles pour le RL agentique de Qwen3-Coder, dont le rapport désigne l'\emph{environment scaling} comme le défi central \cite{qwen3coder2025}, jusqu'à Qwen3.5, dont les gains sont attribués au passage à l'échelle de « pratiquement toutes les tâches et tous les environnements de RL concevables », avec une distribution de tâches à complexité progressive \cite{qwen35_2026}. Kimi K2 systématise cette logique par un pipeline de synthèse en trois étapes, de la constitution d'un répertoire de spécifications d'outils réels et synthétiques à la production de trajectoires multi-tours vérifiées, alimentant un étage de RL conjoint sur environnements simulés et réels \cite{kimik2_2025}. Mistral rapporte de même le rôle d'environnements RL dédiés dans l'entraînement de ses modèles de code agentiques \cite{magistral2025}.

Cette insistance sur la diversité n'est pas rhétorique. Elle répond à un mode d'échec bien identifié du RL agentique : une politique optimisée contre un environnement fixe apprend à exploiter les particularités de sa surface d'outils, de son vérificateur ou de ses gabarits de tâches, de sorte que les gains mesurés ne se transfèrent pas hors de la distribution d'entraînement. La multiplicité des environnements ne joue donc pas le rôle d'un simple élargissement de couverture, mais celui d'un régularisateur : c'est la variété des vérificateurs et des dynamiques qui empêche la politique de sur-ajuster un simulateur particulier et force l'apprentissage de compétences décisionnelles transférables. Ce constat ouvre un axe de recherche en plein essor, qui consiste à traiter l'environnement lui-même comme un artefact à générer, qu'un agent propose ses propres tâches et leurs vérificateurs \cite{zhou2025sca} ou qu'un modèle de monde synthétise intégralement les expériences d'entraînement \cite{dreamgym2025}. Il prolonge directement les approches de génération synthétique de données évoquées en introduction. La génération synthétique d'environnements est du reste une pratique établie en apprentissage robotique : l'agent y évolue dans un jumeau numérique, et l'algorithme ajuste ou régénère l'environnement en fonction des performances courantes de l'agent, jusqu'à la co-évolution ouverte d'environnements et de politiques \cite{wang2019poet}. Cette adaptation de l'environnement aux capacités de l'apprenant relie la génération d'environnements au méta-apprentissage et au curriculum, concepts détaillés en section 2.4.

Pour la présente étude, ce panorama fixe deux exigences et délimite un positionnement. Première exigence : comprendre finement l'anatomie de l'environnement de travail, soit pour TextCraft son espace d'actions compositionnel, son vérificateur exact et sa gradation de difficulté par profondeur de recettes. C'est un préalable à toute interprétation des dynamiques d'apprentissage, car un gain mesuré n'a de sens que rapporté à ce que le vérificateur mesure réellement. Seconde exigence : la valeur d'un résultat obtenu dans un environnement contrôlé se juge à ce qui se transfère hors de sa distribution, faute de quoi l'on ne mesure que l'ajustement à un simulateur. Quant au positionnement, notre sujet n'est pas de générer des environnements, mais de comprendre comment optimiser l'apprentissage d'un agent en leur sein, et d'y étudier les idées d'auto-amélioration et de méta-apprentissage. La génération d'environnements à grande échelle relève avant tout d'une ingénierie logicielle que seuls les constructeurs peuvent soutenir ; l'étude des dynamiques d'apprentissage dans ces environnements constitue la question de recherche, et c'est elle qui correspond au sujet du stage. Ce double constat motive le choix d'un cadre expérimental délibérément restreint mais finement instrumenté (section 3), et l'attention portée aux propriétés de stabilisation et de transfert plutôt qu'aux seuls scores absolus (section 4).

\subsection{Apprentissage par Renforcement en Environnement Interactif : Stabilité et Mesure}

Les algorithmes présentés en section 1.2 ont été conçus et validés dans des régimes mono-tour : contrôle continu pour PPO, alignement RLHF puis raisonnement RLVR pour ses variantes appliquées aux LLMs. Leur transposition aux environnements interactifs est récente, et deux questions transversales structurent la littérature comme notre travail : comment maintenir la stabilité de cet apprentissage, et comment mesurer ce qu'il apprend réellement. Cette sous-section pose ces deux points, que les sections 4.2 et 4.5 développeront expérimentalement.

Le RL profond à récompense sparse est notoirement instable, et le régime multi-tour aggrave chacune de ses pathologies connues. Le détournement de récompense (\emph{reward hacking}) désigne le cas où la politique optimise la récompense mesurée en divergeant de l'intention du concepteur \cite{skalse2022defining} ; les récompenses intermédiaires ajoutées pour densifier le signal en sont un vecteur classique. La dérive de politique désigne l'éloignement progressif de la politique entraînée par rapport à ses capacités d'origine ; la pénalité de divergence KL vers une politique de référence, héritée du RLHF \cite{ouyang2022training}, en est le garde-fou standard, mais le choix de la référence et du coefficient devient critique sur les longs horizons. Enfin, le collapse d'entropie désigne l'effondrement précoce de la diversité des sorties de la politique : Cui et al. \cite{cui2025entropy} montrent que l'entropie chute dès le début de l'entraînement RL, que la performance se paie en entropie, et que la saturation de l'une accompagne l'épuisement de l'autre. Sans intervention, la politique devient quasi déterministe et l'exploration s'éteint, ce qui prive les algorithmes de groupe comme GRPO de tout signal. Ces trois pathologies ne sont pas des curiosités théoriques : nous les avons toutes rencontrées, et la section 4.2 présente les pratiques de stabilisation que cette étude a permis d'identifier, du masquage des observations à l'ancrage KL mobile.

Mesurer ce que le RL apprend exige ensuite un instrument plus fin que le taux de réussite en un tirage. Le pass@$k$, fraction des tâches résolues par au moins un tirage parmi $k$, approche pour $k$ grand le support de la politique, c'est-à-dire l'ensemble des problèmes que le modèle sait résoudre au moins occasionnellement. Cet instrument fonde une thèse forte sur la nature du RLVR : Yue et al. \cite{yue2025rlvr} observent que les modèles entraînés dominent leur modèle de base à petit $k$ mais lui deviennent inférieurs à grand $k$, et concluent que le RL concentre la masse de probabilité sur les chemins déjà récompensés sans étendre la frontière de compétence, en écho au mécanisme de \emph{sharpening} \cite{huang2024self} présenté en introduction. La question de l'attribution s'y ajoute : ce qu'un post-entraînement apporte dépend de ce que le modèle de base contient déjà, le déclenchement de capacités dites émergentes étant conditionné par la présence préalable des structures correspondantes \cite{liu2025drgrpo}, et les générations récentes de modèles ouverts étant massivement exposées aux traces agentiques dès le mid-training \cite{qwen3coder2025, kimik2_2025}. Conformément au positionnement de ce rapport, nous traitons la thèse du plafond pass@$k$ comme une hypothèse à mettre à l'épreuve et non comme un acquis ; la section 4.5 la confronte à nos mesures, ainsi qu'à nos ablations inter-générations (Qwen2.5, Qwen3, Qwen3.5).

\subsection{Curriculum Learning, Autocurriculum et Agents Autotéliques}

\textcolor{red}{Figure illustrative avec les batch randoms de task et les batchs organisés par difficulté par exemple}

Le \emph{curriculum learning} désigne une stratégie d'entraînement dans laquelle les exemples sont présentés au modèle par difficulté croissante, plutôt que dans un ordre arbitraire \cite{bengio2009curriculum}. Dans le cadre du réentraînement par renforcement d'agents LLM, cette idée prend un sens très concret. La récompense y est terminale et binaire : sur une tâche trop difficile pour la politique courante, tous les rollouts d'un groupe échouent, l'avantage relatif est nul et GRPO ne produit aucun gradient ; sur une tâche déjà maîtrisée, tous réussissent et le gradient est nul également. L'ordre dans lequel les tâches sont présentées détermine donc directement la quantité de signal d'apprentissage disponible à chaque étape. C'est ce qui fait du curriculum un levier candidat pour stabiliser et accélérer l'apprentissage face à des horizons longs et des récompenses sparses, et la littérature récente des agents auto-améliorants le répertorie comme l'un de ses paradigmes constitutifs \cite{gao2025selfevolving}.

Une première famille de curriculums est conçue par l'utilisateur (\emph{user-designed}) : l'ordre de présentation des tâches ou les capacités du modèle suivent une heuristique de difficulté fixée avant l'entraînement. La méthode ScalingInter d'AgentGym-RL \cite{xi_agentgym-rl_2025} en est un exemple sur l'axe de l'horizon : le nombre maximal de tours d'interaction est d'abord restreint, ce qui force la politique à exploiter des trajectoires courtes, puis étendu par paliers. Le séquencement par complexité du domaine en est un second, sur l'axe de la tâche elle-même : dans TextCraft, entraîner d'abord sur les recettes de faible profondeur, puis élargir progressivement aux profondeurs supérieures. Ces schémas ont l'avantage de la simplicité et de la lisibilité. Ils ont aussi deux limites. La difficulté qu'ils encodent est celle perçue par le concepteur humain, qui ne coïncide pas nécessairement avec la difficulté effective pour le modèle. Et leur calendrier est fixé à l'avance, aveugle aux dynamiques d'apprentissage réelles : un palier franchi trop tôt prive la politique de signal, un palier maintenu trop longtemps la fait stagner sur des tâches acquises.

Ces limites motivent la seconde famille, l'\emph{autocurriculum}, où la sélection des tâches est ajustée en continu à partir des performances de la politique elle-même. L'idée fondatrice est de sélectionner les tâches selon le progrès d'apprentissage (\emph{learning progress}) qu'elles procurent, c'est-à-dire la vitesse à laquelle la compétence du modèle y évolue \cite{graves2017automated}. L'intuition rejoint la notion de zone proximale de développement : les tâches les plus informatives sont celles de difficulté intermédiaire, ni acquises ni hors de portée, précisément celles où les groupes de rollouts sont mixtes et où GRPO reçoit du signal. Cette approche a récemment reçu des garanties théoriques dans le contexte des LLMs : un autocurriculum réduit de manière exponentielle le besoin de démonstrations supervisées en SFT, et découple le coût du RL de la qualité du modèle de référence \cite{rajaraman2026curriculum}.

L'autocurriculum suppose encore un ensemble de tâches donné, dans lequel l'agent choisit. Le cadre des agents autotéliques franchit un pas de plus : l'agent y définit, poursuit et organise ses propres objectifs au sein d'environnements ouverts \cite{colas2022autotelic, forestier2017imgep}. La sélection de buts devient alors le cœur du problème, et elle se heurte à une difficulté d'échelle : lorsque l'espace des buts est très vaste, estimer la compétence de l'agent sur chaque but par comptage de ses succès devient impraticable, chaque but n'étant visité que rarement. Le framework MAGELLAN \cite{gaven2025magellan} répond à cette difficulté par une approche métacognitive : un estimateur appris, branché sur les représentations internes du LLM, prédit la probabilité de succès de l'agent sur chaque but, et le progrès d'apprentissage est mesuré comme l'évolution de cette prédiction dans le temps. Parce que des buts sémantiquement proches ont des représentations proches, l'estimateur généralise entre buts et fournit une carte de compétence dense à partir d'observations rares. C'est ce mécanisme que nous portons et évaluons en section 4. La même logique de sélection à la frontière des capacités s'étend aux architectures hiérarchiques : dans HERAKLES \cite{carta2025herakles}, une politique haut niveau choisit des sous-buts atteignables que la politique bas niveau apprend à compiler en compétences réutilisables, ce qui constitue un curriculum appliqué aux buts que l'agent se donne lui-même.

Notre étude se situe à la jonction des deux premières familles et emprunte à la troisième son instrument de mesure. Nous comparons, à recette d'entraînement identique, des curriculums user-designed portant sur deux axes orthogonaux de TextCraft, l'horizon d'interaction et la profondeur des recettes, à un autocurriculum fondé sur le progrès d'apprentissage prédit, via notre port de MAGELLAN. Cette comparaison instruit une question que les schémas manuels laissent ouverte : la difficulté perçue par l'humain, ici la profondeur, est-elle la vraie structure de compétence du modèle, ou la machine découvre-t-elle un autre découpage ? La perspective plus lointaine, un générateur d'environnements dont la difficulté suivrait les capacités courantes du modèle, rejoint l'axe de génération d'environnements décrit en section 2.2 ; elle est discutée en ouverture (section 5).

\subsection{Transfert de Compétences : Planification Mono-tour et Contrôle Interactif Multi-tour}

\textcolor{red}{image papier lecun meta learning ?}

Le transfert entre la planification mono-tour en boucle ouverte (\emph{open-loop planning}) et l'exécution interactive multi-tour en boucle fermée (\emph{closed-loop control}) représente un enjeu à la fois économique et algorithmique. D'un point de vue computationnel, l'entraînement par renforcement en mono-tour est nettement plus frugal : il ne nécessite qu'un seul rollout autorégressif par problème, ce qui supprime les dizaines d'allers-retours avec le serveur d'environnement qui dominent le temps de chaque étape d'optimisation en multi-tour. Si des compétences de planification abstraite acquises en mono-tour se transféraient au contrôle interactif, on disposerait d'un régime de pré-alignement à bas coût, susceptible d'accélérer la convergence du RL multi-tour et d'en réduire la facture de calcul.

La littérature sur les capacités de planification des LLMs invite toutefois à la prudence sur le premier terme de ce transfert. Les travaux de Valmeekam, Kambhampati et leurs collaborateurs \cite{valmeekam2023planning, kambhampati2024llmmodulo} établissent que la génération autonome de plans complets et valides reste une faiblesse structurelle des LLMs : le raisonnement transitif sur de longues chaînes de dépendances se dégrade rapidement avec la profondeur du problème. Symétriquement, la littérature du contrôle interactif, de ReAct \cite{yao2023react} aux approches à rétroaction verbale \cite{shinn2023reflexion}, montre que la boucle de retour de l'environnement compense partiellement ces déficits de planification \emph{a priori}. Les messages d'erreur et l'observation de l'état agissent comme un support dynamique que le plan abstrait ne peut fournir seul.

La question du transfert est donc doublement ouverte. Dans le sens mono-tour vers multi-tour : un pré-entraînement en planification accélère-t-il le contrôle interactif ? Dans le sens inverse : le RL interactif améliore-t-il la planification abstraite ? On peut également concevoir des régimes hybrides alternant phases de planification en boucle ouverte et phases d'exécution interactive. La section 4 instrumente et quantifie ces deux directions de transfert sur TextCraft.

Ce couplage entre apprentissage par observation et apprentissage par action dépasse enfin notre cas d'étude et rejoint une question posée par les sciences cognitives. Dupoux, LeCun et Malik \cite{dupoux2026why} partent du même constat de mur des données que notre introduction, et proposent une architecture d'apprentissage autonome inspirée du développement humain : un système qui apprend par observation, un système qui apprend par action en environnement, et un méta-contrôleur qui orchestre les deux. La planification mono-tour en boucle ouverte et le contrôle interactif en boucle fermée en constituent une instanciation naturelle dans le cadre des agents LLM. Un régime qui alternerait ces deux phases et dont l'orchestration serait elle-même apprise relèverait du méta-apprentissage, au même titre que l'autocurriculum de la section 2.4. Nous n'avons pas eu le temps de développer cette direction au-delà des deux mesures de transfert présentées en section 4 ; elle constitue à nos yeux l'un des prolongements les plus naturels de ce travail.

\subsection{Autres Pistes Considérées}

Deux directions complémentaires ont été étudiées au cours du stage, puis dépriorisées au profit du curriculum ; elles ne relèvent pas directement du méta-apprentissage et leur développement est renvoyé en annexe. La première est le \emph{scaling} du calcul à l'inférence : plutôt que d'entraîner davantage, allouer plus de calcul au moment de la génération. Ce levier va de l'apprentissage en contexte, qui adapte le comportement du modèle par simple injection d'exemples sans modification des poids \cite{brown2020gpt3}, aux stratégies d'échantillonnage multiple et de recherche, dont les analyses de Snell et al. \cite{snell2024scaling} fournissent le cadre d'arbitrage : à coût de calcul égal, un modèle plus petit doté d'une stratégie d'inférence adaptée à la difficulté peut surpasser un modèle plus grand interrogé naïvement. Nos expériences few-shot en section 4 relèvent de ce levier, et la frontière entre calcul d'inférence et génération de données d'entraînement s'estompe dans les systèmes générateurs de tâches évoqués en section 2.2 \cite{zhou2025sca, dreamgym2025}.

La seconde piste est la recomposition de trajectoires par échantillonnage préférentiel auto-normalisé (\emph{Self-Normalized Importance Sampling}, SNIS) \cite[chap.~9]{owen2013montecarlo}, formalisée à la suite d'échanges avec Otmane Sakhi (Criteo AI Lab, équipe Performance Science). L'idée consiste à réutiliser hors-politique les sous-composantes des trajectoires d'un même groupe GRPO, en recomposant chaînes de pensée et réponses finales repondérées par des poids auto-normalisés ; le mécanisme exploite le fait qu'évaluer la vraisemblance d'une recomposition coûte une seule passe avant, là où générer une trajectoire nouvelle est coûteux. La transposition au multi-tour se heurte à un verrou structurel que nos essais ont confirmé : l'état intègre l'historique strict des observations du simulateur, et recoller des sous-trajectoires issues de contextes discordants viole la cohérence causale. La formalisation des poids, la dérivation de l'estimateur et le bilan de nos observations sont regroupés en annexe B, avec la variante guidée par un modèle réviseur comme condition d'une reprise éventuelle.











%------------------------------------------------------------------------------------------------------------








\newpage
\section{Cadre Technique : AgentGym-RL et TextCraft}

\textcolor{red}{ajouter quelques figures du papier}
\subsection{L'Écosystème AgentGym : une Plateforme Unifiée d'Environnements}

Les exigences dégagées en section 2, diversité des environnements, vérificateurs exécutables et mesure du transfert, plaident pour un cadre expérimental mutualisé plutôt que pour le développement d'un simulateur ad hoc. Ce travail s'appuie sur l'écosystème AgentGym, construit en deux temps par la même équipe (Fudan NLP). Le premier papier, AgentGym \cite{xi_agentgym_2025}, publié à ACL 2025, répond à l'absence de plateforme unifiée pour \emph{évaluer et entraîner} des agents LLM. Il fédère 14 environnements couvrant 89 tâches réparties en sept grands scénarios, de la navigation web (WebShop, WebArena) aux jeux textuels et de crafting (TextCraft, MAZE, Wordle), en passant par la planification incarnée (ALFWorld, BabyAI), le raisonnement scientifique (SciWorld) et l'utilisation d'outils applicatifs (Weather, Movie, Academia, Sheet, TODOList, BIRD), derrière une API HTTP uniforme qui découple strictement l'agent des serveurs d'environnements. Au-delà de l'infrastructure, la contribution comprend une base d'instructions élargie, une suite d'évaluation (AgentEval), un corpus de trajectoires expertes (AgentTraj) et une méthode d'auto-amélioration, AgentEvol, qui fait évoluer l'agent au-delà des démonstrations par exploration et imitation filtrée de ses propres succès, héritière directe des approches de type ReST présentées en introduction. L'ambition affichée est celle d'agents \emph{généralistes} : un même modèle entraîné et évalué sur l'ensemble des environnements de la suite.

Le second papier, AgentGym-RL \cite{xi_agentgym-rl_2025}, accepté comme présentation orale à ICLR 2026 \cite{xi_agentgymrl_openreview}, prolonge cette infrastructure vers l'apprentissage par renforcement multi-tour de bout en bout. Le constat de départ rejoint exactement notre cadrage : la communauté manque d'un cadre RL interactif unifié capable d'entraîner des agents \emph{from scratch}, sans étape initiale de SFT sur démonstrations expertes, coûteuses et peu généralisables. L'architecture modulaire et découplée permet des rollouts distribués, concurrents et asynchrones, condition du passage à l'échelle du Deep RL en environnement interactif, et supporte les principaux algorithmes d'optimisation de politique (PPO, GRPO, variantes). L'entraînement RL porte sur un sous-ensemble de cinq scénarios et 27 tâches : navigation web, recherche approfondie, tâches incarnées et sciences, jeux dont TextCraft (Minecraft Textuel). Pour contenir l'instabilité du RLVR sur longs horizons, les auteurs introduisent ScalingInter-RL, un curriculum sur la durée d'interaction : l'entraînement débute avec un nombre maximal de tours restreint, ce qui force l'exploitation sur trajectoires courtes, puis l'horizon est progressivement étendu à mesure que la politique se stabilise. Les résultats sont marquants : un Qwen2.5-7B entraîné par ScalingInter-RL atteint 91\,\% de réussite moyenne sur la suite, égalant ou dépassant des modèles propriétaires de premier plan (GPT-4o, Gemini-2.5-Pro, o3) évalués sur les mêmes tâches \cite{xi_agentgym-rl_2025}.

\subsection{Positionnement Critique et Justification du Choix}

Ces résultats appellent cependant une lecture nuancée, qui délimite précisément l'espace de notre contribution. D'abord, la comparaison oppose de petits modèles optimisés directement sur la distribution exacte des environnements d'entraînement à des modèles généralistes évalués à froid, sans adaptation au format d'action ni au comportement du simulateur. Cette asymétrie d'évaluation relativise l'équivalence affichée. La démonstration de faisabilité reste néanmoins économiquement significative pour le déploiement de modèles spécialisés frugaux : concentrer la politique d'un modèle léger sur de longues séquences décisionnelles, avec un transfert partiel vers des tâches non vues. Ensuite, le processus d'évaluation par les pairs, consultable publiquement \cite{xi_agentgymrl_openreview}, situe la contribution davantage du côté de l'infrastructure que de l'algorithmique : ScalingInter-RL est une heuristique de curriculum fixée a priori, dont le gain sur un entraînement à horizon constant varie selon les tâches, et dont le mécanisme n'est ni analysé ni comparé à des alternatives adaptatives. Il faut enfin mesurer l'asymétrie de moyens : ce travail collectif mobilise plusieurs dizaines d'auteurs et une infrastructure multi-GPU, quand la présente étude est une réplication et une extension menées par une seule personne sur un GPU unique. C'est cette contrainte qui donne sa valeur pratique au guide de stabilisation de la section 4.2.

Cette lecture critique n'épuise pourtant pas l'intérêt de l'idée. Faire porter le curriculum sur un paramètre du harnais, ici le budget d'interaction, plutôt que sur la sélection des données elles-mêmes, reste rare dans la littérature du curriculum, qui se concentre généralement sur l'ordre des tâches. L'idée nous semble généralisable au-delà de AgentGym-rl : à coût de calcul égal, des épisodes courts seront toujours préférables, et l'horizon utile sera toujours un bon proxy de difficulté des tâches, de sorte que la règle est plausiblement applicable à tout environnement interactif. C'est d'ailleurs le nombre de tours d'intéraction avec l'environement effectué par Claude dans le papier Toolathlon qui permet aux auteurs de qualifier la difficulté de leurs tâches. Nos expériences la prolongent d'ailleurs sur un second paramètre du même type, le budget de génération par tour (section 4).

ScalingInter-RL n'est qu'un point dans l'espace des régimes de structuration de l'apprentissage décrit en section 2.4 : un curriculum \emph{user-designed}, statique, portant sur une seule variable, l'horizon. L'environnement TextCraft en expose naturellement une seconde, la profondeur des recettes, et la littérature autotélique fournit le cadre d'une troisième famille, adaptative, pilotée par le progrès d'apprentissage. Comparer systématiquement ces régimes à recette d'entraînement identique et à compute contrôlé, ce que le papier ne fait pas, constitue le cœur de notre contribution. Par ailleurs, la plateforme satisfait nos exigences méthodologiques : récompense vérifiable exacte, simulateur peu coûteux autorisant des ablations répétées, et multiplicité d'environnements offrant, à terme, le terrain de mesure du transfert hors distribution exigé en section 2.2.

\subsection{TextCraft : Anatomie du Banc d'Essai, Dataset et Questions Ouvertes}

TextCraft \cite{prasad2023adapt} est un jeu de crafting textuel inspiré de Minecraft, introduit initialement pour évaluer la décomposition de tâches par les LLMs. L'agent dispose de trois actions élémentaires : \texttt{get} pour récupérer un ingrédient de base, \texttt{craft} pour combiner des ingrédients selon une recette, et \texttt{inventory} pour consulter son inventaire. Il doit fabriquer un objet cible dont la recette dépend récursivement d'autres objets. Chaque tâche se caractérise par la \emph{profondeur} de son arbre de dépendances : à profondeur 1, l'objet se fabrique directement à partir d'ingrédients de base ; à profondeur $d$, il requiert la fabrication préalable d'objets de profondeur $d-1$. La figure~\ref{fig:textcraft_episode} présente un épisode annoté, avec le prompt système, l'observation initiale listant les commandes de craft disponibles, et l'alternance de pensées et d'actions de l'agent.

\begin{figure}[htbp]
\centering
\begin{tcolorbox}[
    colback=gray!4, colframe=black!65,
    title=\textbf{Exemple de tâche TextCraft --- prompt système et objectif},
    fonttitle=\small, width=0.96\textwidth,
    boxrule=0.6pt, arc=2mm, left=3mm, right=3mm, top=2mm, bottom=2mm]
{\small\ttfamily
You are given few useful crafting recipes to craft items in Minecraft. Crafting commands are of the format \mbox{"craft} [target object] using [input \mbox{ingredients]"}. Every round I will give you an observation, you have to respond an action based on the state and instruction. You can "get" an object (ingredients) from the inventory or the environment, look-up the game inventory by "inventory", or "craft" (target) using any of the crafting commands. You can use ONLY these crafting commands provided, do not use your own crafting commands. However, if the crafting command uses a generic ingredient like "planks", you can use special types of the same ingredient e.g. "dark oak planks" in the command instead.

\medskip
\textbf{Goal: craft flint and steel.}
}
\end{tcolorbox}
\caption{Prompt système de l'environnement TextCraft et objectif type, tels que fournis à l'agent en début d'épisode (exemple tiré de \cite{xi_agentgym-rl_2025}). L'agent répond ensuite tour par tour par une action \texttt{get}, \texttt{inventory} ou \texttt{craft}, l'environnement renvoyant une observation après chaque action, jusqu'à vérification de l'objet cible dans l'inventaire.}
\label{fig:textcraft_episode}
\end{figure}


L'instance utilisée compte 544 recettes, que nous partitionnons en trois ensembles : 374 recettes d'entraînement, 100 recettes de test et un réservoir de 70 recettes réservées à la construction des exemples few-shot, ce qui garantit l'absence de fuite entre démonstrations et évaluation. La distribution par profondeur est fortement déséquilibrée : l'entraînement compte 109 recettes de profondeur 1, 221 de profondeur 2, 43 de profondeur 3 et une seule de profondeur 4 ; le test en compte respectivement 31, 41, 25 et 3. Ce déséquilibre est une limite du dataset que le papier ne discute pas : avec une seule recette de profondeur 4 à l'entraînement, l'échec éventuel d'une méthode sur les tâches profondes est confondu avec la quasi-absence de support d'entraînement. Nous n'avons pas re-stratifié cette distribution, ce qui aurait constitué une étude à part entière, mais nous la gardons présente dans toutes nos interprétations, et la re-stratification par difficulté figure parmi les travaux futurs les plus directement actionnables (section 5).

Ce banc d'essai concentre, sous forme minimale, les propriétés identifiées en section 2.2 : un espace d'actions compositionnel dont la difficulté est \emph{paramétrable par construction}, la profondeur étant une variable de curriculum naturelle ; un vérificateur exact, l'objet cible figurant ou non dans l'inventaire final ; une observabilité partielle réelle, l'agent devant consulter l'inventaire et découvrir les recettes ; et un horizon ajustable. Sa légèreté computationnelle, sans rendu et à dynamique déterministe, autorise les ablations répétées qu'exigerait difficilement WebArena. Ces propriétés en font l'environnement adapté pour disséquer la dynamique des gradients GRPO avant extension aux autres environnements de la suite, extension envisagée si les régimes de curriculum se montrent concluants, à la manière de la validation multi-environnements de ScalingInter chez les auteurs.

Le tableau~\ref{tab:agentgym_textcraft} reproduit les résultats de référence du papier sur les 100 tâches de test, qui constituent nos lignes de base externes. Sa lecture par profondeur est la plus instructive. La profondeur 4 est un mur quasi universel : GPT-4o, o3 et la totalité des modèles ouverts de moins de 100 milliards de paramètres antérieurs à Qwen3 y échouent intégralement, et seuls Gemini-2.5-Pro, Qwen3-8B, Qwen3-32B et ScalingInter-7B y résolvent une tâche sur trois. La profondeur 3 sépare nettement les générations et les tailles de modèles. Enfin, la comparaison des trois modèles RL du papier montre que l'essentiel du gain de ScalingInter-7B sur AgentGym-RL-7B se joue en profondeur 3 et 4. Précisons une réserve de comparabilité : ces chiffres dépendent du harnais et du prompt d'évaluation des auteurs ; notre propre mesure zéro-shot de Qwen2.5-3B-Instruct sur les mêmes 100 tâches donne 18\,\% avec notre harnais, contre 14\,\% rapportés ici. Les comparaisons externes sont donc indicatives ; toutes nos comparaisons internes utilisent un harnais unique.

Trois questions ouvertes, posées en section 2, prennent alors leur forme opérationnelle. Le curriculum améliore-t-il la vitesse de convergence, le plafond de performance, ou les deux ? Peut-on atteindre les profondeurs 3 et 4 par une gradation appropriée, malgré le support d'entraînement quasi nul en profondeur 4 ? Et les pass@$k$ du modèle nu constituent-ils un plafond infranchissable pour le RL, comme le soutient la thèse du sharpening, ou une bonne structuration de l'apprentissage permet-elle de les dépasser ?


\textcolor{red}{la tabe est archi moche, reprendre le design du paper}
\begin{table}[htbp]
\centering
\caption{Résultats de référence sur TextCraft (100 tâches de test, succès en \%, par profondeur $d$ et global), d'après la Table~3 de \cite{xi_agentgym-rl_2025}.}
\label{tab:agentgym_textcraft}
\begin{tabular}{lccccc}
\toprule
Modèle & $d{=}1$ & $d{=}2$ & $d{=}3$ & $d{=}4$ & Global \\
\midrule
\multicolumn{6}{l}{\emph{Modèles propriétaires}} \\
GPT-4o & 100 & 87{,}80 & 64{,}00 & 0 & 83{,}00 \\
Qwen-Max & 93{,}55 & 75{,}61 & 36{,}00 & 0 & 69{,}00 \\
Gemini-2.5-Flash & 100 & 95{,}12 & 40{,}00 & 0 & 80{,}00 \\
OpenAI o4-mini & 100 & 100 & 84{,}00 & 0 & 93{,}00 \\
OpenAI o3 & 100 & 100 & 84{,}00 & 0 & 93{,}00 \\
Gemini-2.5-Pro & 100 & 100 & 84{,}00 & 33{,}33 & 94{,}00 \\
\midrule
\multicolumn{6}{l}{\emph{Modèles ouverts $\geq$ 100B}} \\
Qwen3-235B-A22B & 100 & 100 & 84{,}00 & 0 & 93{,}00 \\
DeepSeek-V3-0324 & 80{,}65 & 53{,}66 & 40{,}00 & 0 & 57{,}00 \\
DeepSeek-R1-0528 & 100 & 100 & 84{,}00 & 0 & 93{,}00 \\
\midrule
\multicolumn{6}{l}{\emph{Modèles ouverts $<$ 100B}} \\
Qwen2.5-3B-Instruct & 35{,}48 & 7{,}32 & 0 & 0 & 14{,}00 \\
Qwen2.5-7B-Instruct & 80{,}65 & 41{,}46 & 0 & 0 & 42{,}00 \\
Qwen2.5-72B-Instruct & 96{,}77 & 85{,}37 & 48{,}00 & 0 & 77{,}00 \\
Qwen3-4B & 87{,}10 & 36{,}59 & 12{,}00 & 0 & 45{,}00 \\
Qwen3-8B & 100 & 78{,}05 & 40{,}00 & 33{,}33 & 74{,}00 \\
Qwen3-32B & 90{,}32 & 92{,}68 & 72{,}00 & 33{,}33 & 85{,}00 \\
Llama-3.1-8B-Instruct & 74{,}19 & 56{,}10 & 4{,}00 & 0 & 47{,}00 \\
Llama-3.1-70B-Instruct & 100 & 100 & 84{,}00 & 0 & 93{,}00 \\
\midrule
\multicolumn{6}{l}{\emph{Modèles RL du papier}} \\
AgentGym-RL-3B & 100 & 90{,}24 & 28{,}00 & 0 & 75{,}00 \\
AgentGym-RL-7B & 100 & 97{,}56 & 72{,}00 & 0 & 89{,}00 \\
ScalingInter-7B & 100 & 97{,}56 & 76{,}00 & 33{,}33 & 91{,}00 \\
\bottomrule
\end{tabular}
\end{table}




%----------------------------------------------------------------------------------------
\newpage
\section{Expériences et Résultats}

Cette section présente le travail expérimental du stage et ses résultats. L'ordre d'exposition est logique plutôt que chronologique. 

La sous-section 4.1 fixe lecadre et protocole commun à toutes les expériences.

La sous-section 4.2 établit les conditions de stabilité du RL multi-tour et en tire un guide de bonnes pratiques : c'est la première contribution. \textcolor{red}{stop guide, parler d'étude de stabilisation et de ce qu'on a vraiment à partager : LoRA}

Les sous-sections 4.3 et 4.4 portent le cœur de l'étude, la comparaison des régimes de curriculum à recette identique et à compute contrôlé, ainsi que les mesures de transfert entre planification mono-tour et contrôle interactif \textcolor{red}{Non pas de transfert mono multi, à mettre dans ouvertures}: ce sont les deuxième et troisième contributions. La sous-section 4.5 regroupe enfin les expériences d'inférence et les ablations transverses, dont la confrontation de nos résultats à la thèse du plafond pass@$k$ posée en section 2.3 \textcolor{red}{lire le papier en question rapidement}. 







\subsection{Protocole Commun}\label{subsec:protocole}

\paragraph{Environnement et données.} Toutes les expériences portent sur TextCraft, dont l'anatomie et les partitions sont détaillées en section 3.3 : 374 recettes d'entraînement, 100 recettes de test fixes et un réservoir de 70 recettes réservé aux exemples few-shot, sans fuite entre démonstrations et évaluation. La distribution par profondeur de ces partitions, fortement déséquilibrée (une seule recette de profondeur 4 à l'entraînement, \textcolor{red}{à assumer pleinement en section 3 comme limitation de l'étude mais non étudiée ici}), est rappelée dans chaque analyse stratifiée. Le serveur d'environnement est exécuté localement ; l'agent y accède par l'API HTTP d'AgentGym. \textcolor{red}{Ajouter free contribution : changer le dataset avec toutes les données dans train}

\paragraph{Modèles.} Le modèle principal est Qwen2.5-3B-Instruct, la cible du papier de référence, sur lequel portent tous les entraînements. Qwen3.5-4B sert de contrôle inter-générations : nous l'évaluons sans aucun entraînement, avec le même harnais. Pour la génération intermédiaire Qwen3, nous n'avons pas de mesure interne et nous appuyons sur la borne externe du papier (Qwen3-4B : 45\,\% sur son harnais, tableau~\ref{tab:agentgym_textcraft}) \textcolor{red}{mettre directement la bonne capture d'écran de la table ou reprendre le style avec zone grise et best en gras}, avec la réserve de comparabilité énoncée en section 3.3.

\paragraph{Optimisation.} \textcolor{red}{Globalement clair mais énumérations ricches, il faudra rappeler de lire les parties précédentes pour bien comprendre ou mettre des liens vers les sections}L'algorithme est GRPO dans tous les cas : groupes de $G$ rollouts par tâche, avantage relatif intra-groupe, récompense terminale binaire, masquage des tokens d'environnement (\texttt{env\_mask}) et pénalité KL d'ancrage. Deux recettes de référence structurent la section. La première est la réplication de la recette publiée du papier en full finetuning, à hyperparamètres identiques aux leurs : $G=8$, taux d'apprentissage de $1 \times 10^{-6}$, coefficient KL de $1 \times 10^{-3}$ vers une référence fixe (Qwen 2.5 base), température d'échantillonnage de $1{,}0$. La seconde est la recette stabilisée que la section 4.2 construite pas à pas et justifiée choix par choix : adaptation LoRA de petit rang, taux d'apprentissage constant plus élevé, coefficient KL renforcé, et ancre de référence mobile, re-basée périodiquement sur la politique courante. Au sein de la phase curriculum, la taille de groupe passe de $G=8$ à $G=16$ à nombre de trajectoires par mise à jour constant (64) \textcolor{red}{permis par la mémoire libérée par lora}: un banc de débit GPU a montré qu'à 4 tâches par mise à jour au lieu de 8, la génération est environ 18\,\% plus rapide et l'estimation d'avantage par tâche plus précise. La section 4.3 contrôle explicitement l'effet de ce changement \textcolor{red}{chiffres à vérifier}.

\paragraph{Évaluation.} La métrique principale est le pass@$1$ sur les 100 tâches de test, mesuré périodiquement en cours d'entraînement toutes les 47 \textcolor{red}{justifier ce chiffre}mises à jour, par des épisodes complets aux conditions fixes : 30 tours maximum, température $1{,}0$, top-$p$ $1{,}0$ \textcolor{red}{pas compris}, 512 tokens par tour. L'incertitude associée est celle d'une proportion estimée sur $n=100$ tâches : l'écart-type binomial vaut $\sqrt{p(1-p)/n}$, soit au plus 5 points autour de $p = 0{,}5$, ce que confirment les fluctuations observées sur les plateaux de nos courbes \textcolor{red}{pas compris ce calcul d'estimation d'incertitude}. Les écarts entre configurations sont donc systématiquement rapportés à cet écart-type. Le meilleur score d'une courbe appelle la même prudence : c'est un maximum pris sur plusieurs dizaines d'évaluations bruitées, statistique d'extrême qui surestime le niveau réel de l'ordre d'un à deux écarts-types. Nous rapportons en conséquence chaque meilleur score accompagné de la plage des évaluations voisines, et nous avons vérifié empiriquement ce biais : une reprise d'entraînement initialisée sur un pic isolé retombe au niveau des évaluations qui l'entouraient. \textcolor{red}{appuyer sur notre objectivité de résultats et incertitude : on ne parle pas uniquement de chiffres max bruts sans explication du filtrages des données ou de l'incertitude du max. On déconstruit et explore ces méthodes et métriques pas assez transparentes.}

\paragraph{Métriques stratifiées et diagnostiques.} Chaque évaluation périodique enregistre, au-delà du pass@$1$ global, un ensemble de métriques journalisées \textcolor{red}{que veut dire journaliser ? } systématiquement pour toutes les expériences. Cela permet, au delà de la performance brute sur un set de tâches non stratifiées autant en train qu'en test, d'apprécier l'évolution du comportement du modèle par exemple sur les types d'erreurs effectuées, en fonction de la stratégie de curriculum et des hyperparamètres choisis. Côté tâche : le pass@$1$ stratifié par profondeur ($d = 1$ à $4$), le nombre moyen de tours par épisode, et une taxonomie d'erreurs par épisode et par dificulté pour distinguer principalement les erreurs de format des erreurs de raisonnement (erreur de format d'action, recette erronée, ingrédient manquant, objet introuvable, actions multiples dans un même tour). Côté dynamique d'entraînement, chaque mise à jour journalise la moyenne et l'écart-type du reward par groupe, la fraction de groupes à variance nulle, donc sans gradient, l'entropie de la politique, la divergence KL vers l'ancre, les longueurs de génération et la norme du gradient. Ces séries alimentent trois analyses de la suite : l'ordre d'acquisition des profondeurs au fil de l'entraînement, l'ordre d'élimination des types d'erreurs, et les signaux avant-coureurs de collapse (section 4.2). Le support des politiques est enfin mesuré par l'oracle pass@$k$ stratifié par profondeur, dont le budget $k$ et la température sont rapportés à chaque usage (section 4.5) \textcolor{red}{je ne comprends pas de quoi on parle ici}. Chaque configuration correspond à un unique run d'entraînement : le budget mono-GPU interdit les répétitions multi-graines \textcolor{red}{rappelle moi ce qu'est multi graine : ablation ? ou juste seed ? Parce que si c'est seed on peut juste dire que c'est une contribution free}, et l'incertitude est quantifiée par le bruit d'évaluation décrit ci-dessus et par la stabilité temporelle des courbes \textcolor{red}{il faut détailler cette partie incertitude}.

\paragraph{Comparaison à compute contrôlé.} L'axe naturel de comparaison est l'époque, à savoir un passage complet sur les tâches d'entraînement. Deux précautions s'imposent. D'une part, le contenu d'une époque dépend de $G$ : à 64 trajectoires par mise à jour, une époque représente 46 mises à jour pour $G=8$ mais 93 pour $G=16$, soit deux fois plus de trajectoires ; les comparaisons entre tailles de groupe se font donc en mises à jour ou en trajectoires cumulées, jamais en époques. D'autre part, l'époque n'égalise pas le compute : c'est précisément ce que les curriculums modifient, un horizon réduit produisant des épisodes plus courts et moins coûteux. La section 4.4 complète donc chaque comparaison par trois indicateurs relevés dans les journaux d'entraînement : les tokens générés cumulés, qui dominent le coût réel, les heures GPU, et les tours d'environnement consommés. \textcolor{red}{très intéressant, à bosser un peu}

\paragraph{Nomenclature.} Les expériences sont désignées dans la suite par des noms courts ; le tableau~\ref{tab:nomenclature} donne la correspondance avec la numérotation interne du registre expérimental, conservée pour la traçabilité. Les expériences secondaires (évaluations few-shot, planification mono-tour, piste SNIS) sont nommées directement dans les sections concernées.

\begin{table}[htbp]
\centering
\caption{Nomenclature des expériences principales et correspondance avec la numérotation interne.}
\label{tab:nomenclature}
\begin{tabular}{lll}
\toprule
Désignation & Contenu & Numéros internes \\
\midrule
\textsc{Réplication-FT} & recette du papier en full finetuning & exp23 à 23.4 \\
\textsc{Grille-LoRA} & grille rang $\times$ LR $\times$ $\beta$ en LoRA & exp24 \\
\textsc{Ancre-Mobile} & recette stabilisée, baseline $G=8$ & exp25, 25.2 \\
\textsc{Ablation-Ancre} & carré ancre $\times$ $\beta$ & exp30, 31, 31.1 \\
\textsc{Horizon} & curriculum d'horizon d'interaction & exp32 \\
\textsc{Profondeur} & curriculum de profondeur de recettes & exp33 \\
\textsc{Magellan} & autocurriculum à progrès d'apprentissage appris & exp34 \\
\textsc{Budget} & curriculum de budget de génération par tour & exp35 \\
\textsc{Contrôle-G16} & baseline $G=16$ sans curriculum & exp36, 36.1 \\
\bottomrule
\end{tabular}
\end{table}



\subsection{Stabilisation et optimisation de l'apprentissage par renforcement en environnement multi-tours}



\textcolor{red}{Ajouter des figures pour illustrer les propos de chaque partie} 

Cette sous-section établit les conditions de stabilité de l'entraînement, sur lesquelles reposent toutes les comparaisons ultérieures. Elle suit la chronologie logique du problème : d'abord les conditions de possibilité du gradient lui-même, ensuite la réplication de la recette du papier en full finetuning, puis l'échec instructif de sa transposition en LoRA \textcolor{red}{comment ça l'échec on va plus loin que le papier}, enfin le correctif que nous proposons, l'ancre KL mobile, et les signaux avant-coureurs de collapse que l'ensemble de nos runs permet d'identifier. Chacune des trois pathologies annoncées en section 2.3, détournement de récompense, dérive de politique et collapse d'entropie, y trouve une illustration empirique et une contre-mesure. Pour les applications futures de commerce agentiques à Criteo, cette section pratique me semble la plus valorisable et transférable aux problématiques du laboratoire et de l'équipe Agentic Commerce. \textcolor{red}{l'organisation des découvertes est vraiment bien c'est top fable}

\subsubsection{Conditions de Possibilité : le Gradient et la Récompense}

\paragraph{Masquage des observations.} La première condition est passée inaperçue plusieurs semaines, car elle ne produit aucune erreur visible : le run s'exécute, la loss décroît, et la performance régresse sous la baseline. Notre implémentation initiale retournait à la bibliothèque d'entraînement les seuls tokens d'action de l'agent, sans les observations intercalées. Les log-probabilités du backward pass étaient donc recalculées dans un contexte différent de celui de la génération : les ratios d'importance étaient faux, et le gradient avec eux. \textcolor{red}{Alors intro trop story telling et pas besoin de dire qu'on a planché plusieurs semaines dessus. Dire juste la première difficulté c'est le gradient et le masking et que c'est pas évident d'y penser en multi tour. Mettre un schéma explicatif du masking, et vérifier que c'est bien juste les tokens d'agents backprop, et pas tous les tokens sans mask je suis pas certain.}Le correctif consiste à inclure les observations dans la séquence d'entraînement en les couvrant d'un masque \texttt{env\_mask} nul : le backward pass voit le contexte exact de la génération, et le gradient ne porte que sur les tokens produits par la politique. C'est l'analogue trajectoriel exact du masquage d'instruction du SFT (section 1.2), et une condition nécessaire de convergence : avant le correctif, quatre époques d'entraînement faisaient passer le modèle sous sa propre baseline \textcolor{red}{(17, en dessous de 18)} ; après, le même pipeline apprenait normalement. Nous insistons sur le caractère silencieux de ce défaut : aucune métrique standard ne le signale, seule la confrontation systématique au score du modèle non entraîné l'a révélé. \textcolor{red}{faudra dire qqp que se comparer exactement au seed et train set du papier était pour s'assurer de pouvoir comprendre et reproduire les ésultats avant d'ouvrir vers autre chose. SI on avait ouvert dès le début on aurait peut être pas focus ce détail technique. C'est pas le bon endroit pour en parler, section 3 plutot}

\paragraph{Reward shaping} La seconde condition concerne la conception de la récompense. Face à la sparsité du signal terminal, notre première campagne a testé une densification syntaxique : un bonus de $+0{,}02$ par tour à action unique et valide, une pénalité pour les actions invalides \textcolor{red}{retrouver en détail les shape de reward testés} . Le résultat illustre le détournement de récompense au sens strict de \cite{skalse2022defining}. La politique a maximisé la récompense façonnée en raccourcissant ses épisodes : les tâches triviales à une action étaient conservées, les tâches composées à plusieurs crafts étaient abandonnées \textcolor{red}{je comprends pas le lien de causalité entre une action par tour et cet abandon et perte de perf sur les tâches composées} , et le pass@$1$ est passé de 18\,\% à 14\,\%, avec un schéma caractéristique où les épisodes courts concentraient tout le reward positif. La correction du défaut d'implémentation du bonus n'a pas sauvé l'approche : la variante corrigée est descendue à 8\,\% \textcolor{red}{quelle correction} . Dans notre configuration, à vérificateur exact, le signal binaire outcome-only s'est révélé strictement préférable à ce schéma de shaping syntaxique. Nous nous gardons d'en faire une règle générale : densifier sans corrompre est précisément le problème que les modèles de récompense de processus cherchent à résoudre proprement (section 2.1), et notre observation dit seulement que la densification naïve est pire que la sparsité. \textcolor{red}{Définir par bon sens des bonus de rewards qui nous paraissent cohérent et potentiellement stabilisants pour l'apprentissage est souvent une mauvaise idée. Il est nécessaire pour cela d'avoir investiguer les mécanismes d'apprentissages contre intuitifs du RL qui ne bénéficieront pas de cette simplification apparente. Dr GRPO en parle d'ailleurs bien dans la section ... et ouvre un regard critique sur ce point. Dans beaucoup de papiers par ailleurs, les chiffres de reward shaping sont souvent des magic numbers dont les auteurs ne justifient pas la valeur. C'est notamment ce que nous avons vu à Criteo avec le papier SHOP-R1: REWARDING LLMS TO SIMULATE HUMAN BEHAVIOR IN ONLINE SHOPPING VIA REINFORCEMENT LEARNING} 
\textcolor{red}{On avait vu aussi toute une littérature qui parlait du reward : when it's sparse it is better to be dense, tech report de deepseek qui encourage au sparse, dr grpo pour mettre en garde contre des rewars d'apparence bons mais qui poussent à des comportements néfastes.} 
\textcolor{red}{Ce sujet mérite donc une ouverture scientifique et une analyse à peine plus poussée de nos résultats en rapport avec la littérature mentionnée. Nous avons choisi un reward sparse car c'est le plus souvent le cas dans les env et en maths bench et c'est ce qui est encouragé par les constructeurs (deepseek il me semble), mais la reco où les rewards sont pas forcément sparse nécessite un travail plus conséquent dessus et on a pas eu à le faire. Parler aussi des bails où c'est pas du reward shaping mais du reward door ou reward filter : 0 tant que ça passe pas tous les tests, puis sparse (je sais plus quel papier en parle, peut être self challenging agent de Jason Weston, checker la biblio)} 


\subsubsection{Full Finetuning : la Réplication comme Référence de Stabilité}

La lignée \textsc{Réplication-FT} porte la recette du papier, à hyperparamètres identiques (tableau~\ref{tab:hyperparams}) \textcolor{red}{lien ne marche pas} , sur notre pile logicielle mono-GPU. Elle établit deux faits. Le premier est que la recette tient : la performance croît de façon quasi monotone, de 40\,\% après 30 époques à \textbf{69\,\%} après environ 247 époques cumulées, à travers une série de reprises sur meilleur point de contrôle, nécessaire car notre environnement GPU Coder à Criteo connait des coupures hebdomadaires. À aucun moment cette lignée n'a présenté de collapse : la divergence KL vers l'ancre fixe reste faible de bout en bout, avec $\beta = 10^{-3}$ et sans aucun des mécanismes de stabilisation développés plus loin. Le second fait est le coût de cette stabilité : la progression est lente, chaque point de pass@$1$ au-delà de 60\,\% coûtant des dizaines d'heures de GPU, et l'état optimiseur du full finetuning sature la mémoire et le stockage. Cette lignée a aussi fourni la leçon de lecture statistique intégrée au protocole (section 4.1) : une reprise depuis un pic isolé à 64\,\% n'a jamais retrouvé ce niveau, les évaluations voisines oscillant entre 54 et 59\,\%, signature d'un maximum de bruit et non d'un niveau réel.
\textcolor{red}{Dire qu'on a analysé leur code et reproduit exactement à hyperparam égaux. En revanche il faut regarder le nombre d'epochs annoncé par le papier, et ce nb d'epoch est communiqué pour le 7B pas pour le 3B donc encore une possible enfumade de la part de nos amis asiatiques} 
\textcolor{red}{Attend mais on n'a pas un full ft qui arrive à plus de 70 ?} 
\textcolor{red}{Appuyer la stabilité de cette réplication : le reward sur train et test n'a jamais été aussi peu bruité (schéma)} 

\subsubsection{LoRA}

\textcolor{red}{LoRA without Regret, Accélération du training, libérer de l'espace pour augmenter G et affiner le gradient} 

L'adaptation LoRA promettait le même apprentissage pour une fraction du coût mémoire. \textcolor{red}{pas promettait, on s'est penché dessus parce que l'entrainement full ft était beaucoup trop long et ne nous aurait pas permi d'avoir le temps de tester toutes nos expe et le papier LoRA without Regret nous a en plus ajouté des pistes intéressantes d'explication via théorie de l'informations qui nous ont inspiré}. L'argument théorique est solide \textcolor{red}{Non pas solide c'est une idée même dans le papier ils le disent que ça peut être un élément d'explication}  : le RL à récompense terminale ne transmet que $\mathcal{O}(1)$ bits par épisode, de sorte qu'un sous-espace de paramètres restreint suffit à absorber l'information du signal, et la littérature \textcolor{red}{pas la littérature, les auteurs de LWR apportent des preuves expérimentales à leurs hypothèses qui nous ont inspirées} rapporte que LoRA égale le full finetuning en RL même à petit rang, avec un taux d'apprentissage recommandé dix fois supérieur \cite{schulman2025lora}. Cette recette, validée en mono-tour à récompense dense \textcolor{red}{dans leur cas (d'ailleurs elle est pas dense leur récompense si?)} , ne transfère pas \textcolor{red}{nécessairement} à notre régime. La grille \textsc{Grille-LoRA}, croisant rang ($r \in \{16, 32, 64\}$), taux d'apprentissage et $\beta$, à recette du papier par ailleurs, dessine un paysage sans zone de fonctionnement satisfaisante. Le taux d'apprentissage du blog provoque un collapse rapide et complet (pic à 24\,\% puis 0\,\%). À taux réduit, $\beta = 10^{-3}$ monte plus haut (31 à 35\,\% selon le rang) mais collapse systématiquement en seconde moitié de run ; $\beta = 10^{-2}$ ne collapse jamais mais plafonne (27 à 29\,\%), sous la référence full finetuning à budget égal. Le rang n'est pas la variable de sauvetage : $r=16$ fait mieux que $r=32$ et $r=64$, à rebours de l'intuition de capacité.

Ce paysage s'interprète géométriquement, par quatre mécanismes que nous présentons comme des hypothèses cohérentes avec nos observations mais partiellement confondues entre elles, les lignées différant aussi par leurs taux d'apprentissage et facteurs d'échelle \textcolor{red}{et chacune de nos hypothèses pourraient être être l'objet de recherches approfondies}. Premièrement, le full finetuning répartit un gain de reward en déplacements infinitésimaux de milliards de poids, solution proche en espace de fonctions et petite en KL ; le LoRA confine le même gain à un sous-espace de rang faible et doit y voyager plus loin, donc produit plus de KL par point de reward gagné. Deuxièmement, la chaleur effective du pas est amplifiée par le facteur $\alpha / r$. Troisièmement, la paramétrisation en produit de deux matrices couple multiplicativement leurs gradients. Quatrièmement, les gradients du reward sparse sont bruités par nature, et LoRA les concentre là où le full finetuning les moyenne. La conséquence commune de ces mécanismes est la même : à progression égale, la politique LoRA s'éloigne plus vite de sa référence KL, et le rappel vers une ancre fixe devient soit trop faible pour empêcher le collapse, soit assez fort pour l'empêcher mais au prix d'un plafond.
\textcolor{red}{ce deuxième paragraphe mérite d'être plus approfondi : un tiret par explication, et plus de matière pour chaque hypothèse. Backer ces explications par les courbes de collapse et de kl pour une ou plusieurs expé.} 

\subsubsection{L'Ancre KL Mobile : Correctif et Ablation}

\textcolor{red}{C'est notre solution pratique qui a mené à des résultats supérieure au papier d'origine mais nous ne prétendons que cette recette soit universelle ou meilleure objectivement. Des tests de robustesse sur d'autres environnements sont nécessaires et on va les faire grâce à l'env agentgym rl qui nous a fait le travail de créer ces env.} 

Ce diagnostic désigne le correctif. Si le problème est la distance croissante à une référence figée, il faut déplacer la référence \textcolor{red}{mal dit : si la référence est soit trop forte soit trop faible, il faut déplacer la référence. Aller chercher un ou plusieurs papiers qui parlent de moving anchor.} . Notre ancre KL mobile procède par \emph{merge-and-restart} périodique : toutes les quatre époques, l'adaptateur courant est fusionné dans le modèle de base, un adaptateur vierge est réinitialisé, et la pénalité KL se mesure désormais vers ce nouveau point. La politique est inchangée au moment du ré-ancrage ; seul le point de rappel bouge \textcolor{red}{point de rappel ? Et les optimizers sont réinitialisés aussi} . La mécanique du rappel change alors de nature : loin d'une ancre fixe, le gradient de $\beta \, D_{\text{KL}}$ tire la politique vers l'arrière, vers le modèle de base, ce qui dégrade le progrès sans amortir la dynamique locale ; une ancre récente agit au contraire comme une région de confiance locale, qui pénalise la dérive rapide par rapport au dernier point validé, précisément le type de rappel qui prévient un emballement. \textcolor{red}{Remarque : on retrouve dans cette solution des idées très proche des premiers algos de Deep RL à policy gradients : TRPO (trust region policy optimisation) qui travaillait sur des régions de confiance (notre ancre munie de son rappel fort) que ppo avait fait évolué en clipping. Noter qu'il faudra présenter rapidement en annexe TRPO} 

Les résultats valident le mécanisme. La configuration \textsc{Ancre-Mobile} ($r=8$, $\beta = 10^{-2}$, ré-ancrage toutes les 4 époques) monte de 12 à 35\,\% en 15 époques sans le moindre accident, là où la meilleure configuration à ancre fixe de la grille zigzaguait déjà ; poursuivie puis reprise depuis son meilleur point, elle atteint \textbf{65\,\%} après une centaine d'époques et 27 ré-ancrages, puis \textbf{73\,\%} en fin de lignée, sans un seul collapse. L'ablation \textsc{Ablation-Ancre} complète le carré ancre $\times$ $\beta$ (tableau~\ref{tab:anchor_beta}) et permet deux conclusions. D'une part, l'ancre mobile est \emph{nécessaire} : aucune valeur de $\beta$ ne sauve l'ancre fixe, qui plafonne puis décroche à $\beta = 10^{-2}$, collapse à $10^{-3}$, et diverge violemment à $10^{-4}$, la KL atteignant des valeurs numériquement absurdes. D'autre part, l'ablation résout un facteur confondu de la grille initiale : $\beta = 10^{-3}$, qui collapsait systématiquement à ancre fixe, est stable à ancre mobile (53\,\%). C'est donc bien le régime d'ancre, et non la seule valeur de $\beta$, qui gouverne la stabilité ; $\beta = 10^{-2}$ reste néanmoins préférable pour le niveau atteint \textcolor{red}{On a vraiment ce qui faut en terme d'expériences pour claim un truc aussi fort ?} . En résumé : le full finetuning amortit naturellement le RL à reward sparse ; LoRA concentre le même progrès dans un sous-espace étroit à pas amplifié, s'éloigne plus vite en KL, et exige un rappel adapté ; déplacer périodiquement l'ancre transforme ce rappel en région de confiance locale et lève le plafond sans sacrifier la stabilité.

\textcolor{red}{parler de l'avantage que permet LoRA en plus d'avoir des supers scores et accélérer le training : libérer de l'espace pour augmenter G et préciser le signal !} 

\begin{table}[htbp]
\centering
\caption{Carré d'ablation ancre $\times$ $\beta$ (LoRA $r \in \{8, 16\}$, recette stabilisée par ailleurs, pass@$1$ sur les 100 tâches de test). Chaque cellule résume un run unique ; « stable » signifie l'absence de tout épisode de décrochage sur la durée du run.}
\label{tab:anchor_beta}
\begin{tabular}{lccc}
\toprule
 & $\beta = 10^{-2}$ & $\beta = 10^{-3}$ & $\beta = 10^{-4}$ \\
\midrule
Ancre fixe & pic 39 puis décrochage & pic 35 puis collapse & pic 34 puis divergence KL \\
Ancre mobile (4 ép.) & \textbf{65 puis 73, stable} & 53, stable & non couru \\
\bottomrule
\end{tabular}
\end{table}

\begin{figure}[htbp]
    \centering
    \framebox[\textwidth]{\vbox to 5cm{\vfil\centerline{\textbf{[FIGURE : courbes pass@1 test, stabilisation]}}\vfil}}
    \caption{Courbe maîtresse de la stabilisation (Qwen2.5-3B, LoRA, pass@$1$ test au fil des mises à jour) : ancre fixe $\beta=10^{-3}$ (pic puis collapse), ancre fixe $\beta=10^{-2}$ (plafond puis décrochage), ancre mobile $\beta=10^{-2}$ (montée continue jusqu'à 65--73). Runs uniques ; bruit d'évaluation $\pm 4$--5 points.}
    \label{fig:stabilization}
\end{figure}

\subsubsection{Signaux Avant-coureurs et Synthèse}

L'ensemble de nos runs, huit collapses documentés et cinq runs longs stables, permet une analyse des signaux avant-coureurs. Nous la présentons comme diagnostique et non causale : les hyperparamètres varient entre lignées, mais les signatures, elles, sont remarquablement constantes. Le premier signal est l'\textbf{entropie de la politique} : dans tous nos collapses, elle décroche sous environ $0{,}15$ avant l'effondrement du score, la politique devenant quasi déterministe et l'exploration s'éteignant, conformément à la loi empirique de \cite{cui2025entropy} selon laquelle la performance se paie en entropie \textcolor{red}{je sais pas de quoi on parle ici} . Notre cas le plus instructif est le contrôle \textsc{Contrôle-G16} (détaillé en section 4.3) : le passage à $G=16$ \textcolor{red}{on balance le g16 comme ça sans intro ! si c le cas il faut ajouter un paragraphe au dessus dans la sous section LoRA pour expliquer son intérêt dans le calcul du but} accélère l'apprentissage d'un facteur proche de deux, et avec lui le rythme de la déterminisation ; l'entropie passe de $1{,}0$ à $0{,}10$ en quatre époques, puis la KL s'emballe, alors même que le reward d'entraînement se maintient \textcolor{red}{il nous faut des graphs pour appuyer tout ça parce que là c'est le bordel} . Le deuxième signal est la \textbf{longueur des générations}, qui s'écrase (sorties stéréotypées de moins de 250 tokens) pendant la déterminisation, puis explose (plusieurs milliers de tokens de texte dégénéré) après l'emballement. Le troisième est la \textbf{divergence KL} elle-même, dont le franchissement de l'unité a toujours été, chez nous, irréversible. Notons enfin que l'ancre mobile ne protège pas de ce mode d'échec : en ré-ancrant sur une politique en cours de déterminisation, elle fige la dégénérescence au lieu de la freiner. La stabilité exige donc à la fois un rappel adapté et une source d'entropie, rôle que la section 4.3 attribuera au curriculum. \textcolor{red}{Il faudrait analyser les courbes d'erreurs aussi pour voir si on a des corrélations entre les erreurs et le collapse. D'ailleurs une étude de corrélation entre les signaux cités ci dessus et le collapse qui suit serait bienvenue (matrice de corrélation)} 

La régulation du taux d'apprentissage mérite une mention, car elle a mobilisé une lignée entière d'essais \textcolor{red}{trop familier dans le ton}. Les paliers calendaires stabilisent mais coupent des dynamiques de progression en cours ; la coupe adaptative déclenchée sur stagnation du reward se déclenche trop souvent sur du bruit, et surtout, couper consolide l'état courant, qui après un pic est souvent une dérive ; la variante la plus saine combine paliers longs et restauration du meilleur état du palier écoulé. Ces mécanismes, ainsi que la politique de persistance des meilleurs points de contrôle avec leur état optimiseur, sont documentés en annexe C. La recette stabilisée finale (tableau~\ref{tab:hyperparams}) n'en a finalement plus besoin : à taux constant, l'ancre mobile suffit. \textcolor{red}{ah oui là tu parles du lr décroissant ? Il faudrait peut être un graph avec les expé concernées, mais si il y a pas assez de matière tant pis ta dernière phrase balaie bien le sujet} 

Pour terminer cette partie, résumons les princpaux findings. Masquer les observations et vérifier tout pipeline contre le score du modèle non entraîné, car le défaut est silencieux. Préférer la récompense binaire exacte à toute densification syntaxique naïve \textcolor{red}{j'aurai plutot dit bien investiguer le fonctionnement du rl et la propagation du signal reward dans les mécanismes d'apprentissage avant de faire du reward shaping} . Avec LoRA, ne pas importer les recettes du mono-tour dense \textcolor{red}{ça ne veut rien dire on parle d'un papier c'est pas des recettes uniberselles} ; réduire le taux d'apprentissage, monter $\beta$, et déplacer l'ancre périodiquement. Surveiller l'entropie comme signal principal, avec la longueur des générations comme confirmation, et considérer tout franchissement durable de KL $> 1$ comme un point de non-retour. Enfin, ne lire les scores qu'au travers du bruit d'évaluation, et ne jamais faire confiance à un pic isolé. 
\textcolor{red}{J'aime pas ce dernier paragraphe : trop familier, on est pas un compte insta qui donne des tips. Reformuler en conclusion de partie qui rappelle les résultats, oubli cette idée de guide ou tuto je sais pas quoi.} 






\subsection{Expériences RL : Curriculums et Transfert}
\label{sec:curriculums}

La section~4.2 a montréque l'ancre KL mobile fournie un rappel adapté pour stabiliser GRPO multi-tour. Une source d'entropie est également nécessaire \textcolor{red}{une source d'entropie ? qu'est ce que tu racontes} . La présente section examine si le curriculum peut tenir ce second rôle \textcolor{red}{il faut dire que cette deuxième partie expérimentale peut se lire de deux façons : la suite de nos travaux de stabilisation et l'évaluation des curriculums dans ce role, et aussi une étude du curriculum dans certaines questions plus recherche et meta learning type on peut pbattre la passk et on peut faire du transfert learning et est ce qeu ça améliore, au delà du simple stabilisant. C'est pour cela que les curriculums sont un sujet intéressant à la frontière entre optimisation rationnelle utile et intérêt marqué pour la curiosité en rl sur l'exploration ! tu le dis déjà bien dans ce paragraphe mais je pense qu'il faut que cette dualité en soit l'idée centrale. En gros trop intéressant mais risquée sur meta learning, et au pire ça stabilise}, tout en instruisant la question centrale du stage : ordonner l'expérience présentée à la politique accélère-t-il l'apprentissage, en élève-t-il l'asymptote, et le gain acquis sur une difficulté se transfère-t-il aux autres \textcolor{red}{et aussi permet il d'identifier des comportements intéressants généralisables dans par exemple la vitesse d'apprentissage, la généralisation, la limitations des erreurs, etc} ? Nous décrivons d'abord quatre régimes de curriculum comparés à recette commune, puis leurs résultats et la lecture qu'ils autorisent, avant d'aborder le transfert entre planification mono-tour et contrôle multi-tour \textcolor{red}{on n'a pas eu le temps de faire cette dernière expé il faudra la mettre dans ouverture, en précisant que c'est une forme de meta learning et ouvrir sur le meta learning en général au delà des curriculums, avec le papier de lecun et dupoux par exemple} .

\subsubsection{Quatre régimes de curriculum à recette commune}

\textcolor{red}{faire des bullet points par type de curriculum arpès, une par curriculum} 
Un curriculum désigne ici une politique qui module, au fil des époques, la difficulté effective des épisodes d'entraînement, par contrainte de budget ou par échantillonnage des tâches \textcolor{red}{essayer de trouver une définition dans un ou plusieurs papiers, puis mettre en lumière que c'est pas juste ordonner les tâches ça peut être aussi limiter les capacités du modèle.} . Les quatre régimes partagent la recette stabilisée en 4.2 : LoRA de rang 8 ($\alpha=32$), taux d'apprentissage $3\times10^{-6}$, $\beta=0{,}01$, ancre mobile toutes les 4 époques, groupes de taille $G=16$ et 64 trajectoires par mise à jour. Ils se comparent à deux références sans curriculum : \textsc{Ancre-Mobile} ($G=8$, 65 puis 73~\%) et \textsc{Contrôle-G16}, qui reprend la recette commune à l'identique et isole ainsi l'effet de la taille de groupe \textcolor{red}{faut dire que ça n'a pas tenu, témoin de la stabilité du curriculum !}. Conformément au protocole de la section 4.1, chaque configuration correspond à un run unique et aucun hyperparamètre n'est réglé par bras \textcolor{red}{c'est quoi un bras ? } : les conclusions sont conditionnelles à cette recette commune.

Les deux premiers régimes agissent sur le budget alloué à l'agent. \textsc{Horizon} reprend le principe de ScalingInter, proposé par le papier de référence pour étendre progressivement l'espace d'interaction \cite{xi_agentgym-rl_2025} : le plafond de tours vaut 10 dès l'époque 0, 20 dès la 15, 30 dès la 30. Un plafond bas force des solutions courtes et rend les tâches profondes temporairement insolubles ; chaque relèvement les ré-introduit en cours d'entraînement.\textcolor{red}{voilà ça typiquement ici il faut faire un bullet point par type de curriculum. Pour le précédent spécifiquement, être plus clair en disant nombre de tour, et en interprétant et en donnant l'hypothèse qui est que restreindre le nombre de tour pousse le modèle à réussir les tâches résolvables en peu de tour, donc soit d'abord les plus faciles soit les moyennes mais d'une manière très efficiente, et certainement les deux en meme temps, et on analysera donc les courbes de solved par depth pour comparer les différents régimes}  \textsc{Budget} transpose l'idée au budget de génération par tour : 256 tokens dès l'époque 0, 512 dès la 15, 1\,024 dès la 35, l'horizon restant fixé à 30 tours. \textcolor{red}{un peu la même idée mais sur le nombre de tokens générés au lieu du nb de tours. Penser à écrire en conclusion de cette sous partie et en ouverture qu'on pourait totalement imaginer de combiner les curriculums pour voir si certaines combinaisons permettent d'encore meilleurs résultats} 

\textsc{Profondeur} ordonne au contraire la difficulté par la structure des tâches sans toucher aux capacités du modele \textcolor{red}{définition la plus proche du curriculum} . L'échantillonnage est uniforme sur les tâches de profondeur inférieure ou égale au palier courant : $\leq 1$ à l'époque 0, $\leq 2$ à la 12, $\leq 3$ à la 25, $\leq 4$ à la 37. Le mécanisme est un échantillonneur pondéré opérant sur le dataset entier, et non un découpage en sous-datasets : la définition des époques et le reste de la boucle d'entraînement restent identiques aux autres bras \textcolor{red}{je sais toujours pas ce qu'est un bras, et pour cette expé on a changé le fonctionnement en mettant la possibilité de passer au pallier suivant si train reward stagne au dessus de 80 pendant une epoch je crois, à vérifier le pourcentage exact } .


\textcolor{red}{Expliquer l'utilité de l'autocurriculum : on n'a la plupart du temps pas accès aux difficultés intrinsèques des tâches, et meme si on en a une idée humaine rien ne dit que la notion de difficulté pour le modèle soit totalement alignée sur la notre. Un sélecteur automatiqeu de tâche est essentiel, surtout dans des tâches de reco ou la difficulté peut être dure à déterminer (nombre d'appel à outils, nb de tour réalisé pour finir la tache par un llm sota propriétaire (comme toolathlon utilise opus)} 
Le quatrième régime, \textsc{Magellan}, remplace les paliers fixés à la main par un autocurriculum. La section 2.4 a présenté le progrès d'apprentissage (\emph{learning progress}) comme critère de sélection des tâches, introduit par les travaux sur le curriculum automatique de Graves et al., où un bandit alloue l'entraînement aux tâches dont la performance évolue le plus \cite{graves2017automated} ; des garanties théoriques récentes soutiennent ce principe pour le raisonnement \cite{rajaraman2026curriculum}. MAGELLAN rend ce critère praticable à l'échelle d'un espace de tâches en apprenant à \emph{prédire} le progrès à partir des représentations internes du LLM, ce qui généralise l'estimation aux tâches peu visitées \cite{gaven2025magellan} \textcolor{red}{pour le moment je n'ai aucune idée de ce que raconte les autres papiers que magellan ci dessus, ni véritablement de ce que raconte magellan précisément et de notre implémentation ilf faut que je regarde en détail} . Notre port en suit fidèlement le mécanisme : l'échantillonnage des tâches est proportionnel au progrès d'apprentissage prédit ; l'estimateur est une tête MLP posée sur l'état caché du dernier token de l'énoncé, portée par des adaptateurs dédiés séparés de la politique, et entraînée en continu sur les couples (tâche, succès) observés ; le progrès d'une tâche est l'écart absolu entre la prédiction courante et celle d'une copie retardée de l'estimateur, distante d'environ 100 mises à jour ; un mélange $\varepsilon$-greedy, décru de 1{,}0 à 0{,}2, garantit une exploration résiduelle de tout le support \textcolor{red}{la typiquement moi qui suit pas hyper familier du papier je comoprends rien donc mes encadrants encore moins. Il faudrait faire une présentation un peu plus claire du papier et introduire les notions dont tu parles, ou alros se contenter d'un paragrpahe court qui renvoie à une annexe pour présenter plus honorablement magellan qui est un vrai travail de recherche à part entière bien plus conséquent que juste avoir l'idée de scale le nb de tokens par gen} . Ce régime porte aussi un enjeu propre : la profondeur est une difficulté apparente pour l'humain, et rien ne garantit qu'elle recoupe la structure de compétence du modèle. La partition des tâches induite par l'estimateur appris pourra être comparée à la partition par profondeur ; une divergence indiquerait que le curriculum humain optimise le mauvais axe. Une prédiction a été pré-enregistrée avant le lancement : l'autocurriculum doit désinvestir la profondeur 4, dont le support d'entraînement est quasi nul (un seul item, cf.~section 3.3).\textcolor{red}{je comprends pas cette dernière phrase, qui a fait quelle prédiction qu'est ce que tu racontes} 

\subsubsection{Résultats : le curriculum comme régularisateur d'exploration}

\textsc{Horizon} est le seul régime terminé au 28 août \textcolor{red}{on est le 2 sept mets à jour les expé qui ont tourné} , et il établit le record du projet. Le meilleur score, 82~\%, est atteint deux fois au cours du run, ce qui écarte l'hypothèse d'un pic isolé ; la montée passe par environ 67~\% vers l'époque 29 et 82~\% vers l'époque 66, avant un plateau final entre 75 et 82 qui a motivé un arrêt volontaire vers l'époque 80, \textcolor{red}{rappeler ce chiffre en lumière des 200 et qulques du full fine tuning rl } . Le run traverse 20 ré-ancrages sans collapse, avec une entropie stable entre 0{,}6 et 0{,}7. Ce score dépasse les 75~\% du modèle AgentGym-RL-3B du papier (tableau~\ref{tab:agentgym_textcraft}) et les 73~\% de la référence \textsc{Ancre-Mobile} à $G=8$, au bruit d'évaluation près (section 4.1) ! \textcolor{red}{encore une fois ce bruit d'évaluation c'est un truc que tu as sorti tout seul je n'en ai jamais parlé, je sais pas si c'est utile, je sais pas si agentgym en parlent, est ce que ton estimation est rigoureuse, basée sur quoi ? } 

\textcolor{red}{Peut être commencer cette sousousection résultats par le run contrôle} Le contrôle éclaire l'origine de ce gain. Le premier run de \textsc{Contrôle-G16} a collapsé : montée rapide de 12 à 54~\%, plus vite que la référence $G=8$ à trajectoires égales (récompense d'entraînement 0{,}45 contre 0{,}30 au même budget), puis effondrement de l'entropie (médiane passée de 1{,}0 à 0{,}59 puis 0{,}10 en quatre époques), emballement de la KL, évaluation retombée de 54 à 13 et générations dégénérées \textcolor{red}{faut éviter ce genre d'enchainement d'idées par virgules avec des acronymes, ça vaut pour tout le rapport dès que tu détectes ce genre de formulation faut changer et refaire des phrases normales verbeuses, en rappelant de temps en temps ce que sont les acronymes} . Ce profil reproduit le mécanisme de collapse d'entropie identifié en RLVR, où la politique se déterminise plus vite que le rappel ne la freine \cite{cui2025entropy}, ici malgré l'ancre mobile (section 2.3). Une relance avec ré-ancrage toutes les 12 époques est en file, \textcolor{red}{néanmoins nous n'avons pas eu le temps de nous pencher davantage là dessus étant donné que ce collapse est en soi pour nous un résultat à comparer avec nos curriculums, et ce collapse illustre parfaitement les capacités de stabilisation des curriculums.} .

Ces deux runs livrent la lecture centrale de la section. À recette commune, $G=16$ accélère l'apprentissage et, dans le même mouvement, la déterminisation de la politique ; sans curriculum, la recette commune est instable à cette vitesse. Le curriculum d'horizon agit alors comme un régularisateur d'exploration : chaque relèvement du plafond ré-injecte de la difficulté, maintient des groupes GRPO mixtes, donc des avantages non nuls, et soutient l'entropie. Il n'est pas qu'un accélérateur. Ce constat reste conditionnel : un run par bras \textcolor{red}{arrête avec ce mot} , aucun réglage par bras, et le contrôle propre de l'effet de $G$ attend la relance en file. La question de savoir si le gain de \textsc{Horizon} étend le support du modèle ou concentre un support préexistant est instruite en section 4.5 avec l'outillage pass@$k$. \textcolor{red}{pas clair cette fin de paragraphe} 

\textsc{Profondeur}, \textcolor{red}{mets à jour avec nos résultats} en cours au 28 août, fournit déjà une observation attendue par la section 2.4. Sur le palier $\leq 1$, la récompense d'entraînement sature à 1{,}0 des époques 6 à 12 environ : tous les groupes deviennent uniformes, l'avantage est nul, le GRPO ne produit plus de gradient, et l'entropie descend vers 0{,}19. C'est la démonstration empirique de la limite des paliers statiques maintenus trop longtemps. Au passage au palier $\leq 2$ à l'époque 12, la récompense retombe à 0{,}25 et l'entropie remonte vers 0{,}33 : la difficulté ré-injectée rétablit le signal. L'évaluation se situe entre 28 et 29~\% vers l'époque 17. \textsc{Magellan} \textcolor{red}{mets à jour, dis que on espère avoir le temps de run magellan} et \textsc{Budget} sont en file au 28 août ; la comparaison complète des régimes, à compute contrôlé, fait l'objet de la section 4.4.

Deux analyses de transfert, extraites des logs existants, seront produites au gel expérimental ; le texte y renvoie sans en anticiper les résultats. \textcolor{red}{lesquelles sont elles ces analyses de transfert exactement ? } 

\begin{figure}[htbp]
  \centering
  \framebox[0.9\textwidth]{\rule{0pt}{5cm} figure à produire au gel}
  \caption{Pass@1 par profondeur au fil de l'entraînement, pour chaque régime de curriculum et les deux références sans curriculum. L'ordre d'acquisition des profondeurs indique si le GRPO uniforme suit un curriculum émergent (profondeur 1, puis 2, puis 3) et si les curriculums explicites modifient cet ordre.}
  \label{fig:acquisition_depths}
\end{figure}

\begin{figure}[htbp]
  \centering
  \framebox[0.9\textwidth]{\rule{0pt}{5cm} figure à produire au gel}
  \caption{Taxonomie des erreurs au fil des évaluations périodiques, par profondeur et par régime (aires empilées). L'ordre d'élimination des classes d'erreurs (format, recette erronée, planification) indique ce que chaque régime corrige en premier.}
  \label{fig:elimination_erreurs}
\end{figure}


\subsection{Comparaison à Compute Contrôlé}
\label{subsec:compute-controle}

La section~\ref{subsec:protocole} a fixé la convention de comptage : chaque mise à jour de gradient (màj) consomme 64 trajectoires, quelle que soit la taille de groupe $G$. Une époque vaut donc 46~màj à $G=8$ et 93~màj à $G=16$. Comparer deux runs en époques n'a de sens qu'à $G$ égal ; entre $G$ différents, la comparaison se fait en màj \textcolor{red}{je veux pas ce genre d'acronymes} ou en trajectoires cumulées. Même à $G$ constant, l'époque n'égalise pas le compute : les curriculums modifient précisément la longueur des épisodes, donc le volume généré par màj \textcolor{red}{bah oui mais c'est l'avantage des curriculums ça il faut le prendre en compte sinon tu montre pas l'intérêt des curriculums dans ta comparaison} . Nous retenons les trois indicateurs posés en~\ref{subsec:protocole} : tokens générés cumulés (ils dominent le coût), heures GPU \textcolor{red}{c'est vraiment pertinent heures gpu ?} et tours d'environnement consommés \textcolor{red}{comment tu justifies que c'est une mesure intéressante ?} . Les chiffres sont arrêtés au 28/08 ; les runs encore en cours sont signalés \textcolor{red}{propose moi déjà des résultats et premières interprétations} .



Avant de comparer les runs, il faut isoler l'effet de $G$ sur la vitesse brute. Un banc dédié (8~màj réelles par configuration, 64~trajectoires par màj) donne 57{,}1~s/màj à $G=8$ avec le réglage mémoire historique, contre 46{,}9~s/màj à $G=16$ avec cache élargi, soit environ 18~\% de gain. Le mécanisme est le partage du préfixe : les 16 \emph{rollouts} d'une même tâche réutilisent le cache KV de son prompt. Élargir davantage la mémoire du cache est contre-productif (61{,}5~s/màj). Ce banc mesure une vitesse à conditions fixées ; la vitesse moyenne d'un run réel dépend en outre de la longueur des épisodes, donc de la politique et du régime de curriculum. Le tableau~\ref{tab:compute} récapitule le coût des runs principaux.


\textcolor{red}{le tableau suivant dépasse à droite} 
\begin{table}[htbp]
\centering
\caption{Coût en compute des runs principaux (chiffres arrêtés au 28/08 ; 64~trajectoires par màj pour tous les runs).}
\label{tab:compute}
\begin{tabular}{lcccccc}
\toprule
Run & $G$ & màj & Heures GPU & s/màj moyen & Tokens générés ($\times 10^8$) & Pass@1 best \\
\midrule
\textsc{Ancre-Mobile} (initial) & 8 & $\sim$5\,100 & \textcolor{red}{[à extraire au gel]} & \textcolor{red}{[à extraire au gel]} & 4{,}51\textsuperscript{a} & 65 \\
\textsc{Ancre-Mobile} (reprise) & 8 & 2\,726 & $\sim$20 & $\sim$26 & 1{,}68 & 73 \\
\textsc{Horizon} & 16 & $\sim$7\,470 & $\sim$48 & $\sim$23 & 4{,}53 & \textbf{82} \\
\textsc{Contrôle-G16} & 16 & $\sim$1\,420 & 15{,}7 & $\sim$40 & 1{,}55 & 54\textsuperscript{b} \\
\textsc{Profondeur} & 16 & en cours & en cours & $\sim$41\textsuperscript{c} & 0{,}92\textsuperscript{d} & en cours \\
\bottomrule
\end{tabular}

\smallskip
{\footnotesize \textsuperscript{a}~Lignée complète jusqu'à l'époque $\sim$111, dernière valeur journalisée avant l'interruption. \textsuperscript{b}~Best avant collapse. \textsuperscript{c}~Valeur typique en début de run (horizon constant de 30~tours). \textsuperscript{d}~Cumul au 28/08, run en cours.}
\end{table}

La première lecture porte sur le curriculum d'horizon. À recette et $G$ identiques, \textsc{Horizon} domine \textsc{Contrôle-G16} sur les deux axes : le score final (82 contre 54, ce dernier suivi d'un collapse) et le coût par màj ($\sim$23 contre $\sim$40~s/màj). L'explication est mécanique : le palier initial plafonne les épisodes à 10~tours, là où le contrôle joue 30~tours dès le départ. La vitesse de \textsc{Profondeur} ($\sim$41~s/màj à horizon constant) recoupe cette lecture. Comme les époques à $G=16$ comptent le même nombre de màj, le curriculum d'horizon est aussi moins cher à époques égales. Rapporté à la màj, il a généré près de moitié moins de tokens que le contrôle, au caveat près que le collapse de ce dernier allonge ses générations finales. \textcolor{red}{les mots caveats, bras et maj tu arretes je veux plus les entendre} La question de savoir si le curriculum accélère seulement l'apprentissage ou en relève l'asymptote reste ouverte jusqu'aux runs restants. La figure~\ref{fig:pass1-tokens} met ces trajectoires sur l'axe qui égalise le compute.

\begin{figure}[htbp]
\centering
\framebox[0.9\textwidth]{\rule{0pt}{5cm} Figure à produire au gel}
\caption{Pass@1 périodique en fonction des tokens générés cumulés, pour \textsc{Horizon} ($G=16$), \textsc{Contrôle-G16} ($G=16$), \textsc{Ancre-Mobile} ($G=8$, reprise incluse) et \textsc{Profondeur} ($G=16$, en cours). Chaque point est une évaluation sur les 100 items de test. À abscisse égale, les runs ont consommé le même compute de génération ; l'écart vertical mesure l'efficacité en compute de chaque régime.}
\label{fig:pass1-tokens}
\end{figure}

La seconde lecture concerne l'effet de $G$ seul, à trajectoires cumulées égales. Un unique couple pré-collapse est disponible : \textsc{Ancre-Mobile} ($G=8$) contre \textsc{Contrôle-G16} ($G=16$), à recette par ailleurs identique. Au même budget de trajectoires, le second atteint un reward d'entraînement de 0{,}45 contre 0{,}30 pour le premier. Cette montée plus rapide est cohérente avec une meilleure estimation de l'avantage à $G=16$, mais elle repose sur un seul couple de runs : c'est une observation, pas une loi, \textcolor{red}{et ouvre vers des pistes d'études} . Deux limites bornent enfin toute cette comparaison : chaque configuration n'a été courue qu'une fois, et le bruit d'évaluation ($\pm$4-5~pts) rend les écarts de moins de 10~points peu concluants \textcolor{red}{ah oui mais là typiquement ce genre d'argument ça décrédibilise tout notre taf et le fait qu'on ait réussi à battre le paper. Ils parletn de bruit d'eval dans agentgym rl ou pas, j'ai pas souvenir perso} . \textsc{Profondeur} et les contrôles restants complèteront le tableau au gel des expériences \textcolor{red}{update avec l'état actuel} .









\subsection{Expériences d'Inférence et Ablations Transverses}
\label{sec:inference_ablations}

\textcolor{red}{j'ai ajouté la première sousousection, change et introduit là ici} 
Les sections précédentes ont porté sur l'entraînement. Celle-ci rassemble les mesures obtenues sans entraînement, ou en marge de celui-ci \textcolor{red}{ou qu'on n'a pas eu le temps de refaire dans le dernier setup de stabilisation ou qu'on a pas pu mener plus en détail parce que ça sortait du cadre central du projet} , qui donnent leur cadre aux résultats RL : l'oracle pass@$k$ qui met à l'épreuve l'hypothèse du plafond posée en section 2.3, le prompt few-shot comme injection de support exogène, une ablation entre générations de modèles, et le bilan de la piste SNIS \textcolor{red}{et donc le transfert learning} . Le tableau~\ref{tab:recap_section4} récapitule l'ensemble des configurations de la section 4.






\subsubsection{Oracle pass@$k$ et mise à l'épreuve de l'hypothèse du plafond}

Le pass@$k$ mesure la probabilité qu'au moins un tirage parmi $k$ résolve la tâche. Notre oracle tire, pour chacun des 100 items du test, $N=20$ trajectoires indépendantes à température 1{,}0 et 30 tours au plus, puis calcule le pass@$k$ par l'estimateur non biaisé usuel, globalement et par profondeur. La thèse de Yue et al.~\cite{yue2025rlvr} \textcolor{red}{faudrait lire ce papier et vérifier que c'est bien ce qu'il raconte}, posée en section 2.3 comme hypothèse à tester, fait du pass@$k$ du modèle nu un plafond pour le RLVR. L'oracle en fournit la mesure de départ.

Sur Qwen2.5-3B nu, le pass@1 vaut 10{,}3~\% \textcolor{red}{alors faut savoir parce que desfois dans le rapport on dit 18 desfois 10 et agentgym rl disent 14} et le pass@20 vaut 47~\%. La décomposition par profondeur sépare deux régimes. À $k=20$, la profondeur 1 atteint 94~\% et la profondeur 2 atteint 44~\%, contre un pass@1 de 4~\% à ce niveau : la profondeur 2 est un mur de \emph{fiabilité}, avec une large marge récupérable. Les profondeurs 3 et 4 restent à 0~\% (0 succès sur 500 tirages en profondeur 3, 0 sur 60 en profondeur 4) : c'est un mur de \emph{capacité}. La conséquence pour GRPO est directe : sans aucun succès, les groupes sont uniformément en échec, l'avantage relatif est nul et le gradient aussi~\cite{liu2025drgrpo}. Le passage au 7B nu déplace les niveaux sans lever le mur : pass@1 33{,}2~\%, pass@20 68~\%, et par profondeur à $k=20$ : 94, 76 et 32~\%, puis 0 sur les 3 items de profondeur 4. \textcolor{red}{si on en parle pas ailleurs faudra dire aussi que les modèles plus gros performent mieux et que nous on a choisit le 3B pour des raisons de compute initialement et aussi pour criteo ça se justifie aussi par les contraintes de rapidité d'inférence pour que le user ne doivent pas attendre trop longtemps sa reco}

Un modèle RL intermédiaire de la lignée full-FT \textcolor{red}{motiver ce choix parce que là ça sort de nulle part, dire par ex meme sans prendre le modele le plus performant on voit que l'hypothèse est dépassée. C'est d'ailleurs pour ça qu'il faut vraiment vérifier le papier qui annonce ce passk oracle parce que c vraiment une hypothèse qui marche pas je sais pas comment ils ont backé ça}(pass@1 de 32~\%) apporte un premier indice contre une lecture stricte du plafond. Son pass@20 atteint 65~\%, et sa frontière d'atteignabilité s'est étendue : 14 items de profondeur 2 et 2 items de profondeur 3, jamais résolus par le modèle nu en 20 tirages, le sont désormais au moins une fois. Le RL n'a donc pas seulement fiabilisé des succès existants, il a élargi le support mesuré. \textcolor{red}{enumeration pas assez verbeuse, il faut pas aller aussi vite, mes lecteurs ne sont pas familiers de tout ça.}

Le test le plus direct vient de \textsc{Horizon}. Ce run atteint un pass@1 de 82~\% (section 4.3, $G=16$), alors que le modèle nu plafonne à un pass@20 de 47~\% : le pass@1 après RL dépasse le pass@20 avant RL. Nous présentons ce résultat comme une mise en tension de l'hypothèse du plafond, pas comme une réfutation. Plusieurs réserves s'imposent. Le budget $k=20$ est fini et le pass@$k$ croît avec $k$ ; Yue et al. mesurent leurs croisements à des budgets bien supérieurs. La mesure dépend aussi de la température de l'oracle (1{,}0 ici). L'analyse par profondeur reste à produire au gel des expériences, en croisant les journaux d'évaluation de \textsc{Horizon} avec ceux de l'oracle : sur quels items le modèle entraîné réussit-il là où le modèle nu échoue 20 fois ? La profondeur 4 reste par ailleurs à 0, et le confound du jeu d'entraînement (un seul item de profondeur 4, section 3.3) borne ce que TextCraft peut prouver sur ce mur. Si le calendrier le permet, un complément tranchera entre extension et concentration du support au sens du \emph{sharpening}~\cite{huang2024self} : mesurer le pass@20 du meilleur modèle \textsc{Horizon}.

\subsubsection{Few-shot : gain immédiat, format et amorçage}

Le prompt few-shot consiste à fournir au modèle $k$ exemples résolus avant la tâche courante, sans mise à jour des poids~\cite{brown2020gpt3}. Sur Qwen2.5-3B, $k=5$ exemples insérés en historique de dialogue portent le pass@1 de 18 à 31~\%. Le format conditionne le signe du gain : les mêmes exemples concaténés en bloc de texte dans le prompt système donnent 5 à 8~\%, sous le zéro-shot. Le modèle imite alors la forme du bloc et hallucine les observations de l'environnement au lieu d'attendre les siennes.

Au-delà du gain immédiat, le few-shot joue un rôle d'amorçage pour le RL. En zéro-shot, le pass@20 vaut 0~\% en profondeur 3 : aucun groupe GRPO mixte n'est possible à ce niveau. Avec $k=20$ exemples, le pass@10 y atteint 16~\%. Les budgets sont inégaux (10 tirages contre 20), la comparaison est donc conservatrice. Les trois valeurs de $k$ mobilisées répondent à trois usages : 5 est le meilleur point du balayage pour le gain direct, 20 caractérise le régime d'amorçage, 10 est la valeur utilisée à l'entraînement.

L'interaction few-shot~$\times$~RL ne dispose que d'une paire propre : deux runs identiques à paliers longs, dont seul le prompt diffère, donnent 41~\% en zéro-shot contre 49~\% avec $k=10$ (+8 points). Ces runs ont été interrompus par des pannes d'infrastructure ; le résultat est à lire avec prudence, au bruit d'évaluation près ($\pm$5 points). Un run dédié (\textsc{FewShot-RL}, $G=16$, sans curriculum) est en file d'attente. En l'état, la lecture prudente est la suivante : le few-shot amorce et accélère les premières époques, sans garantie de stabilité ni de plus-value par rapport au GRPO simple, avec un risque de sur-imitation des recettes des exemples au détriment de la tâche courante. \textcolor{red}{c'est un gain rapide à étudier }

\subsubsection{Ablation inter-générations}

Le mur mesuré dépend-il du RL ou du modèle de base ? Dans notre harnais, Qwen2.5-3B nu obtient 18~\% et Qwen3.5-4B nu obtient 77~\%, au-dessus de l'objectif du papier (75~\%) sans aucun entraînement. Comme borne externe, le papier rapporte 45~\% pour Qwen3-4B dans son propre harnais~\cite{xi_agentgym-rl_2025} (tableau~\ref{tab:agentgym_textcraft}). Cette progression entre générations s'accorde avec l'exposition croissante des modèles récents aux traces agentiques dès le mid-training (section 2.3) : une partie du mur relève du modèle de base, pas de l'algorithme de RL. Ce que le RL ajouterait à chaque génération n'a en revanche pas été mesuré : aucun run RL n'a été couru sur Qwen3.5, et la question reste ouverte. \textcolor{red}{ça tombe un peu comme un cheveu dans la soupe cette partie : qu'est ce qu'on veut dire ? Oui faut prendre les modèles les plus récents parce que c ceux qui ont eu des training agentiques massifs et c de l'easy win, nous on est resté sur 2.5 pour avoir un comparratif honnete avec le paper de base et toutes les conditions de stabilisations restent valables mais les modèles récents seront certainement moins instables vu leur entrainement massif sur des env rl.}

\subsubsection{Piste SNIS}

La campagne exploratoire a testé le recollement stochastique de sous-trajectoires multi-tours repondérées par SNIS~\cite{owen2013montecarlo} (section 2.6) : les séquences reconstruites sans alignement causal violent la cohérence des observations et conduisent à un effondrement de l'apprentissage, confirmation empirique du verrou théorique identifié. La voie exécuteur-réviseur n'a pas atteint le stade d'une formulation viable des poids de transition entre tours ; nous la documentons en annexe~B à titre indicatif.

\subsubsection{Piste Transfert entre planification mono-tour et contrôle multi-tour}

\textcolor{red}{cette sousousection je l'ai décalée du 4.3 vers 4.5} 
Le curriculum ordonne le transfert de savoir entre difficultés ; une seconde forme de transfert, posée en section 2.5, sépare deux compétences. La planification mono-tour demande au modèle d'émettre en une seule complétion le plan d'actions complet, rejoué ensuite pas à pas dans l'environnement sans rétroaction. Le contrôle multi-tour choisit au contraire chaque action après observation du résultat de la précédente. \textcolor{red}{reparler du papier qui dit que en env un llm se débrouille bcp mieux qu'en full planning} \textcolor{red}{expliquer aussi clairement ce qu'on cherche à étudier : est ce que apprendre en mono tour ça aide à apprendre en multi tour et vice versa ? est ce que y a pas une optimisation via meta learning entre ces deux types d'apprentissage qui améliorerai les perfs et réduirai les couts d'entraineent ? c'est principalement motivé par le nombre de calculs beaucoup moins important en mono tour, et la possible mobilisation de toute la littérature type scpo et west of n initiale de l'offre de stage.} 

L'écart entre les deux se mesure d'abord sur un modèle non entraîné. Qwen3.5-4B atteint 54~\% en planification mono-tour contre 77~\% en interactif : la boucle de rétroaction vaut 23 points. Cette asymétrie rejoint les limites documentées de la planification en boucle ouverte des LLMs, dont les plans autonomes se dégradent dès que la tâche exige une consistance longue \cite{valmeekam2023planning}, et l'argument selon lequel la fiabilité vient d'une boucle où un vérificateur externe, ici l'environnement, corrige le modèle \cite{kambhampati2024llmmodulo}.

Le transfert entre les deux régimes d'entraînement est asymétrique. Dans le sens interactif vers mono-tour, la politique Qwen2.5-3B entraînée par RL interactif améliore sa planification mono-tour d'environ 12 à environ 19{,}5~\%, soit 7 points ; ces chiffres proviennent d'un balayage de températures et se lisent comme indicatifs, sans test statistique revendiqué. Dans le sens inverse, \textsc{Plan-Mode}, un RL mono-tour qui optimise directement la qualité du plan, atteint 37~\% en mono-tour mais retombe à 9~\% en multi-tour, contre 18~\% pour le modèle non entraîné : l'interférence est catastrophique, le modèle ne respectant plus le protocole d'une action par tour ; les profondeurs 3 et 4 restent inchangées. Améliorer le plan ne suffit donc pas, et peut coûter la compétence interactive. \textcolor{red}{mais ça c'est évidant en fait on apprend au modele à tout planifier. Non il faut regarder si alterner entre les deux permet d'accélérer le training global en terme de cout ou autre. Faut présenter ces résultats en disant on savait que ça allait pas marcher mais on a quand meme testé une première expé basée sur le prompting en espérant rien mais parce qeu c'était gratuit et rapide, on la décrit ici juste pour inspirer des expériences futures. En fait cette section très exploratoire je sais pas si ça vaudrait plutot le coup de la mettre ailleurs dans une partie plus ouverture ou considérations. Ou alors dans pistes et on présente SNIS et Transfert mono multi je vais le mettre là. Peut être il faudrait faire un tout petit paragraphe comme pour SNIS et renvoyer vers une annexe} 

Les régimes hybrides, qui alterneraient les deux entraînements, et le pré-alignement mono-tour avant RL interactif n'ont pas été courus et restent des perspectives \textcolor{parce que trop loin de ce qu'on travaille ici et trop expérimental}{aa} . La section 2.5 en donne le cadre, avec la position de Dupoux, LeCun et Malik pour qui les compétences d'agent s'acquièrent dans l'interaction avec l'environnement plutôt que par prédiction hors boucle \cite{dupoux2026why} ; nos mesures d'asymétrie vont dans ce sens sans le démontrer.

\begin{table}[htbp]
\centering
\caption{Récapitulatif des configurations de la section 4 (pass@1 sur les 100 items du test ; bruit d'évaluation $\pm$5 points ; chiffres arrêtés au 28/08). Les pass@$k$ sont mesurés à température 1{,}0.}
\label{tab:recap_section4}
\begin{tabular}{lccl}
\toprule
Configuration & pass@1 (\%) & pass@$k$ (budget) & Statut / nature \\
\midrule
Qwen2.5-3B zéro-shot & 18 & 47 ($k=20$) & baseline interne \\
Qwen2.5-3B few-shot $k=5$ (dialogue) & 31 & --- & baseline interne \\
Qwen2.5-3B few-shot (bloc système) & 5--8 & --- & baseline interne \\
\textsc{Réplication-FT} & 69 & --- & RL interne, full-FT \\
\textsc{Ancre-Mobile} & 73 & --- & RL interne, LoRA \\
\textsc{Horizon} & \textbf{82} & à produire & RL interne, record \\
\textsc{Contrôle-G16} & 54 & --- & RL interne, collapse après le pic \\
\textsc{Profondeur} & --- & --- & RL interne, en cours \\
Qwen3.5-4B nu & 77 & $\approx$80 ($k=18$) & référence inter-générations \\
AgentGym-RL-3B & 75 & --- & externe (papier) \\
Gemini 3.5 Flash (API) & 99 & --- & externe, borne haute \\
\bottomrule
\end{tabular}
\end{table}











%---------------------------------------------------------------------------------------
\newpage
\section{Délimitation du Périmètre et Ouverture}

L'auto-amélioration des agents LLM est un domaine vaste, à la croisée de l'ingénierie système, de l'apprentissage par renforcement et de l'architecture cognitive. Afin de garantir la cohérence scientifique d'un stage de six mois, le périmètre d'étude a été délimité autour des leviers algorithmiques d'optimisation paramétrique de la politique, à savoir le RLVR, le curriculum et les stratégies d'échantillonnage, dans un environnement contrôlé unique. Cette délimitation est un choix méthodologique, non un jugement de pertinence : plusieurs axes majeurs de la littérature récente, volontairement écartés, constituent autant de prolongements naturels de ce travail.

\begin{itemize}
    \item \textbf{Mémoire long-terme et amélioration non-paramétrique.} L'auto-amélioration des agents se divise en deux grands paradigmes : l'approche \emph{paramétrique}, par mise à jour des poids via RL, objet de cette étude, et l'approche \emph{non-paramétrique}, où l'agent s'améliore à poids constants en accumulant de l'expérience dans une mémoire externe consultable, de la gestion hiérarchique du contexte à la MemGPT \cite{packer_memgpt_2023} aux banques de compétences récupérées à l'inférence. Les deux paradigmes sont complémentaires plutôt que concurrents. Nos résultats few-shot (section 4.5) montrent d'ailleurs que le contenu du contexte peut débloquer ce que le gradient seul ne peut atteindre, ce qui suggère qu'un couplage mémoire--RL serait une suite naturelle.
    \item \textbf{Architectures multi-agents et optimisation du harnais.} L'orchestration de sous-agents coordonnés, déployée à l'échelle industrielle \cite{kimi_k25_2026}, et la découverte automatique d'architectures agentiques \cite{hu_adas_2025} relèvent de l'optimisation du système plutôt que de la politique. Ce travail s'est concentré sur la capacité décisionnelle d'un agent unique, condition préalable à toute composition : un orchestrateur ne peut pas compenser des sous-agents dont la politique élémentaire est instable.
    \item \textbf{Algorithmes RL alternatifs et densification de la récompense.} Les approches à réseau critique complet, le fine-tuning supervisé itératif sur trajectoires filtrées, et la densification du signal par modèles de récompense de processus \cite{lightman2023verify} n'ont pas été retenus, au profit de la simplicité et de la frugalité mémoire de GRPO sans critique. Notre expérience de reward shaping (section 4.2) éclaire ce choix en creux : densifier le signal sans le corrompre est précisément le problème que les PRM cherchent à résoudre proprement, et leur coût d'estimation en multi-tour reste un verrou ouvert. Le mélange SFT puis RL et l'extension multi-environnements systématique n'ont pas non plus été courus ; la tentative de réplication sur le modèle 7B, interrompue par un défaut logiciel, est documentée en annexe C.
    \item \textbf{Génération synthétique d'environnements et modèles de monde.} Traiter l'environnement comme un artefact à générer, qu'un agent propose ses propres tâches vérifiables \cite{zhou2025sca} ou qu'un modèle de monde synthétise intégralement les expériences d'entraînement \cite{dreamgym2025}, constitue la frontière de recherche la plus directement alignée avec le constat de ce rapport : si le passage à l'échelle des environnements est le moteur des capacités agentiques (sections 1.1 et 2.2), leur génération automatique en est le prolongement logique. L'infrastructure de données qu'exigent ces approches dépasse le cadre de ce stage ; elles en dessinent l'horizon plutôt que le contenu.
\end{itemize}

Deux travaux futurs se dégagent avec un caractère directement actionnable. Le premier est la re-stratification du jeu d'entraînement : la distribution des profondeurs (une seule recette de profondeur 4 à l'entraînement, section 3.3) est le principal facteur confondant de l'étude, et l'arbre de craft permet de générer un jeu équilibré par difficulté. Re-courir les régimes de curriculum sur un tel jeu serait le test le plus direct de l'hypothèse du mur de capacité. Le second est le couplage entre apprentissage par planification mono-tour et apprentissage interactif, dont la section 2.5 a posé le cadre et dont nos mesures d'asymétrie (section 4.3) motivent l'étude ; son orchestration apprise relèverait du méta-apprentissage.

Ces frontières dessinent en retour la portée des contributions : les régimes de curriculum et le transfert mono-tour/multi-tour étudiés ici sont des propriétés de la politique et de son signal d'apprentissage, indépendantes du harnais, de la mémoire et du mécanisme de génération des tâches, donc réutilisables dans chacun des cadres écartés. La jonction la plus prometteuse reste celle esquissée en section 2.4 : brancher un autocurriculum piloté par le progrès d'apprentissage sur un générateur de tâches ou d'environnements, de sorte que la difficulté proposée suive en permanence les capacités courantes du modèle. Elle réunirait en un seul système la structuration de l'apprentissage étudiée dans ce rapport et la génération d'environnements qui en constitue l'horizon.

\newpage
\section{Conclusion}

Ce rapport a étudié l'auto-amélioration d'agents LLM par renforcement en environnement interactif multi-tour. Son point de départ est un constat désormais partagé par les constructeurs de modèles : le facteur limitant du RL agentique n'est plus le volume de données annotées mais l'environnement d'entraînement lui-même, sa difficulté, sa diversité et la qualité de son vérificateur. Ne pouvant prétendre au passage à l'échelle des environnements accessible à l'industrie, nous en avons pris le contre-pied méthodologique : comprendre finement, dans un environnement contrôlé à récompense vérifiable, ce qui conditionne la stabilité et l'efficience de cet apprentissage, et mettre à l'épreuve l'hypothèse selon laquelle le RL ne ferait qu'aiguiser un support préexistant sans jamais l'étendre.

La première contribution est un guide de stabilisation du GRPO multi-tour, adossé à l'ensemble de nos runs. Les conditions de possibilité d'abord : le masquage des observations dans le calcul du gradient, dont le défaut est silencieux, et le rejet motivé du reward shaping syntaxique après mise en évidence d'un détournement de récompense caractérisé. Le cœur ensuite : la recette full finetuning du papier est stable mais lente et coûteuse (69\,\% après environ 247 époques) ; sa transposition LoRA selon les recettes de la littérature mono-tour collapse systématiquement ; le correctif que nous proposons, une ancre KL mobile re-basée périodiquement sur la politique courante, stabilise l'entraînement à petit rang et porte le modèle à 73\,\%. L'ablation du carré ancre-coefficient KL montre que le régime d'ancre, davantage que la valeur du coefficient, gouverne la stabilité. Enfin, la méta-analyse de nos huit collapses identifie des signaux avant-coureurs constants, l'entropie de la politique en premier lieu, et le collapse de notre contrôle à grande taille de groupe rappelle que le rappel KL seul ne suffit pas : la stabilité exige aussi une source d'entropie.

La deuxième contribution porte sur les curriculums, et son résultat principal est établi : à recette commune, le curriculum d'horizon porte Qwen2.5-3B à 82\,\% de réussite, au-dessus des 75\,\% publiés par le papier de référence sur infrastructure multi-GPU, pour deux jours de calcul sur un GPU unique. La comparaison avec son contrôle sans curriculum, qui a collapsé après une montée plus rapide, suggère que le curriculum agit comme un régularisateur d'exploration et non comme un simple accélérateur ; cette lecture reste conditionnelle à la recette commune et à la relance du contrôle. Les autres régimes, curriculum de profondeur, autocurriculum à progrès d'apprentissage appris et curriculum de budget, sont en cours ou en file au moment de la rédaction : leurs résultats seront figés au gel expérimental, et nous nous abstenons d'en préjuger. Le curriculum de profondeur a déjà fourni une observation utile, la saturation complète du signal sur un palier maintenu trop longtemps, qui illustre la limite des paliers statiques et motive les régimes adaptatifs.

La troisième contribution rassemble les mesures qui donnent leur cadre à ces résultats. L'oracle pass@$k$ établit la structure du problème : mur de fiabilité en profondeur 2, mur de capacité en profondeurs 3 et 4. Deux observations mettent l'hypothèse du plafond pass@$k$ sous tension sans la réfuter : un modèle entraîné étend sa frontière d'atteignabilité à des items que le modèle nu ne résout jamais en vingt tirages, et le pass@1 du meilleur run dépasse le pass@20 du modèle nu, sous les réserves de budget fini et d'analyse par profondeur encore à produire. Le transfert entre planification et contrôle interactif est asymétrique : la boucle de rétroaction vaut 23 points sur un modèle non entraîné, le RL interactif améliore modestement la planification, et le RL de planification détruit la compétence interactive. Le few-shot, enfin, agit comme catalyseur d'exploration, capable de créer du signal là où le gradient GRPO est identiquement nul.

Pour Criteo, ce travail constitue une double brique : une synthèse structurée d'une littérature récente et morcelée, du formalisme du Deep RL aux pratiques industrielles des constructeurs de modèles, et une infrastructure expérimentale opérationnelle, pipeline GRPO multi-tour stabilisé, environnement instrumenté et recette documentée, qui a porté un modèle de 3 milliards de paramètres au-delà du résultat publié sur la même tâche. Elle est directement réutilisable le jour où les cas d'usage internes, de la recommandation conversationnelle aux agents opérant sur les API propriétaires, appelleront la spécialisation frugale de modèles légers, dont ce rapport documente à la fois le potentiel et les conditions de réussite.


















\newpage



\newpage

\begin{thebibliography}{99}

\bibitem{vaswani2017attention}
Ashish Vaswani, Noam Shazeer, Niki Parmar, Jakob Uszkoreit, Llion Jones, Aidan~N. Gomez, Łukasz Kaiser, and Illia Polosukhin.
\newblock Attention is all you need.
\newblock {\em Advances in Neural Information Processing Systems (NeurIPS)}, 30, 2017.

\bibitem{ouyang2022training}
Long Ouyang, Jeffrey Wu, Xu~Jiang, Diogo Almeida, Carroll Wainwright, Pamela Mishkin, Chong Zhang, Sandhini Agarwal, Katarina Slama, Alex Ray, et~al.
\newblock Training language models to follow instructions with human feedback.
\newblock {\em Advances in Neural Information Processing Systems (NeurIPS)}, 35:27730--27744, 2022.

\bibitem{zelikman2022star}
Eric Zelikman, Yuhuai Wu, Jesse Mu, and Noah Goodman.
\newblock STaR: Bootstrapping reasoning with reasoning.
\newblock {\em Advances in Neural Information Processing Systems (NeurIPS)}, 35:15476--15488, 2022.

\bibitem{singh2024beyond}
Avi Singh, John~D. Co-Reyes, Rishabh Agarwal, et~al.
\newblock Beyond human data: Scaling self-training for problem-solving with language models ($\text{ReST}^{\text{EM}}$).
\newblock {\em Transactions on Machine Learning Research (TMLR)}, 2024. arXiv:2312.06585.

\bibitem{huang2024self}
Audrey Huang, Adam Block, Dylan~J. Foster, Dhruv Madeka, Jordan~T. Ash, Akshay Krishnamurthy, Max Simchowitz, and Alexander Rakhlin.
\newblock Self-improvement in language models: The sharpening mechanism.
\newblock {\em arXiv preprint arXiv:2412.01951}, 2024.

\bibitem{schulman2017proximal}
John Schulman, Filip Wolski, Prafulla Dhariwal, Alec Radford, and Oleg Klimov.
\newblock Proximal policy optimization algorithms.
\newblock {\em arXiv preprint arXiv:1707.06347}, 2017.

\bibitem{shao2024deepseekmath}
Zhihong Shao, Peiyi Wang, Qihao Zhu, Runxin Xu, Junxiao Song, Mingchuan Zhang, Y.K. Li, Y.~Wu, and Daya Guo.
\newblock DeepSeekMath: Pushing the limits of mathematical reasoning in open language models.
\newblock {\em arXiv preprint arXiv:2402.03300}, 2024.

\bibitem{rafailov2024direct}
Rafael Rafailov, Archit Sharma, Eric Mitchell, Christopher~D. Manning, Stefano Ermon, and Chelsea Finn.
\newblock Direct preference optimization: Your language model is secretly a reward model.
\newblock {\em Advances in Neural Information Processing Systems (NeurIPS)}, 36, 2024.

\bibitem{deepseekr1_2025}
DeepSeek-AI.
\newblock DeepSeek-R1: Incentivizing reasoning capability in LLMs via reinforcement learning.
\newblock {\em arXiv preprint arXiv:2501.12948}, 2025.

\bibitem{openai_o1_2024}
OpenAI.
\newblock Learning to reason with LLMs (OpenAI o1).
\newblock {\em Blog OpenAI}, \url{https://openai.com/index/learning-to-reason-with-llms/}, 2024.

\bibitem{liu2025drgrpo}
Zichen Liu, Changyu Chen, Wenjun Li, Penghui Qi, Tianyu Pang, Chao Du, Wee~Sun Lee, and Min Lin.
\newblock Understanding R1-Zero-like training: A critical perspective (Dr.~GRPO).
\newblock {\em arXiv preprint arXiv:2503.20783}, 2025.

\bibitem{yao2023react}
Shunyu Yao, Jeffrey Zhao, Dian Yu, Nan Du, Izhak Shafran, Karthik Narasimhan, and Yuan Cao.
\newblock ReAct: Synergizing reasoning and acting in language models.
\newblock {\em International Conference on Learning Representations (ICLR)}, 2023.

\bibitem{xi_agentgym_2025}
Zhiheng Xi, Yiwen Ding, Wenxiang Chen, et~al.
\newblock AgentGym: Evaluating and training large language model-based agents across diverse environments.
\newblock {\em Proceedings of ACL 2025}. arXiv:2406.04151.

\bibitem{xi_agentgym-rl_2025}
Zhiheng Xi, Jixuan Huang, Chenyang Liao, et~al.
\newblock AgentGym-RL: Training LLM agents for long-horizon decision making through multi-turn reinforcement learning.
\newblock {\em arXiv preprint arXiv:2509.08755}, 2025.

\bibitem{xi_agentgymrl_openreview}
OpenReview.
\newblock AgentGym-RL: An open-source framework to train LLM agents for long-horizon decision making via multi-turn RL.
\newblock \url{https://openreview.net/forum?id=ZgCCDwcGwn}, ICLR 2026.

\bibitem{prasad2023adapt}
Archiki Prasad, Alexander Koller, Mareike Hartmann, Peter Clark, Ashish Sabharwal, Mohit Bansal, and Tushar Khot.
\newblock ADaPT: As-needed decomposition and planning with language models.
\newblock {\em Findings of NAACL 2024}. arXiv:2311.05772.

\bibitem{carta2023glam}
Thomas Carta, Clément Romac, Thomas Wolf, Sylvain Lamprier, Olivier Sigaud, and Pierre-Yves Oudeyer.
\newblock Grounding large language models in interactive environments with online reinforcement learning (GLAM).
\newblock {\em International Conference on Machine Learning (ICML)}, 2023. arXiv:2302.02662.

\bibitem{zhang2023boss}
Jesse Zhang, Jiahui Zhang, Karl Pertsch, Ziyi Liu, Xiang Ren, Minsuk Chang, Shao-Hua Sun, and Joseph~J. Lim.
\newblock Bootstrap your own skills: Learning to solve new tasks with large language model guidance (BOSS).
\newblock {\em Conference on Robot Learning (CoRL)}, 2023. arXiv:2310.10021.

\bibitem{gaven2025magellan}
Loris Gaven, Thomas Carta, Clément Romac, Cédric Colas, Sylvain Lamprier, Olivier Sigaud, and Pierre-Yves Oudeyer.
\newblock MAGELLAN: Metacognitive predictions of learning progress guide autotelic LLM agents in large goal spaces.
\newblock {\em arXiv preprint arXiv:2502.07709}, 2025.

\bibitem{carta2025herakles}
Thomas Carta, Clément Romac, Loris Gaven, Pierre-Yves Oudeyer, Olivier Sigaud, and Sylvain Lamprier.
\newblock HERAKLES: Hierarchical skill compilation for open-ended LLM agents.
\newblock {\em arXiv preprint arXiv:2508.14751}, 2025. % Auteurs et ID à recouper sur arXiv

\bibitem{foster2026autocurriculum}
Dylan~J. Foster, Naman Rajaraman, and Alexander Rakhlin.
\newblock Theoretical guarantees for autocurriculum learning in interactive reinforcement learning.
\newblock {\em arXiv preprint arXiv:2601.00000}, 2026. % ID placeholder — À RÉSOUDRE avant rendu

\bibitem{forestier2017imgep}
Sébastien Forestier, Rémy Portelas, Yoan Mollard, and Pierre-Yves Oudeyer.
\newblock Intrinsically motivated goal exploration processes with automatic curriculum learning.
\newblock {\em arXiv preprint arXiv:1708.02190}, 2017.

\bibitem{colas2022autotelic}
Cédric Colas, Tristan Karch, Olivier Sigaud, and Pierre-Yves Oudeyer.
\newblock Autotelic agents with intrinsically motivated goal-conditioned reinforcement learning: A short survey.
\newblock {\em Journal of Artificial Intelligence Research (JAIR)}, 74:1159--1199, 2022.

\bibitem{bengio2009curriculum}
Yoshua Bengio, Jérôme Louradour, Ronan Collobert, and Jason Weston.
\newblock Curriculum learning.
\newblock {\em International Conference on Machine Learning (ICML)}, 2009.

\bibitem{zhou2025sca}
Yifei Zhou, Sergey Levine, Jason Weston, Xian Li, and Sainbayar Sukhbaatar.
\newblock Self-challenging language model agents.
\newblock {\em arXiv preprint arXiv:2506.01716}, 2025.

\bibitem{dreamgym2025}
Zhaorun Chen, et~al.
\newblock Scaling agent learning via experience synthesis (DreamGym).
\newblock {\em arXiv preprint arXiv:2511.03773}, 2025.

\bibitem{valmeekam2023planning}
Karthik Valmeekam, Alberto Olmo, Sarath Sreedharan, and Subbarao Kambhampati.
\newblock On the planning abilities of large language models: A critical investigation.
\newblock {\em Advances in Neural Information Processing Systems (NeurIPS)}, 36, 2023.

\bibitem{kambhampati2024llmmodulo}
Subbarao Kambhampati, Karthik Valmeekam, Lin Guan, et~al.
\newblock LLMs can't plan, but can help planning in LLM-modulo frameworks.
\newblock {\em arXiv preprint arXiv:2402.01817}, 2024.

\bibitem{brown2020gpt3}
Tom Brown, et~al.
\newblock Language models are few-shot learners.
\newblock {\em Advances in Neural Information Processing Systems (NeurIPS)}, 33, 2020.

\bibitem{cobbe2021verifiers}
Karl Cobbe, et~al.
\newblock Training verifiers to solve math word problems.
\newblock {\em arXiv preprint arXiv:2110.14168}, 2021.

\bibitem{wang2022selfconsistency}
Xuezhi Wang, et~al.
\newblock Self-consistency improves chain of thought reasoning in language models.
\newblock {\em arXiv preprint arXiv:2203.11171}, 2022.

\bibitem{yao2023tot}
Shunyu Yao, et~al.
\newblock Tree of Thoughts: Deliberate problem solving with large language models.
\newblock {\em arXiv preprint arXiv:2305.10601}, 2023.

\bibitem{madaan2023selfrefine}
Aman Madaan, et~al.
\newblock Self-Refine: Iterative refinement with self-feedback.
\newblock {\em arXiv preprint arXiv:2303.17651}, 2023.

\bibitem{shinn2023reflexion}
Noah Shinn, Federico Cassano, Edward Berman, Ashwin Gopinath, Karthik Narasimhan, and Shunyu Yao.
\newblock Reflexion: Language agents with verbal reinforcement learning.
\newblock {\em arXiv preprint arXiv:2303.11366}, 2023.

\bibitem{snell2024scaling}
Charlie Snell, Jaehoon Lee, Kelvin Xu, and Aviral Kumar.
\newblock Scaling LLM test-time compute optimally can be more effective than scaling model parameters.
\newblock {\em arXiv preprint arXiv:2408.03314}, 2024.

\bibitem{lightman2023verify}
Hunter Lightman, et~al.
\newblock Let's verify step by step.
\newblock {\em arXiv preprint arXiv:2305.20050}, 2023.

\bibitem{skalse2022defining}
Joar Skalse, Nikolaus Howe, Dmitrii Krasheninnikov, and David Krueger.
\newblock Defining and characterizing reward hacking.
\newblock {\em Advances in Neural Information Processing Systems (NeurIPS)}, 35, 2022.

\bibitem{owen2013montecarlo}
Art~B. Owen.
\newblock {\em Monte Carlo Theory, Methods and Examples}.
\newblock Stanford University, 2013. \url{https://artowen.su.domains/mc/}.

\bibitem{schulman2025lora}
John Schulman, et~al.
\newblock LoRA without regret.
\newblock {\em Blog Thinking Machines Lab}, \url{https://thinkingmachines.ai/blog/lora/}, 2025.

\bibitem{qwen3_2025}
Qwen Team.
\newblock Qwen3 technical report.
\newblock {\em arXiv preprint arXiv:2505.09388}, 2025.

\bibitem{qwen3coder2025}
Qwen Team.
\newblock Qwen3-Coder: Agentic coding in the world.
\newblock {\em Blog Qwen}, \url{https://qwenlm.github.io/blog/qwen3-coder/}, 2025.

\bibitem{qwen35_2026}
Qwen Team.
\newblock Qwen3.5 release notes.
\newblock 2026. % URL exacte à vérifier sur qwen.ai

\bibitem{kimik2_2025}
Kimi Team.
\newblock Kimi K2: Open agentic intelligence.
\newblock {\em arXiv preprint arXiv:2507.20534}, 2025.

\bibitem{kimi_k25_2026}
Kimi Team.
\newblock Kimi K2.5: Visual agentic intelligence.
\newblock {\em arXiv preprint arXiv:2602.02276}, 2026.

\bibitem{magistral2025}
Mistral AI.
\newblock Magistral.
\newblock {\em arXiv preprint arXiv:2506.10910}, 2025.

\bibitem{taubench2024}
Shunyu Yao, Noah Shinn, Pedram Razavi, and Karthik Narasimhan.
\newblock $\tau$-bench: A benchmark for tool-agent-user interaction in real-world domains.
\newblock {\em arXiv preprint arXiv:2406.12045}, 2024.

\bibitem{tau2bench2025}
Victor Barres, et~al.
\newblock $\tau^2$-bench: Evaluating conversational agents in a dual-control environment.
\newblock {\em arXiv preprint arXiv:2506.07982}, 2025.

\bibitem{toolathlon2025}
Junlong Li, et~al.
\newblock The Tool Decathlon: Benchmarking language agents for diverse, realistic, and long-horizon task execution.
\newblock {\em ICLR 2026}. arXiv:2510.25726.

\bibitem{mcpuniverse2025}
Salesforce AI Research.
\newblock MCP-Universe: Benchmarking large language models with real-world Model Context Protocol servers.
\newblock {\em arXiv preprint arXiv:2508.14704}, 2025.

\bibitem{mcpatlas2026}
Scale AI.
\newblock MCP-Atlas: A large-scale benchmark for tool-use competency with real MCP servers.
\newblock {\em arXiv preprint arXiv:2602.00933}, 2026.

\bibitem{zhou2023webarena}
Shuyan Zhou, et~al.
\newblock WebArena: A realistic web environment for building autonomous agents.
\newblock {\em arXiv preprint arXiv:2307.13854}, 2023.

\bibitem{shridhar2020alfworld}
Mohit Shridhar, Xingdi Yuan, Marc-Alexandre Côté, Yonatan Bisk, Adam Trischler, and Matthew Hausknecht.
\newblock ALFWorld: Aligning text and embodied environments for interactive learning.
\newblock {\em arXiv preprint arXiv:2010.03768}, 2020.

\bibitem{jimenez2023swebench}
Carlos~E. Jimenez, et~al.
\newblock SWE-bench: Can language models resolve real-world GitHub issues?
\newblock {\em arXiv preprint arXiv:2310.06770}, 2023.

\bibitem{openai2024swebenchverified}
OpenAI.
\newblock Introducing SWE-bench Verified.
\newblock {\em Blog OpenAI}, \url{https://openai.com/index/introducing-swe-bench-verified/}, 2024.

\bibitem{anthropic2024agents}
Anthropic.
\newblock Building effective agents.
\newblock \url{https://www.anthropic.com/research/building-effective-agents}, 2024.

\bibitem{anthropic2024mcp}
Anthropic.
\newblock Introducing the Model Context Protocol.
\newblock \url{https://www.anthropic.com/news/model-context-protocol}, 2024.

\bibitem{hu_adas_2025}
Shengran Hu, Cong Lu, and Jeff Clune.
\newblock Automated design of agentic systems (ADAS).
\newblock {\em ICLR 2025}. arXiv:2408.08435.

\bibitem{packer_memgpt_2023}
Charles Packer, Vivian Fang, Shishir~G. Patil, Kevin Lin, Sarah Wooders, and Joseph~E. Gonzalez.
\newblock MemGPT: Towards LLMs as operating systems.
\newblock {\em arXiv preprint arXiv:2310.08560}, 2023.

\bibitem{gao2025selfevolving}
Huan-ang Gao, Jiayi Geng, Wenyue Hua, Mengkang Hu, et~al.
\newblock A survey of self-evolving agents: What, when, how, and where to evolve on the path to artificial super intelligence.
\newblock {\em arXiv preprint arXiv:2507.21046}, 2025.

\bibitem{rajaraman2026curriculum}
Nived Rajaraman, Audrey Huang, Miroslav Dudík, Robert~E. Schapire, Dylan~J. Foster, and Akshay Krishnamurthy.
\newblock Learning to reason with curriculum~I: Provable benefits of autocurriculum.
\newblock {\em Conference on Learning Theory (COLT)}, 2026. arXiv:2603.18325.

\bibitem{graves2017automated}
Alex Graves, Marc~G. Bellemare, Jacob Menick, Rémi Munos, and Koray Kavukcuoglu.
\newblock Automated curriculum learning for neural networks.
\newblock {\em International Conference on Machine Learning (ICML)}, 2017.

\bibitem{yue2025rlvr}
Yang Yue, Zhiqi Chen, Rui Lu, Andrew Zhao, Zhaokai Wang, Shiji Song, and Gao Huang.
\newblock Does reinforcement learning really incentivize reasoning capacity in LLMs beyond the base model?
\newblock {\em Advances in Neural Information Processing Systems (NeurIPS)}, 2025. arXiv:2504.13837.

\bibitem{wang2019poet}
Rui Wang, Joel Lehman, Jeff Clune, and Kenneth~O. Stanley.
\newblock Paired open-ended trailblazer (POET): Endlessly generating increasingly complex and diverse learning environments and their solutions.
\newblock {\em arXiv preprint arXiv:1901.01753}, 2019.




% --- Entrées NON VÉRIFIÉES, citées uniquement dans la section 5 (périmètre).
% --- Vérifier leur existence sur arXiv ou retirer les citations correspondantes.
% \bibitem{hero_2025}
% Zhi Tao, Yu-Feng Li, and Zhouchen Lin.
% \newblock HERO: Hybrid exploitation and reward optimization for LLM agents.
% \newblock {\em arXiv preprint arXiv:2501.04567}, 2025.
%
% \bibitem{feng_rewardflow_2026}
% Ruishi Feng, Tianyi Zhang, and Yilun Du.
% \newblock Reward-Flow: Dense process reward estimation for multi-turn agents.
% \newblock {\em arXiv preprint arXiv:2602.01234}, 2026.
%
% \bibitem{maes_worldmodel_2026}
% Francis Maes, Jean-Baptiste Alayrac, and Timothée Lesort.
% \newblock LeWorldModel: Generative environment modeling for interactive RL.
% \newblock {\em arXiv preprint arXiv:2601.04321}, 2026.

\end{thebibliography}














\appendix



\newpage
% =================================================================
% ANNEXE 1 — Fondements du Deep RL, du cadre POMDP à GRPO
% Condensé structuré du cours CS224R (Stanford, Deep Reinforcement Learning)
% Prérequis LaTeX : amsmath, amssymb
% =================================================================
\section{Annexe 1 : Fondements du Deep Reinforcement Learning, du cadre POMDP à GRPO}
\label{annexe:maths_rl}

Cette annexe présente de manière condensée le socle théorique de l'apprentissage par renforcement profond nécessaire à la lecture de ce document. Elle s'inspire grandement et suit la progression du cours CS224R de Stanford (\emph{Deep Reinforcement Learning}, \href{https://cs224r.stanford.edu/}{cs224r.stanford.edu}) et du cours de Berkley de Sergei Levine: formalisme de décision séquentielle, apprentissage par imitation, gradients de politique, méthodes Acteur-Critique et leurs versions \emph{off-policy} menant à PPO, apprentissage par valeurs, RL hors-ligne, apprentissage de la récompense, et enfin application aux LLMs jusqu'à GRPO et au raisonnement.

% =================================================================
\subsection*{A.1 Cadre Formel : MDP et POMDP}
% =================================================================
L'apprentissage par renforcement modélise un problème de décision séquentielle : un agent observe, agit, observe le résultat de son action, et itère. À chaque pas de temps $t$, l'agent se trouve dans un état $s_t$ (l'état du « monde »), choisit une action $a_t$, reçoit une récompense instantanée $r(s_t, a_t)$ quantifiant la qualité locale du couple état-action, et l'environnement transitionne vers un état $s_{t+1}$ selon une dynamique inconnue $p(s_{t+1} \mid s_t, a_t)$. Cette dynamique satisfait la propriété de Markov : l'état suivant ne dépend que de l'état et de l'action courants, indépendamment du passé. Le tuple $(\mathcal{S}, \mathcal{A}, p, r)$, muni d'une distribution initiale $p(s_1)$, définit un processus de décision markovien (\emph{Markov Decision Process}, MDP).

Le comportement de l'agent est représenté par une politique $\pi_\theta(a \mid s)$, une distribution de probabilité sur les actions conditionnée par l'état, paramétrée par un réseau de neurones. Une séquence $\tau = (s_1, a_1, s_2, a_2, \dots, s_T, a_T)$ est appelée trajectoire (ou \emph{rollout}, ou épisode), et sa probabilité sous la politique se factorise :
\begin{equation}
p_\theta(\tau) = p(s_1) \prod_{t=1}^{T} \pi_\theta(a_t \mid s_t)\, p(s_{t+1} \mid s_t, a_t)
\end{equation}
L'objectif du RL est de maximiser l'espérance de la somme des récompenses :
\begin{equation}
\max_\theta \; J(\theta) = \mathbb{E}_{\tau \sim p_\theta(\tau)} \left[ \sum_{t=1}^{T} r(s_t, a_t) \right]
\label{eq:rl_objective}
\end{equation}
L'espérance est indispensable car deux sources de stochasticité rendent le gain aléatoire : la dynamique de l'environnement et la politique elle-même. La stochasticité de la politique n'est pas un défaut mais une nécessité, tant pour l'exploration (essayer des actions différentes est la condition de l'apprentissage par essai-erreur) que pour la modélisation de comportements variés.

Lorsque l'agent n'a pas accès à l'état complet du monde mais seulement à une observation $o_t$ (une image caméra, l'historique d'une conversation), le cadre devient un processus de décision markovien partiellement observable (\emph{Partially Observable MDP}, POMDP). La conséquence formelle est déterminante : en marginalisant les états, les observations futures dépendent de tout l'historique, $p(o_{t+1} \mid o_{1:t}, a_{1:t}) \neq p(o_{t+1} \mid o_t, a_t)$. La politique doit alors être dotée de mémoire, $\pi_\theta(a_t \mid o_{1:t})$ — ce que réalise naturellement un LLM dont le contexte accumule la trace complète des interactions. C'est ce cadre POMDP qui décrit les agents LLM étudiés dans ce document.

Ce formalisme unifie des applications très diverses : en robotique, $s$ regroupe les images et positions articulaires et $a$ une commande motrice ; pour un agent conversationnel, $o$ est le dernier message de l'utilisateur, $a$ la réponse générée, et $\tau$ la trace de la conversation. Face à l'objectif (\ref{eq:rl_objective}), les familles d'algorithmes se distinguent par leur stratégie : imiter une politique performante (imitation), différencier directement l'objectif (gradients de politique), estimer la valeur de la politique courante pour l'améliorer (Acteur-Critique), estimer la valeur de la politique optimale (méthodes par valeurs), ou apprendre la dynamique elle-même (méthodes par modèle). Chacune fait des compromis différents sur le coût de collecte des données, la stabilité et les hypothèses requises.

% =================================================================
\subsection*{A.2 Apprentissage par Imitation}
% =================================================================
L'approche la plus directe consiste à imiter des démonstrations $\mathcal{D} = \{(s_1, a_1, \dots, s_T, a_T)\}$ collectées par un expert. Le clonage comportemental (\emph{Behavior Cloning}, BC) traite ce problème par apprentissage supervisé génératif, en maximisant la vraisemblance des actions de l'expert sous la politique :
\begin{equation}
\min_\theta \; -\mathbb{E}_{(s,a) \sim \mathcal{D}} \left[ \log \pi_\theta(a \mid s) \right]
\label{eq:bc}
\end{equation}
Cette formulation probabiliste (plutôt qu'une régression $\ell_2$) est essentielle dès que les données proviennent de plusieurs démonstrateurs : une régression vers la moyenne de comportements multimodaux produit une action qui n'appartient à aucun mode, tandis qu'un modèle génératif expressif (distribution catégorielle, mélange de gaussiennes, diffusion) capture la multimodalité.

Le BC est un algorithme \emph{offline} : il n'utilise que des données existantes, sans interaction. Sa limite structurelle est le décalage de covariables (\emph{covariate shift}) : contrairement au cadre supervisé classique où les entrées sont i.i.d., les actions prédites déterminent les états futurs. Une petite erreur emmène l'agent hors de la distribution des démonstrations, où ses erreurs s'amplifient — le phénomène d'erreurs composées (\emph{compounding errors}), $p_{\text{expert}}(s) \neq p_\pi(s)$. Les remèdes reposent sur la collecte de données correctives en ligne (DAgger et ses variantes), où un expert ré-annote ou corrige les états visités par la politique apprise. Par construction, l'imitation ne peut pas dépasser le démonstrateur et n'offre aucun cadre pour s'améliorer par la pratique.

Ce cadre est le pendant exact du SFT des LLMs : la perte (\ref{eq:bc}) est celle du SFT, les démonstrations étant les paires instruction-réponse. Les erreurs composées y ont aussi leur analogue, un modèle dérivant hors distribution au fil d'une génération longue ou d'un dialogue multi-tour.

% =================================================================
\subsection*{A.3 Gradients de Politique (REINFORCE)}
% =================================================================
Les gradients de politique attaquent directement l'objectif (\ref{eq:rl_objective}), avec une difficulté : la récompense n'est en général pas différentiable (elle peut être un vérificateur symbolique, un simulateur ou un jugement humain), et l'espérance porte sur une distribution qui dépend de $\theta$. L'astuce de la dérivée du logarithme, $\nabla_\theta p_\theta(\tau) = p_\theta(\tau) \nabla_\theta \log p_\theta(\tau)$, résout les deux problèmes à la fois. En notant $R(\tau) = \sum_t r(s_t, a_t)$ et en observant que les termes de dynamique de $\log p_\theta(\tau)$ ne dépendent pas de $\theta$ :
\begin{equation}
\nabla_\theta J(\theta) = \mathbb{E}_{\tau \sim p_\theta(\tau)} \left[ \sum_{t=1}^{T} \nabla_\theta \log \pi_\theta(a_t \mid s_t) \, R(\tau) \right]
\label{eq:reinforce}
\end{equation}
Le gradient est désormais une espérance sous $p_\theta$, estimable par échantillonnage Monte-Carlo de $N$ trajectoires. L'interprétation est éclairante : (\ref{eq:reinforce}) est exactement le gradient de l'imitation (\ref{eq:bc}) appliqué aux propres générations de la politique, pondéré par la récompense. Si $R > 0$, la vraisemblance des actions de la trajectoire est augmentée ; si $R < 0$, elle est diminuée. C'est la formalisation de l'apprentissage par essai-erreur : faire davantage ce qui a bien fonctionné.

L'estimateur brut souffre d'une variance élevée, que deux raffinements classiques réduisent sans introduire de biais. D'abord la causalité : l'action au temps $t$ ne peut influencer les récompenses passées, on remplace donc $R(\tau)$ par le retour restant (\emph{reward-to-go}) $\sum_{t' \geq t} r(s_{t'}, a_{t'})$. Ensuite la \emph{baseline} : soustraire une constante $b$ (typiquement la récompense moyenne) ne change pas le gradient en espérance mais recentre le signal, produisant des gradients négatifs pour les comportements sous la moyenne :
\begin{equation}
\nabla_\theta J(\theta) \approx \frac{1}{N} \sum_{i=1}^{N} \sum_{t=1}^{T} \nabla_\theta \log \pi_\theta(a_{i,t} \mid s_{i,t}) \left( \Big( \sum_{t' = t}^{T} r(s_{i,t'}, a_{i,t'}) \Big) - b \right)
\label{eq:pg_baseline}
\end{equation}
Même raffiné, cet estimateur reste bruité — un point critique en récompense \emph{sparse}, où la plupart des trajectoires ne portent aucun signal différencié. Il est de plus strictement \emph{on-policy} : le gradient exige des échantillons de la politique courante, donc chaque pas de gradient impose de recollecter des données, ce qui est prohibitif lorsque la génération est coûteuse.

% =================================================================
\subsection*{A.4 Fonctions de Valeur et Méthodes Acteur-Critique}
% =================================================================
Les méthodes Acteur-Critique améliorent le gradient de politique en apprenant à estimer ce qui est bon plutôt qu'en se contentant de l'observer. Trois objets sont centraux. La fonction de valeur et la Q-fonction d'une politique $\pi$ sont les espérances du retour :
\begin{align}
V^\pi(s) &= \mathbb{E}_\pi \Big[ \textstyle\sum_{t' \geq t} \gamma^{t'-t}\, r(s_{t'}, a_{t'}) \,\Big|\, s_t = s \Big] \\
Q^\pi(s, a) &= \mathbb{E}_\pi \Big[ \textstyle\sum_{t' \geq t} \gamma^{t'-t}\, r(s_{t'}, a_{t'}) \,\Big|\, s_t = s, \, a_t = a \Big]
\end{align}
où $\gamma \in (0, 1]$ est un facteur d'actualisation. $V^\pi(s)$ ne dépend que de l'état : c'est le gain moyen attendu en partant de $s$ et en suivant $\pi$. $Q^\pi(s,a)$ conditionne en plus sur la première action, potentiellement différente de celle que $\pi$ aurait choisie. Elles sont reliées par $V^\pi(s) = \mathbb{E}_{a \sim \pi(\cdot \mid s)}[Q^\pi(s,a)]$, ce qui définit l'avantage :
\begin{equation}
A^\pi(s, a) = Q^\pi(s, a) - V^\pi(s)
\end{equation}
soit le surcroît de gain de l'action $a$ par rapport au comportement moyen de $\pi$ en $s$. D'espérance nulle sous $\pi$, l'avantage est la \emph{baseline} idéale, dépendante de l'état : substitué au retour brut dans (\ref{eq:pg_baseline}), il produit l'estimateur de gradient le moins bruité de cette famille :
\begin{equation}
\nabla_\theta J(\theta) \approx \frac{1}{N} \sum_{i=1}^{N} \sum_{t=1}^{T} \nabla_\theta \log \pi_\theta(a_{i,t} \mid s_{i,t}) \, \hat{A}^\pi(s_{i,t}, a_{i,t})
\label{eq:actor_critic_grad}
\end{equation}
En pratique, seul un réseau Critique $V_\phi \approx V^\pi$ est appris, par régression supervisée. Deux régimes d'estimation s'opposent pour construire ses cibles. L'estimation Monte-Carlo régresse $V_\phi(s_t)$ sur les retours observés $\sum_{t' \geq t} r_{t'}$ : elle est sans biais mais à forte variance. Le \emph{bootstrapping} (apprentissage par différence temporelle, TD) utilise l'estimation courante du réseau lui-même comme cible, $y_t = r(s_t, a_t) + \gamma V_\phi(s_{t+1})$ : la variance est faible mais un biais est introduit puisque la cible est elle-même approximative. Les cibles à $n$ pas, $y_t = \sum_{t'=t}^{t+n-1} \gamma^{t'-t} r_{t'} + \gamma^n V_\phi(s_{t+n})$, interpolent entre les deux et dominent souvent en pratique. L'avantage s'estime alors par $\hat{A}_t = r(s_t, a_t) + \gamma V_\phi(s_{t+1}) - V_\phi(s_t)$ ou sa généralisation à $n$ pas.

% =================================================================
\subsection*{A.5 Acteur-Critique Off-Policy et PPO}
% =================================================================
L'algorithme Acteur-Critique de la section précédente reste \emph{on-policy}. Pour amortir le coût de collecte, on souhaite effectuer plusieurs pas de gradient sur un même lot de trajectoires — mais dès le premier pas, la politique $\pi_\theta$ diffère de la politique $\pi_{\theta_{\text{old}}}$ ayant généré les données. L'échantillonnage préférentiel (\emph{importance sampling}) corrige ce décalage : pour toute distribution de proposition $q$, $\mathbb{E}_{x \sim p}[f(x)] = \mathbb{E}_{x \sim q}[\frac{p(x)}{q(x)} f(x)]$. Appliqué au niveau des trajectoires, le ratio $\prod_t \pi_\theta(a_t \mid s_t)/\pi_{\theta_{\text{old}}}(a_t \mid s_t)$ explose ou s'annule avec l'horizon ; l'approximation standard, en passant à une espérance par pas de temps, ne conserve que le ratio local :
\begin{equation}
r_t(\theta) = \frac{\pi_\theta(a_t \mid s_t)}{\pi_{\theta_{\text{old}}}(a_t \mid s_t)}
\end{equation}
Un danger subsiste : les avantages ayant été estimés sous $\pi_{\theta_{\text{old}}}$, la politique est incitée à s'éloigner arbitrairement d'elle pour maximiser des avantages devenus obsolètes — surapprentissage caractéristique du régime \emph{off-policy}. Deux mécanismes le contiennent. Le premier est une contrainte explicite de région de confiance, $\mathbb{E}_s[D_{\text{KL}}(\pi_\theta(\cdot \mid s) \parallel \pi_{\theta_{\text{old}}}(\cdot \mid s))] \leq \varepsilon$. Le second, retenu par \emph{Proximal Policy Optimization} (PPO) \cite{schulman2017proximal}, borne le ratio par écrêtage (\emph{clipping}) : dès que $r_t(\theta)$ sort de $[1-\epsilon, 1+\epsilon]$ dans la direction favorable, le gradient s'annule et le token ou l'action concernés cessent de contribuer à la mise à jour. Le minimum avec l'objectif non clippé garantit que l'écrêtage n'opère que lorsqu'il rend l'objectif plus conservateur :
\begin{equation}
\mathcal{L}_{\text{PPO}}(\theta) = -\mathbb{E}_{(s_t, a_t) \sim \pi_{\theta_{\text{old}}}} \left[ \min\Big( r_t(\theta)\, \hat{A}_t, \; \text{clip}\big(r_t(\theta),\, 1-\epsilon,\, 1+\epsilon\big)\, \hat{A}_t \Big) \right]
\label{eq:ppo}
\end{equation}
En pratique ($\epsilon \approx 0{,}2$), chaque lot de trajectoires sert à plusieurs époques de mises à jour avant recollecte. PPO n'est que modérément \emph{off-policy} — il réutilise le lot courant, non l'ensemble de l'expérience passée — et c'est précisément ce compromis qui fait sa stabilité, au prix d'une moindre efficacité en données que les méthodes à tampon de rejeu (\emph{replay buffer}) comme SAC. Cette stabilité explique son statut d'algorithme de référence, en robotique comme pour les LLMs.

% =================================================================
\subsection*{A.6 Aparté : Méthodes par Valeurs (Q-Learning)}
% =================================================================
Une famille alternative se passe entièrement de politique explicite. Si l'on dispose d'une estimation exacte de $Q^{\pi^*}$ pour la politique optimale, la politique gloutonne $\pi(s) = \arg\max_a Q(s, a)$ est elle-même optimale. Le Q-learning exploite l'équation d'optimalité de Bellman :
\begin{equation}
Q^{\pi^*}(s, a) = r(s, a) + \gamma \, \mathbb{E}_{s' \sim p(\cdot \mid s, a)} \Big[ \max_{a'} Q^{\pi^*}(s', a') \Big]
\end{equation}
en régressant itérativement $Q_\phi$ vers les cibles bootstrappées $y = r + \gamma \max_{a'} Q_\phi(s', a')$, sur des transitions issues d'un tampon de rejeu — l'algorithme est nativement \emph{off-policy}, donc très efficace en données. Sa stabilisation requiert un ensemble d'artifices (réseaux cibles gelés, double Q-learning contre la surestimation systématique du $\max$, cibles à $n$ pas). Le $\max_{a'}$ le rend surtout adapté aux espaces d'actions discrets de petite taille : il est inapplicable en l'état à la génération de texte, où l'« action » est une séquence de tokens dans un espace combinatoire. Nous ne le développons pas davantage, mais ses concepts (équation de Bellman, bootstrapping, surestimation des valeurs) sont les fondements de la section suivante.

% =================================================================
\subsection*{A.7 RL Hors-Ligne (Offline RL)}
% =================================================================
Le RL hors-ligne pose la question suivante : peut-on apprendre une politique performante à partir d'un jeu de données statique $\mathcal{D}$, collecté par une politique de comportement inconnue $\pi_\beta$, sans aucune interaction nouvelle ? L'enjeu pratique est considérable (réutiliser les données existantes, éviter une collecte coûteuse ou risquée), et l'avantage théorique sur l'imitation est net : en exploitant les récompenses, le RL hors-ligne peut dépasser $\pi_\beta$ et recombiner des fragments de bonnes trajectoires (\emph{trajectory stitching}) qu'aucune trajectoire du jeu de données ne réalise en entier.

Le verrou central est la surestimation des valeurs hors distribution. Si l'on applique naïvement un algorithme \emph{off-policy}, l'amélioration de politique interroge $Q_\phi$ sur des actions $a'$ absentes du jeu de données, où le réseau est arbitrairement peu fiable ; la politique recherche activement ces zones de sur-optimisme, et l'erreur s'auto-amplifie à chaque itération — sans données nouvelles pour la corriger. Toutes les méthodes de RL hors-ligne organisent leur conservatisme autour de ce problème. La ligne la plus simple, l'imitation filtrée, ne conserve que les meilleures trajectoires du jeu de données et les clone. Plus fine, la régression pondérée par l'avantage (\emph{Advantage-Weighted Regression}, AWR) pondère l'imitation de chaque transition par la qualité estimée de l'action :
\begin{equation}
\theta \leftarrow \arg\max_\theta \; \mathbb{E}_{(s,a) \sim \mathcal{D}} \left[ \log \pi_\theta(a \mid s) \, \exp\big( \tfrac{1}{\alpha} \hat{A}(s, a) \big) \right]
\label{eq:awr}
\end{equation}
Cette forme approxime la solution d'un problème d'amélioration sous contrainte KL envers $\pi_\beta$, et surtout ne rétropropage que sur des actions présentes dans les données — le sur-optimisme hors distribution est évité par construction. Des méthodes comme IQL affinent l'estimation de valeur (régression par expectiles asymétriques) pour évaluer une politique meilleure que $\pi_\beta$ sans jamais interroger d'action hors support. Ce cadre éclaire directement certains algorithmes LLM : la perte (\ref{eq:awr}) est la forme générale dont les méthodes d'auto-amélioration par filtrage (imiter ses propres trajectoires réussies, telles STaR ou ReST) sont le cas binaire, et DPO s'analysera comme un algorithme hors-ligne sur données de préférences.

% =================================================================
\subsection*{A.8 Apprentissage de la Récompense}
% =================================================================
Tout ce qui précède suppose une fonction de récompense donnée. Or, hors des jeux et des tâches vérifiables, la récompense ne va pas de soi : elle doit être apprise, et tout signal appris peut être exploité. Un classifieur de succès pré-entraîné, utilisé comme récompense, est systématiquement piraté par l'optimisation RL, qui découvre les états hors de sa distribution d'entraînement où il se trompe avec confiance (\emph{reward hacking}) ; la parade générique est adversariale, en ré-entraînant continûment le signal sur les états que la politique visite.

Pour des juges humains, l'expérience montre que les notes absolues sont bruitées et mal calibrées, tandis que les comparaisons par paires sont fiables. Le modèle de Bradley-Terry postule que la probabilité de préférer $\tau_w$ à $\tau_l$ est $\sigma(r_\theta(\tau_w) - r_\theta(\tau_l))$, où $\sigma$ est la sigmoïde, et le modèle de récompense s'apprend par maximum de vraisemblance sur les préférences annotées :
\begin{equation}
\mathcal{J}_{\text{RM}}(\phi) = -\mathbb{E}_{(x, y_w, y_l) \sim \mathcal{D}} \left[ \log \sigma\big( \text{RM}_\phi(x, y_w) - \text{RM}_\phi(x, y_l) \big) \right]
\label{eq:bradley_terry}
\end{equation}
Le juge peut aussi être un autre LLM (\emph{RL from AI Feedback}, RLAIF), l'observation clé étant que la critique est plus facile que la génération. Dans tous les cas, le modèle de récompense demeure une approximation faillible des préférences réelles — son sur-ajustement par la politique (\emph{reward over-optimization}) est documenté et motive les garde-fous de la section suivante.

% =================================================================
\subsection*{A.9 RL pour les LLMs : RLHF et DPO}
% =================================================================
Le pipeline RLHF \cite{ouyang2022training} assemble les briques précédentes. Un LLM est une politique $\pi_\theta(y \mid x)$ générant une réponse $y = (y_1, \dots, y_T)$ token par token ; le MDP sous-jacent a pour état le contexte $(x, y_{<t})$, pour action le token suivant, et pour dynamique la simple concaténation — déterministe et connue. Après pré-entraînement et SFT (l'étape d'imitation), on collecte des comparaisons par paires sur des réponses échantillonnées, on entraîne un modèle de récompense via (\ref{eq:bradley_terry}), puis on optimise la politique par PPO contre ce modèle. La récompense étant apprise donc imparfaite, l'objectif est régularisé par une pénalité d'ancrage à la politique de référence $\pi_{\text{ref}}$ (le modèle SFT) :
\begin{equation}
\max_\theta \; \mathbb{E}_{y \sim \pi_\theta(\cdot \mid x)} \left[ \text{RM}_\phi(x, y) \right] - \beta \, D_{\text{KL}}\big( \pi_\theta(\cdot \mid x) \parallel \pi_{\text{ref}}(\cdot \mid x) \big)
\label{eq:rlhf_objective}
\end{equation}
Cette KL d'ancrage est distincte de la région de confiance de PPO (section A.5) : la première est un terme de l'objectif, dirigé vers une référence figée, préservant les capacités générales et limitant le piratage de la récompense ; la seconde est un mécanisme d'optimisation, dirigé vers la politique du lot précédent.

\emph{Direct Preference Optimization} (DPO) \cite{rafailov2024direct} court-circuite ce pipeline en observant que le problème (\ref{eq:rlhf_objective}) admet une solution en forme fermée, $\pi^*(y \mid x) \propto \pi_{\text{ref}}(y \mid x) \exp(\frac{1}{\beta}\text{RM}(x,y))$, que l'on peut inverser pour exprimer la récompense en fonction de la politique : $\text{RM}(x, y) = \beta \log \frac{\pi^*(y \mid x)}{\pi_{\text{ref}}(y \mid x)} + \beta \log Z(x)$. En injectant cette reparamétrisation dans la perte de Bradley-Terry (\ref{eq:bradley_terry}) — où la constante de normalisation $Z(x)$ s'annule par différence — on obtient une perte de classification binaire portant directement sur les paramètres du LLM :
\begin{equation}
\mathcal{L}_{\text{DPO}}(\theta; \pi_{\text{ref}}) = -\mathbb{E}_{(x, y_w, y_l)} \left[ \log \sigma \left( \beta \log \frac{\pi_\theta(y_w \mid x)}{\pi_{\text{ref}}(y_w \mid x)} - \beta \log \frac{\pi_\theta(y_l \mid x)}{\pi_{\text{ref}}(y_l \mid x)} \right) \right]
\end{equation}
DPO élimine le modèle de récompense explicite, le Critique et l'échantillonnage en ligne : c'est un algorithme hors-ligne au sens de la section A.7, simple et efficace, mais qui n'exploite pas de données fraîches de la politique courante.

% =================================================================
\subsection*{A.10 RL pour le Raisonnement : RLVR et GRPO}
% =================================================================
La fragilité des récompenses apprises trouve sa résolution la plus radicale dans les domaines à vérification déterministe : en mathématiques ou en programmation, un vérificateur symbolique délivre une récompense binaire impossible à pirater. Ce RL à récompense vérifiable (RLVR) est le moteur des modèles de raisonnement : la politique génère une longue chaîne de pensée avant sa réponse finale, et seule la justesse de cette dernière est récompensée — l'intégralité de la séquence de raisonnement recevant le signal terminal. Ce paradigme a ouvert un second axe de \emph{scaling}, celui du calcul à l'inférence : l'entraînement RLVR fait croître conjointement la précision et la longueur des chaînes de pensée, le modèle apprenant de lui-même à « réfléchir plus longtemps » sur les problèmes difficiles.

L'algorithme emblématique de ce régime est \emph{Group Relative Policy Optimization} (GRPO) \cite{shao2024deepseekmath}. Dans PPO appliqué aux LLMs, le Critique $V_\phi$ est un second réseau de la taille de l'Acteur, coûteux en mémoire, et son estimation de valeur est de qualité douteuse sous récompense terminale. GRPO le supprime en estimant la \emph{baseline} statistiquement : pour chaque prompt $q$, on échantillonne un groupe de $N$ réponses $\{y^{(1)}, \dots, y^{(N)}\}$ sous $\pi_{\theta_{\text{old}}}$, un vérificateur attribue les scores $\{R_1, \dots, R_N\}$, et l'avantage de chaque réponse est son score normalisé au sein du groupe :
\begin{equation}
\hat{A}_i = \frac{R_i - \text{mean}(\mathbf{R})}{\text{std}(\mathbf{R}) + \delta}
\end{equation}
Chaque token de la réponse $i$ hérite du même avantage scalaire $\hat{A}_i$, injecté dans l'objectif clippé de PPO, auquel s'ajoute la pénalité d'ancrage :
\begin{equation}
\mathcal{L}_{\text{GRPO}}(\theta) = -\frac{1}{N} \sum_{i=1}^{N} \frac{1}{|y^{(i)}|} \sum_{t=1}^{|y^{(i)}|} \min\Big( r_{i,t}(\theta)\, \hat{A}_i, \; \text{clip}\big(r_{i,t}(\theta), 1-\epsilon, 1+\epsilon\big)\, \hat{A}_i \Big) + \beta \, D_{\text{KL}}\big( \pi_\theta \parallel \pi_{\text{ref}} \big)
\label{eq:grpo}
\end{equation}
avec $r_{i,t}(\theta) = \pi_\theta(y_t^{(i)} \mid x, y_{<t}^{(i)}) / \pi_{\theta_{\text{old}}}(y_t^{(i)} \mid x, y_{<t}^{(i)})$. Deux propriétés méritent l'attention. D'une part, si toutes les réponses du groupe obtiennent le même score — typiquement $N$ échecs sur un problème trop difficile, ou $N$ succès sur un problème trop facile —, alors $\text{std}(\mathbf{R}) = 0$ et le gradient est nul : aucun apprentissage n'a lieu, ce qui fonde l'importance du calibrage de la difficulté (et motive les approches par curriculum étudiées dans ce document). D'autre part, la diffusion uniforme de $\hat{A}_i$ sur tous les tokens renonce à toute attribution de crédit fine au sein de la trajectoire : c'est le compromis qui rend l'algorithme frugal, et la limite que les modèles de récompense de processus (PRM) cherchent à lever.

En contexte agentique multi-tour, un ajustement supplémentaire s'impose : la séquence contient des tokens émis par l'environnement (observations, résultats d'outils) qui ne sont pas des actions de la politique. Un masque binaire $M_t \in \{0,1\}$ ($1$ sur les tokens générés par l'agent, $0$ sur ceux de l'environnement) restreint la perte aux seules actions, sous peine d'entraîner le modèle à prédire les réponses déterministes de l'environnement et de dégrader la politique — l'exact analogue trajectoriel du masquage de l'instruction en SFT.

Ce socle — gradient de politique, avantage relatif de groupe, ancrage KL, masquage des observations — constitue le cadre algorithmique de l'ensemble des expériences présentées dans ce document.










\newpage


% =================================================================
% ANNEXE B — Dérivation du SNIS (suppose \appendix actif)
% =================================================================
\section{Dérivation Mathématique du SNIS dans l'Optimisation de Politique de Groupe}
\label{sec:annexe_snis}

Cette annexe détaille le développement mathématique de l'échantillonnage préférentiel auto-normalisé (\emph{Self-Normalized Importance Sampling}, SNIS) appliqué à la réestimation du gradient dans les algorithmes de type GRPO. Le formalisme s'appuie sur la théorie des méthodes Monte Carlo \cite[chap.~9]{owen2013montecarlo} et découle d'échanges avec Otmane Sakhi (Criteo AI Lab, équipe Performance Science).

\subsection*{B.1 Cadre et Formulation Initiale}
Soit $x$ l'instruction initiale, $z$ la séquence de pensée et $a$ la réponse finale. On note $\pi_\theta(z, a \mid x) = \pi_\theta(z \mid x)\, \pi_\theta(a \mid x, z)$ la distribution jointe de génération. La fonction de valeur conditionnelle à maximiser s'écrit :
\begin{equation}
V(\pi_\theta \mid x) = \mathbb{E}_{(z,a) \sim \pi_\theta(\cdot \mid x)} \left[ r(x,a) \right]
\end{equation}
où $r(x,a)$ est la récompense vérifiable de la réponse $a$. Par l'astuce de la dérivée du logarithme (\emph{score function estimator}) :
\begin{equation}
\nabla_\theta V(\pi_\theta \mid x) = \mathbb{E}_{(z,a) \sim \pi_\theta(\cdot \mid x)} \left[ \nabla_\theta \log \pi_\theta(z,a \mid x) \, r(x,a) \right]
\label{eq:snis_gradient_base}
\end{equation}

\subsection*{B.2 Changement de Mesure et Recombinaison}
L'objectif est de réévaluer la pertinence de tout couple recomposé $(z_j, a_i)$, formé à partir de la pensée d'un rollout $j$ et de la réponse d'un rollout $i$, au sein d'un même groupe de $G$ générations issues du prompt $x$. En décomposant l'espérance de (\ref{eq:snis_gradient_base}) suivant les lois marginale et conditionnelle, puis en changeant la mesure sur $a$ :
\begin{align}
\nabla_\theta V(\pi_\theta \mid x) &= \mathbb{E}_{z \sim \pi_\theta(\cdot \mid x)} \left[ \mathbb{E}_{a \sim \pi_\theta(\cdot \mid x, z)} \left[ \nabla_\theta \log \pi_\theta(z,a \mid x) \, r(x,a) \right] \right] \\
&= \mathbb{E}_{z \sim \pi_\theta(\cdot \mid x)} \left[ \mathbb{E}_{a \sim \pi_\theta(\cdot \mid x)} \left[ \frac{\pi_\theta(a \mid x, z)}{\pi_\theta(a \mid x)} \nabla_\theta \log \pi_\theta(z,a \mid x) \, r(x,a) \right] \right]
\end{align}
Le ratio de densité $\pi_\theta(a \mid x, z) / \pi_\theta(a \mid x)$ permet d'exprimer $a$ sous la marginale $\pi_\theta(a \mid x)$ plutôt que sous la conditionnelle : les réponses peuvent alors provenir de n'importe quel rollout du groupe.

\subsection*{B.3 Estimateur Monte Carlo et Interversion des Sommes}
À chaque étape GRPO, le modèle génère $G$ complétions indépendantes $(z_i, a_i)_{i \in [1, G]}$. L'approximation Monte Carlo de la double espérance s'écrit :
\begin{equation}
\widehat{\nabla_\theta V}(\pi_\theta \mid x) = \frac{1}{G} \sum_{j=1}^G \frac{1}{G} \sum_{i=1}^G \frac{\pi_\theta(a_i \mid x, z_j)}{\pi_\theta(a_i \mid x)} \nabla_\theta \log \pi_\theta(z_j, a_i \mid x) \, r(x, a_i)
\end{equation}
Les sommes étant finies, on regroupe les termes par réponse $a_i$ :
\begin{equation}
\widehat{\nabla_\theta V}(\pi_\theta \mid x) = \frac{1}{G} \sum_{i=1}^G r(x, a_i) \cdot \frac{1}{G} \sum_{j=1}^G \frac{\pi_\theta(a_i \mid x, z_j)}{\pi_\theta(a_i \mid x)} \nabla_\theta \log \pi_\theta(z_j, a_i \mid x)
\end{equation}

\subsection*{B.4 Auto-Normalisation : des Marginales Inaccessibles aux Poids SNIS}
La marginale $\pi_\theta(a_i \mid x) = \mathbb{E}_{z \sim \pi_\theta(\cdot \mid x)}[\pi_\theta(a_i \mid x, z)]$ n'est pas accessible sous forme fermée, mais les $z_l$ du groupe étant précisément tirés de $\pi_\theta(\cdot \mid x)$, elle admet l'estimateur Monte Carlo $\hat{\pi}_\theta(a_i \mid x) = \frac{1}{G} \sum_{l=1}^G \pi_\theta(a_i \mid x, z_l)$. En substituant cet estimateur au dénominateur, les facteurs $1/G$ se compensent et l'estimateur de gradient prend la forme auto-normalisée :
\begin{equation}
\widehat{\nabla_\theta V}(\pi_\theta \mid x) = \frac{1}{G} \sum_{i=1}^G r(x, a_i) \sum_{j=1}^G \tilde{w}_{ij} \, \nabla_\theta \log \pi_\theta(z_j, a_i \mid x),
\qquad
\tilde{w}_{ij} = \frac{\pi_\theta(a_i \mid x, z_j)}{\sum_{l=1}^G \pi_\theta(a_i \mid x, z_l)}
\end{equation}

\paragraph{Interprétation en termes de cible et de proposition.} Ces poids réalisent un échantillonnage préférentiel auto-normalisé au sens strict : pour chaque réponse $a_i$ fixée, la distribution \emph{cible} est le postérieur $\pi_\theta(z \mid x, a_i) \propto \pi_\theta(z \mid x)\, \pi_\theta(a_i \mid x, z)$ et la distribution de \emph{proposition} est la marginale $\pi_\theta(z \mid x)$ dont sont issus les $z_l$. Les poids d'importance non normalisés sont alors proportionnels à la vraisemblance $\pi_\theta(a_i \mid x, z_j)$, et leur normalisation sur le groupe donne exactement $\tilde{w}_{ij}$.

\paragraph{Propriétés statistiques.} Comme tout estimateur auto-normalisé — ratio de deux estimateurs Monte Carlo corrélés —, $\widehat{\nabla_\theta V}$ est \emph{biaisé à $G$ fini} mais \emph{asymptotiquement consistant} (le biais décroît en $\mathcal{O}(1/G)$), en échange d'une réduction substantielle de variance par rapport à l'échantillonnage préférentiel classique, dont le dénominateur exact est ici de toute façon inaccessible \cite[chap.~9]{owen2013montecarlo}.

\paragraph{Tempérage.} Pour lisser ou accentuer l'attribution de crédit entre pensées du groupe, une température $\alpha > 0$ peut être appliquée \emph{avant} normalisation — l'appliquer après casserait la contrainte $\sum_j \tilde{w}_{ij} = 1$ :
\begin{equation}
\tilde{w}_{ij}^{(\alpha)} = \frac{\pi_\theta(a_i \mid x, z_j)^\alpha}{\sum_{l=1}^G \pi_\theta(a_i \mid x, z_l)^\alpha}
\end{equation}
On retrouve les poids SNIS pour $\alpha = 1$, une pondération uniforme quand $\alpha \to 0$, et une affectation dure à la pensée la plus vraisemblable quand $\alpha \to \infty$. Pour $\alpha \neq 1$, l'estimateur perd l'interprétation de la section précédente et introduit un biais supplémentaire ; le tempérage doit donc être vu comme une heuristique de régularisation, non comme un raffinement statistique.

\subsection*{B.5 Implications Computationnelles et Verrou Multi-tour}
\begin{itemize}
    \item \textbf{Asymétrie computationnelle.} La génération autorégressive de $G$ nouvelles trajectoires coûte $\mathcal{O}(G \cdot L)$ passes de décodage séquentiel ($L$ longueur de séquence), tandis que l'évaluation des $G^2$ vraisemblances $\pi_\theta(a_i \mid x, z_j)$ s'effectue en passes avant parallèles par simple \emph{teacher forcing} sur les paires recomposées. Le SNIS enrichit donc le signal de gradient — $G^2$ paires pondérées au lieu de $G$ trajectoires — à coût marginal faible.
    \item \textbf{Verrou en environnement multi-tour.} En raisonnement mono-tour, la pensée $z_j$ et la réponse $a_i$ ne dépendent que de $x$. En contrôle interactif, l'action au tour $t$ dépend de l'historique strict des observations $o_{1:t-1}$ renvoyées par le simulateur : recomposer $(z_j, a_i)$ à partir de rollouts distincts juxtapose des contextes d'exécution discordants, ce qui viole la cohérence causale de la trajectoire et invalide l'hypothèse de support commun entre cible et proposition. C'est ce verrou, confirmé empiriquement par l'effondrement de l'apprentissage sous recollement naïf (section 4), qui a maintenu cette piste au stade exploratoire ; la voie exécuteur-réviseur esquissée en section 2.6 — restreindre la recomposition à des sous-séquences causalement cohérentes extraites par un modèle juge — reste à formaliser, notamment l'expression des poids de transition entre tours.
\end{itemize}




\newpage
% =================================================================
% ANNEXE C — Analyse critique
% =================================================================
\section{Analyse Critique : Parcours, Reproduction et Dette Technique}
\label{sec:annexe_critique}

L'exploration menée durant ce stage s'est inscrite dans une double complexité : l'appropriation d'une littérature récente à la croisée du Deep RL et des LLMs, et le surcoût d'ingénierie inhérent à l'expérimentation en environnement décisionnel interactif. Cette annexe en retrace les choix, les impasses et les leçons.

\subsection*{C.1 Cadrage de Recherche et Réorientation Stratégique}
La première phase du stage a été consacrée à un travail d'intégration au sein des équipes de Criteo (problématiques ACT, AIAC et Performance Science), axé sur l'état de l'art en systèmes de recommandation (RecSys), les publicités conversationnelles et les identifiants sémantiques. Après un mois et demi d'exploration, un constat s'est imposé : la littérature actuelle en RecSys repose quasi exclusivement sur des formulations mono-tours statiques, sans utilisation d'outils externes. Établir un pont direct entre recommandation et agents LLM interactifs multi-tours s'est révélé prématuré pour un stage de six mois, en l'absence de fondations établies à ce croisement académique.

Cette observation a motivé une réorientation délibérée vers l'étude fondamentale des agents autonomes et du réentraînement par renforcement à récompense vérifiable. Plutôt que de disperser les efforts sur une intersection immature, ce choix a permis de construire un cadre théorique et méthodologique rigoureux autour d'AgentGym-RL et TextCraft. Le travail fournit ainsi une synthèse structurée d'une littérature morcelée et une brique technologique et conceptuelle directement réexploitable par Criteo pour ses futurs développements agentiques.

\subsection*{C.2 Reproduction de la Baseline : Coût et Enseignements}
Une part significative du temps expérimental a été consacrée à tenter de retrouver les performances rapportées par le papier AgentGym-RL sur TextCraft, afin d'asseoir nos comparaisons sur une base équitable. Cet effort s'est révélé beaucoup plus coûteux qu'anticipé, pour des raisons qui rejoignent l'analyse de la section 2.1 : en RL agentique, le résultat dépend autant du harnais — format de prompt, parseur d'actions, budget de tours, gestion des erreurs — que des poids et de l'algorithme, et ces choix sont rarement documentés au niveau de détail qu'exigerait une reproduction exacte. La poursuite du chiffre publié, au-delà de l'établissement d'une baseline saine, a un rendement scientifique décroissant : les semaines investies ont produit peu de connaissances généralisables sur nos questions de recherche.

L'exercice n'a cependant pas été stérile. Il a imposé une exploration systématique des leviers de stabilisation — régimes de taux d'apprentissage, restauration au meilleur état, adaptation LoRA, formats de prompt — dont plusieurs sont devenus des composantes de la configuration de référence (section 4.2), et il a forgé une compréhension fine de la sensibilité du RL multi-tour à des choix d'implémentation en apparence anodins. La leçon méthodologique que nous en tirons, et que nous appliquons dans la feuille de route : traiter la baseline publiée comme un ordre de grandeur, stabiliser notre propre configuration de référence, et concentrer l'effort de rigueur sur les comparaisons \emph{internes} à protocole constant — seules porteuses de conclusions fiables.

\subsection*{C.3 Obstacles Techniques et Leçons Algorithmiques}
Contrairement à l'apprentissage supervisé sur jeux de données statiques, le réentraînement d'un agent en environnement interactif impose de surmonter une dette d'ingénierie importante. Les difficultés rencontrées — matérielles, logicielles ou algorithmiques — se résument en cinq points :

\begin{itemize}
    \item \textbf{Réutilisation d'infrastructure.} Le développement d'environnements de simulation étant un travail d'ingénierie à part entière, le choix d'AgentGym-RL a permis de s'appuyer sur des serveurs HTTP existants et de concentrer l'effort de recherche sur l'optimisation de la politique.
    \item \textbf{Masquage des observations (\texttt{env\_mask}) sous TRL.} L'absence initiale de masquage des tokens renvoyés par l'environnement entraînait le recalcul erroné des log-probabilités, des divergences KL et des ratios d'échantillonnage préférentiel sur des tokens non générés par le modèle, corrompant silencieusement le gradient. Le diagnostic de ce bug — dont le symptôme n'était qu'une convergence anormalement lente puis un effondrement — a représenté l'un des débogages les plus coûteux du stage.
    \item \textbf{Synchronisation vLLM/TRL sous contraintes système.} L'articulation entre l'inférence rapide (vLLM, génération des rollouts) et l'optimisation (TRL, perte GRPO) s'est heurtée à des incompatibilités d'environnement (version ancienne de la \texttt{glibc}), imposant un mécanisme de synchronisation personnalisé entre moteur de génération et pipeline de rétropropagation, avant une migration vers le transfert direct des poids d'adaptateurs en mémoire GPU.
    \item \textbf{Reward shaping et reward hacking.} L'ajout de récompenses intermédiaires denses (bonus de $+0{,}02$ par action syntaxiquement valide) a induit un détournement de récompense caractérisé \cite{skalse2022defining} : la politique a appris à capituler dès le premier tour pour empocher le bonus de conformité, faisant chuter le taux de réussite de 18\,\% à 8\,\%. Dans notre configuration — TextCraft, vérificateur exact, shaping syntaxique —, le signal binaire outcome-only s'est révélé strictement préférable ; la généralisation de ce constat à d'autres schémas de densification reste ouverte (voir PRM, section 5).
    \item \textbf{Stabilisation par LoRA ($r=16$).} Face à la faible densité d'information véhiculée par une récompense terminale sparse, le passage du finetuning complet à une adaptation de rang faible a dimensionné l'espace des paramètres modifiables à l'information réellement transmise par le RLVR \cite{schulman2025lora}, stoppant l'effondrement des gradients et permettant à Qwen2.5-3B d'atteindre 54\,\% de réussite (pic évalué à 58\,\%).
\end{itemize}


\end{document}